%% titre: Cours d'Optimisation
%% matiere: Optimisation
%% niveau: Licence 3
%% annee: 2026-2027
%% ordre: 2
%% description: Cours d'optimisation licence 3 2026-2027.
\documentclass[a4paper,11pt]{report}
\usepackage[utf8]{inputenc}
\usepackage[T1]{fontenc}
\usepackage[french]{babel}
\usepackage{lmodern}
\usepackage[a4paper,margin=2.5cm,headheight=14pt]{geometry}
\linespread{1.05}
\usepackage{microtype}
\usepackage[table]{xcolor}
\usepackage{amsmath,amssymb,amsthm,amsfonts,mathtools}
\usepackage{array,booktabs,multirow,multicol}
\usepackage{enumitem}
\setlist{nosep,leftmargin=*}
\usepackage{graphicx,float}
\usepackage{tikz}
\usetikzlibrary{patterns,arrows.meta,calc}
\usepackage[most]{tcolorbox}
\usepackage{fancyhdr}
\usepackage{titlesec}
% ========================================
\usepackage[table]{xcolor}
\usepackage{colortbl}

% Couleurs personnalisées
\definecolor{maincolor}{RGB}{41, 128, 185}
\definecolor{accentcolor}{RGB}{236, 240, 241}
\definecolor{darkcolor}{RGB}{44, 62, 80}
\definecolor{titlecolor}{RGB}{0, 51, 102}
\definecolor{linecolor}{RGB}{70, 130, 180}
\definecolor{titleblue}{RGB}{0, 70, 140}
\definecolor{bggray}{RGB}{245, 245, 245}
\definecolor{keywordred}{RGB}{180, 0, 0}
\definecolor{vargreen}{RGB}{0, 120, 0}
\definecolor{mathblue}{RGB}{0, 0, 120}
\definecolor{headerbg}{RGB}{0, 84, 159}
\definecolor{rowbg}{RGB}{240, 247, 255}

\usepackage[colorlinks=true,linkcolor=blue!60!black,urlcolor=blue!60!black,pdftitle={Optimisation - Synthèse de cours}]{hyperref}

\definecolor{titlecolor}{RGB}{0,51,102}
\definecolor{linecolor}{RGB}{70,130,180}

% --- titres
\titleformat{\chapter}[hang]{\Large\bfseries\color{titlecolor}}{Chapitre \thechapter.}{0.6em}{}[\vspace{2pt}\color{linecolor}\titlerule]
\titlespacing*{\chapter}{0pt}{-10pt}{14pt}
\titleformat{\section}{\large\bfseries\color{titlecolor}}{\thesection}{0.6em}{}
\titlespacing*{\section}{0pt}{10pt}{4pt}
\titleformat{\subsection}{\normalsize\bfseries\color{titlecolor}}{\thesubsection}{0.6em}{}
\titlespacing*{\subsection}{0pt}{7pt}{3pt}
\setcounter{secnumdepth}{2}
\setcounter{tocdepth}{1}

% --- en-têtes
\pagestyle{fancy}
\fancyhf{}
\renewcommand{\headrulewidth}{0.4pt}
\renewcommand{\footrulewidth}{0.4pt}
\fancyhead[L]{ULBO}\fancyhead[C]{Optimisation -- Synthèse}\fancyhead[R]{L3S5 Mathématiques}
\fancyfoot[L]{Dr A. Tougma}\fancyfoot[C]{\thepage}\fancyfoot[R]{2025--2026}
\fancypagestyle{plain}{\fancyhf{}\renewcommand{\headrulewidth}{0pt}\fancyfoot[C]{\thepage}\renewcommand{\footrulewidth}{0pt}}

% --- théorèmes
\theoremstyle{plain}
\newtheorem{theorem}{Théorème}[chapter]
\newtheorem{proposition}[theorem]{Proposition}
\newtheorem{lemma}[theorem]{Lemme}
\newtheorem{corollary}[theorem]{Corollaire}
\theoremstyle{definition}
\newtheorem{definition}[theorem]{Définition}
\newtheorem{exemple}[theorem]{Exemple}
\theoremstyle{remark}
\newtheorem{remarque}[theorem]{Remarque}
\newtheorem{hypothese}[theorem]{Hypothèse}
\renewcommand{\proofname}{Démonstration}
\newtheorem*{NB}{NB}

\tcbset{thmstyle/.style={boxrule=0.6pt,arc=2pt,before skip=7pt,after skip=7pt,breakable,left=4pt,right=4pt,top=2pt,bottom=2pt}}
\tcolorboxenvironment{theorem}{thmstyle,colback=red!5!white,colframe=red!50!black}
\tcolorboxenvironment{proposition}{thmstyle,colback=red!5!white,colframe=red!50!black}
\tcolorboxenvironment{lemma}{thmstyle,colback=red!5!white,colframe=red!50!black}
\tcolorboxenvironment{corollary}{thmstyle,colback=red!5!white,colframe=red!50!black}
\tcolorboxenvironment{definition}{thmstyle,colback=blue!5!white,colframe=blue!50!black}
\tcolorboxenvironment{exemple}{thmstyle,colback=green!5!white,colframe=green!50!black}
\tcolorboxenvironment{remarque}{thmstyle,colback=yellow!6!white,colframe=orange!60!black}
\newtcolorbox{retenir}{thmstyle,colback=gray!8,colframe=gray!60!black,title={À retenir},fonttitle=\bfseries\small,coltitle=white,colbacktitle=gray!60!black}

% --- commandes
\newcommand{\R}{\mathbb{R}}
\newcommand{\N}{\mathbb{N}}
\newcommand{\norm}[1]{\left\lVert#1\right\rVert}
\newcommand{\abs}[1]{\left\lvert#1\right\rvert}
\newcommand{\grad}{\nabla}
\newcommand{\hess}{\nabla^2}
\newcommand{\T}{^{\!\top}}
\newcommand{\sd}{\mathrm{s.c.}}
\newcommand{\scal}[2]{\langle #1,#2\rangle}
\newcommand{\M}{\mathcal{L}}
\DeclareMathOperator{\rg}{rang}
\renewcommand{\arraystretch}{1.15}
\setlength{\parindent}{0pt}
\setlength{\parskip}{3pt}
\addto\captionsfrench{\renewcommand{\chaptername}{Chapitre}\renewcommand{\contentsname}{Table des matières}}

% ========================================
% BOÎTES ET CADRES
% ========================================
\usepackage[most]{tcolorbox}
\tcbuselibrary{theorems,skins,breakable}
\usepackage{mdframed,framed}

% Style de l'encadré coloré
\newenvironment{colorframed}{%
	\def\FrameCommand{\fboxsep=10pt \fcolorbox{black}{bggray}}%
	\MakeFramed{\advance\hsize-\width \FrameRestore}%
}{\endMakeFramed}
\begin{document}
	\begin{titlepage}
		\thispagestyle{empty}
		\centering
		\vspace*{2cm}
		{\large\textsc{Université Lédéa Bernard Ouédraogo}\par}
		\vspace{1.2cm}
		{\color{linecolor}\rule{0.8\textwidth}{1.2pt}\par}
		\vspace{1cm}
		{\Huge\bfseries\color{titlecolor} Optimisation\par}
		\vspace{0.4cm}
		{\Large Synthèse de cours et exemples corrigés\par}
		\vspace{1cm}
		{\color{linecolor}\rule{0.8\textwidth}{1.2pt}\par}
		\vspace{2cm}
		{\Large\textbf{Licence 3 Mathématiques}\par}
		\vspace{0.3cm}
		{\large Semestre 5 -- 2025/2026\par}
		\vspace{0.3cm}
		{\large Cours magistral 20 h -- Travaux dirigés 6 h\par}
		\vspace{2.5cm}
		{\large\textbf{Dr A. Tougma}\par}
		\vfill
		{\small Version condensée du cours : notations, hypothèses et résultats du cours original ; erreurs signalées en annexe (Errata).\par}
	\end{titlepage}
	
	\tableofcontents
	
	\chapter*{Introduction}
	\addcontentsline{toc}{chapter}{Introduction}
	Un \textbf{problème d'optimisation} consiste à chercher, parmi un ensemble $D\subset\R^n$ de points \emph{admissibles}, un point $x^*$ où une fonction \emph{objectif} $f:\R^n\to\R$ atteint sa plus petite (ou plus grande) valeur :
	\[
	(P)\quad \min_{x\in D} f(x),\qquad \max f=-\min(-f).
	\]
	Si $D=\R^n$, le problème est \emph{sans contraintes}. Sinon $D=\{x\mid g_j(x)\le0,\ h_i(x)=0\}$ et $(P)$ est un \emph{programme mathématique}. Il est \emph{convexe} si $f$ et les $g_j$ sont convexes et les $h_i$ affines (alors $D$ est convexe) ; \emph{différentiable} si $f,g_j,h_i$ le sont ; \emph{linéaire} si tout est affine.
	
	Un point $x^*\in D$ est un \textbf{minimum local} s'il existe une boule ouverte $B$ centrée en $x^*$ telle que $f(x^*)\le f(x)$ pour tout $x\in B\cap D$ ; c'est un \textbf{minimum global} si $f(x^*)\le f(x)$ pour tout $x\in D$. Le problème est \emph{réalisable} si $D\ne\emptyset$ ; $f(x^*)$ est la \emph{valeur optimale}.
	
	\textbf{Plan.} Chapitre 1 : outils d'analyse (différentiabilité, convexité, formes quadratiques). Chapitre 2 : existence, unicité et optimisation sans contrainte. Chapitres 3--4 : contraintes d'égalité (Lagrange) puis d'égalité et d'inégalité (Karush--Kuhn--Tucker). Chapitres 5--8 : programmation linéaire (modélisation, résolution graphique, simplexe, dualité).
	
	\chapter{Rappels de calcul différentiel et d'analyse convexe}
	
%	\begin{retenir}
%		\textbf{Outils utilisés dans tout le cours} \emph{(rappels)}.
%		\begin{itemize}
%			\item Produit scalaire $\scal uv=u\T v$, norme $\norm u=\sqrt{\scal uu}$, \emph{Cauchy--Schwarz} : $\abs{\scal uv}\le\norm u\,\norm v$.
%			\item $\varphi(h)=o(\norm h^k)$ signifie $\varphi(h)/\norm h^k\to0$ quand $h\to0$ : ``négligeable devant $\norm h^k$''.
%			\item Une matrice symétrique $A$ a des valeurs propres réelles et une base orthonormée de vecteurs propres $v_1,\dots,v_n$ ; si $x=\sum\alpha_iv_i$, alors $x\T Ax=\sum\lambda_i\alpha_i^2$ et $\norm x^2=\sum\alpha_i^2$. $A$ est \emph{semi-définie positive} ($A\succeq0$) si $x\T Ax\ge0$ pour tout $x$, \emph{définie positive} ($A\succ0$) si $x\T Ax>0$ pour tout $x\ne0$ ; cela équivaut à : toutes les valeurs propres $\ge0$ (resp.\ $>0$). Alors $x\T Ax\ge\lambda_{\min}\norm x^2$.
%			\item Dérivée directionnelle : $\lim_{t\to0}\frac{f(x+td)-f(x)}t=\grad f(x)\T d$ si $f$ est différentiable en $x$.
%			\item Un ensemble fermé borné de $\R^n$ est \emph{compact} ; toute suite bornée a une sous-suite convergente (Bolzano--Weierstrass).
%		\end{itemize}
%	\end{retenir}
	
	\section{Différentielle, gradient, hessienne}
	\begin{definition}[Différentiabilité]
		Soit $U\subset\R^n$ ouvert et $f:U\to\R^m$. $f$ est \textbf{différentiable} en $x_0\in U$ s'il existe une application linéaire $L$ telle que $f(x_0+h)=f(x_0)+L(h)+o(\norm h)$ quand $h\to0$. $L=df(x_0)$ est unique ; sa matrice dans les bases canoniques est la \textbf{jacobienne} $J_f(x_0)=\left( \dfrac{\partial f_k}{\partial x_i(x_0)}\right) `_{k,i}$, et $df(x_0)(h)=J_f(x_0)h$.
	\end{definition}
	\begin{exemple}
		$f(x,y)=(x^2+y,\;xy)$ en $x_0=(1,1)$ : $J_f(1,1)=\begin{pmatrix}2&1\\1&1\end{pmatrix}$ et $f(x_0+h)-f(x_0)-J_f h=(h_1^2,\,h_1h_2)$, dont la norme vaut $\abs{h_1}\norm h=o(\norm h)$ : $f$ est différentiable en $(1,1)$.
	\end{exemple}
	\begin{definition}[Gradient]
		Si $f:U\to\R$ est différentiable en $x$, son \textbf{gradient} est l'unique vecteur $\grad f(x)$ tel que $df(x)(h)=\scal{\grad f(x)}{h}$ pour tout $h$ ; ses composantes sont les dérivées partielles.
	\end{definition}
	\textbf{Interprétation.} $\grad f(x)$ pointe dans la direction de plus forte croissance de $f$, avec pente maximale $\norm{\grad f(x)}$ ; il est orthogonal à la ligne de niveau $\{y\mid f(y)=f(x)\}$. Exemple : $f=x^2+2y^2$, $\grad f(2,1)=(4,4)$, orthogonal à l'ellipse $f=6$ en $(2,1)$.
	\begin{theorem}[Règle de la chaîne]
		Soit $g:U\to V$ différentiable en $x_0$ et $f:V\to\R^p$ différentiable en $g(x_0)$. Alors $f\circ g$ est différentiable en $x_0$ et $J_{f\circ g}(x_0)=J_f(g(x_0))\,J_g(x_0)$. Si $p=1$ : $\grad(f\circ g)(x)=J_g(x)\T\,\grad f(g(x))$.
	\end{theorem}
%	\begin{proof}
%		\textbf{Notations.} $y_0=g(x_0)$, $J_g=J_g(x_0)$, $J_f=J_f(y_0)$. \textbf{Rappel} : $g$ différentiable en $x_0$ signifie $g(x_0+k)-g(x_0)=J_gk+\varepsilon_g(k)\norm k$ avec $\varepsilon_g(k)\to0$ quand $k\to0$ ; de même $f(y_0+u)-f(y_0)=J_fu+\varepsilon_f(u)\norm u$ avec $\varepsilon_f(u)\to0$ quand $u\to0$.
%		
%		\textbf{Étape 1 : l'accroissement de $g$.} Posons $u(k)=g(x_0+k)-g(x_0)$. Alors $g(x_0+k)=y_0+u(k)$ et $u(k)=J_gk+\varepsilon_g(k)\norm k$. Comme $g$ est continue en $x_0$, $u(k)\to0$ quand $k\to0$.
%		
%		\textbf{Étape 2 : on applique la différentiabilité de $f$ à cet accroissement.}
%		\[f(g(x_0+k))-f(g(x_0))=f(y_0+u(k))-f(y_0)=J_f\,u(k)+\varepsilon_f(u(k))\norm{u(k)}.\]
%		
%		\textbf{Étape 3 : on remplace $u(k)$} dans le premier terme :
%		\[f(g(x_0+k))-f(g(x_0))=\underbrace{J_fJ_g\,k}_{\text{terme linéaire}}+\underbrace{J_f\,\varepsilon_g(k)\norm k}_{R_1}+\underbrace{\varepsilon_f(u(k))\norm{u(k)}}_{R_2}.\]
%		
%		\textbf{Étape 4 : $R_1$ et $R_2$ sont des $o(\norm k)$.}
%		\begin{itemize}
%			\item $\norm{R_1}\le\norm{J_f}\,\norm{\varepsilon_g(k)}\,\norm k$ et $\norm{\varepsilon_g(k)}\to0$ : donc $\norm{R_1}/\norm k\to0$.
%			\item Pour $k$ assez petit, $\norm{\varepsilon_g(k)}\le1$, donc $\norm{u(k)}\le\norm{J_g}\norm k+\norm k=(\norm{J_g}+1)\norm k$. Alors $\norm{R_2}\le(\norm{J_g}+1)\,\norm{\varepsilon_f(u(k))}\,\norm k$, et $\varepsilon_f(u(k))\to0$ car $u(k)\to0$ : donc $\norm{R_2}/\norm k\to0$.
%		\end{itemize}
%		\textbf{Conclusion.} $f\circ g(x_0+k)-f\circ g(x_0)=J_fJ_g\,k+o(\norm k)$ : $f\circ g$ est différentiable en $x_0$ de jacobienne $J_f(g(x_0))J_g(x_0)$. Si $p=1$, la jacobienne de $f$ est la ligne $\grad f\T$ ; on transpose $(\grad f\T J_g)\T=J_g\T\grad f$ pour obtenir le gradient.
%	\end{proof}
	\begin{exemple}
		Soit $f(y_1,y_2)=y_1^2+y_2^3$ et $g(x)=(\sin x,\cos x)\T$, $h=f\circ g=\sin^2x+\cos^3x$. On a $\grad f=(2y_1,3y_2^2)\T$ et $J_g=(\cos x,-\sin x)\T$, donc
		\[h'(x)=J_g(x)\T\grad f(g(x))=\cos x\cdot2\sin x-\sin x\cdot3\cos^2x=2\sin x\cos x-3\sin x\cos^2x,\]
		ce qui coïncide avec la dérivation directe. Cette règle est le moteur des algorithmes de descente de gradient (apprentissage de réseaux de neurones).
	\end{exemple}
	\begin{definition}[Hessienne]
		Si $f:U\to\R$ est de classe $\mathcal C^2$, sa \textbf{hessienne} $$\hess f(x)=\left( \dfrac{\partial^2f}{\partial x_i\partial x_j(x)}\right) _{ij}$$ est symétrique (\emph{théorème de Schwarz} : les dérivées secondes croisées, supposées continues au voisinage de $x_0$, coïncident en $x_0$).
	\end{definition}
	\begin{theorem}[Taylor à l'ordre 2]
		Soit $f\in\mathcal C^2(U)$. 
		\begin{enumerate}
			\item \textbf{Taylor--Young} : $f(x+h)=f(x)+\grad f(x)\T h+\tfrac12h\T\hess f(x)h+o(\norm h^2)$. 
			\item \textbf{Taylor--Lagrange} : si le segment $[x,x+h]\subset U$, il existe $c\in]0,1[$ tel que $f(x+h)=f(x)+\grad f(x)\T h+\tfrac12h\T\hess f(x+ch)h$.
		\end{enumerate}
	\end{theorem}
	Au voisinage d'un point critique ($\grad f(x^*)=0$), le comportement de $f$ est donc gouverné par le signe de la forme quadratique $h\mapsto h\T\hess f(x^*)h$.
	
%	\begin{exemple}[Un développement de Taylor concret]
%		$f(x,y)=x^2y$ en $(1,1)$ : $\grad f=(2xy,\,x^2)=(2,1)$, $\hess f=\begin{pmatrix}2y&2x\\2x&0\end{pmatrix}=\begin{pmatrix}2&2\\2&0\end{pmatrix}$. Donc
%		\[f(1+h_1,1+h_2)=1+2h_1+h_2+\tfrac12(2h_1^2+4h_1h_2)+o(\norm h^2)=1+2h_1+h_2+h_1^2+2h_1h_2+o(\norm h^2).\]
%		Vérification par calcul exact : $(1+h_1)^2(1+h_2)=1+2h_1+h_2+h_1^2+2h_1h_2+h_1^2h_2$, et le reste $h_1^2h_2$ est bien un $o(\norm h^2)$ (il est d'ordre 3).
%	\end{exemple}
	
	\begin{remarque}[Lecture de la formule de Taylor]
		Dans $f(x+h)=f(x)+\grad f(x)\T h+\frac12h\T\hess f(x)h+o(\norm h^2)$, le premier terme est la valeur de base, le deuxième l'approximation linéaire (plan tangent), le troisième l'approximation quadratique qui décrit la \emph{courbure}. Le reste peut aussi s'écrire sous forme intégrale : $f(x+h)=f(x)+\grad f(x)\T h+\int_0^1(1-t)\,h\T\hess f(x+th)h\,dt$. La formule de Taylor--Lagrange généralise le théorème des accroissements finis et sert surtout aux preuves théoriques.
	\end{remarque}
	
	\section{Ensembles et fonctions convexes}
	\begin{definition}[]\quad
		\begin{itemize}
			\item $C\subset\R^n$ est \textbf{convexe} si $\lambda x+(1-\lambda)y\in C$ pour tous $x,y\in C$, $\lambda\in[0,1]$. 
			\item Une fonction $f:C\to\R$ est 
			\begin{itemize}
				\item \textbf{convexe} si
				$f(\lambda x+(1-\lambda)y)\le\lambda f(x)+(1-\lambda)f(y)$ (inégalité de Jensen), 
				\item \textbf{strictement convexe} si l'inégalité est stricte pour $x\ne y$, $\lambda\in]0,1[$.
			\end{itemize}
		\end{itemize}
	\end{definition}
	Les ensembles $\{x\mid g(x)\le0\}$ avec $g$ convexe, les demi-espaces et les hyperplans sont convexes. Une caractérisation équivalente de la convexité de $f$ est la convexité de son \textbf{épigraphe} $\{(x,t)\in C\times\R\mid t\ge f(x)\}$ : géométriquement, le graphe de $f$ est sous chacune de ses cordes.
	\begin{theorem}[Caractérisations différentielles]\label{thm:conv}
		Soit $f$ différentiable sur un ouvert contenant le convexe $C$.
		\begin{enumerate}
			\item $f$ est convexe sur $C$ $\iff$ $f(y)\ge f(x)+\grad f(x)\T(y-x)$ pour tous $x,y\in C$ ; strictement convexe $\iff$ inégalité stricte pour $x\neq y$.
			\item Si $f$ est $\mathcal C^2$ et $C$ ouvert : $f$ convexe $\iff$ $\hess f(x)$ semi-définie positive pour tout $x\in C$ ; si $\hess f(x)$ est définie positive pour tout $x$, alors $f$ est strictement convexe (\textbf{la réciproque est fausse} : $x\mapsto x^4$ est strictement convexe mais $f''(0)=0$).
		\end{enumerate}
	\end{theorem}
%	\begin{proof}
%		\textbf{(1), sens $\Rightarrow$.} On suppose $f$ convexe et on fixe $x,y\in C$.
%		\begin{enumerate}[label=\emph{Étape \arabic*.},leftmargin=*,itemsep=2pt]
%			\item Pour $\theta\in]0,1]$, le point $x+\theta(y-x)=(1-\theta)x+\theta y$ est dans $C$ ; l'inégalité de Jensen donne $f(x+\theta(y-x))\le(1-\theta)f(x)+\theta f(y)$.
%			\item On retranche $f(x)$ et on divise par $\theta>0$ : $\dfrac{f(x+\theta(y-x))-f(x)}{\theta}\le f(y)-f(x)$.
%			\item On fait $\theta\to0^+$. Le membre de gauche est un taux d'accroissement de $f$ en $x$ dans la direction $y-x$ : il tend vers $\grad f(x)\T(y-x)$ car $f$ est différentiable. Une inégalité large passe à la limite, d'où $\grad f(x)\T(y-x)\le f(y)-f(x)$.
%		\end{enumerate}
%		\emph{Version stricte} ($f$ strictement convexe, $x\ne y$) : avec $z=\frac{x+y}2$, l'étape précédente appliquée à $(x,z)$ donne $f(z)\ge f(x)+\frac12\grad f(x)\T(y-x)$ ; la stricte convexité donne $f(z)<\frac12f(x)+\frac12f(y)$. En combinant, $\frac12\grad f(x)\T(y-x)<\frac12\bigl(f(y)-f(x)\bigr)$, soit l'inégalité stricte.
%		
%		\textbf{(1), sens $\Leftarrow$.} On suppose $f(y)\ge f(x)+\grad f(x)\T(y-x)$ pour tous $x,y$. Soit $x,y\in C$, $\theta\in[0,1]$ et $z=(1-\theta)x+\theta y$.
%		\begin{enumerate}[label=\emph{Étape \arabic*.},leftmargin=*,itemsep=2pt]
%			\item On applique l'hypothèse avec le point de base $z$, successivement vers $x$ puis vers $y$ : $f(x)\ge f(z)+\grad f(z)\T(x-z)$ et $f(y)\ge f(z)+\grad f(z)\T(y-z)$.
%			\item On multiplie la première par $(1-\theta)\ge0$, la seconde par $\theta\ge0$ (les inégalités ne changent pas de sens) et on additionne : $(1-\theta)f(x)+\theta f(y)\ge f(z)+\grad f(z)\T\bigl[(1-\theta)(x-z)+\theta(y-z)\bigr]$.
%			\item Le crochet vaut $(1-\theta)x+\theta y-z=0$. Il reste $(1-\theta)f(x)+\theta f(y)\ge f(z)$ : c'est l'inégalité de Jensen. Si l'hypothèse est stricte pour $x\ne y$, l'inégalité obtenue l'est aussi pour $\theta\in]0,1[$.
%		\end{enumerate}
%		
%		\textbf{(2), sens $\Leftarrow$.} Soit $x\ne y$ dans $C$ ; posons $h=y-x$. Taylor--Lagrange donne un $z\in]x,y[$ tel que $f(y)=f(x)+\grad f(x)\T h+\frac12h\T\hess f(z)h$. Si $\hess f(z)\succeq0$, le dernier terme est $\ge0$ ; si $\hess f(z)\succ0$ et $h\ne0$ il est $>0$. On obtient l'inégalité (large, resp. stricte) de (1), donc la convexité (resp. stricte).
%		
%		\textbf{(2), sens $\Rightarrow$ (convexe $\Rightarrow$ hessienne $\succeq0$).} Fixons $x\in C$, $h\in\R^n$ et $t>0$ petit. Par (1), $f(x+th)-f(x)-t\grad f(x)\T h\ge0$. Par Taylor--Young, le membre de gauche vaut $\frac{t^2}2h\T\hess f(x)h+o(t^2)$. On divise par $t^2$ : $\frac12h\T\hess f(x)h+\frac{o(t^2)}{t^2}\ge0$ ; avec $t\to0$, $h\T\hess f(x)h\ge0$.
%	\end{proof}

	
	\section{Fonctionnelles quadratiques et critère de Sylvester}
	Soit $f(x)=\frac12x\T Ax-b\T x+c$, avec $A$ symétrique (quitte à remplacer $A$ par $\frac12(A+A\T)$). Alors
	\[\grad f(x)=Ax-b,\qquad \hess f(x)=A.\]
	Par le théorème~\ref{thm:conv}, $f$ est convexe $\iff A$ semi-définie positive, strictement convexe $\iff A$ définie positive (ici l'équivalence est vraie car $\hess f$ est constante). Si $A$ est définie positive, l'unique point critique $x^*=A^{-1}b$ est l'unique minimum global.
	
	\emph{Calcul.} En développant $f(x)=\frac12\sum_{i,j}A_{ij}x_ix_j-\sum_kb_kx_k+c$ :
	$\frac{\partial f}{\partial x_k}=\frac12\sum_jA_{kj}x_j+\frac12\sum_iA_{ik}x_i-b_k=(Ax)_k-b_k$ par symétrie de $A$, puis $\frac{\partial^2f}{\partial x_l\partial x_k}=A_{kl}$.
	\textbf{Points critiques.} $\grad f(x^*)=Ax^*-b=0$. Si $A$ est inversible, il existe un unique point critique $x^*=A^{-1}b$ ; si de plus $A\succ0$, c'est l'unique minimum global ; si $A$ n'est pas semi-définie positive, le point critique (s'il existe) n'est \emph{pas} un minimum.
	
	\begin{definition}
		Le $k$-ième \textbf{mineur principal dominant} $\Delta_k$ de $A$ est le déterminant de la sous-matrice $k\times k$ extraite du coin supérieur gauche.
	\end{definition}
	\begin{theorem}[Critère de Sylvester]
		Soit $A$ symétrique. $A$ est définie positive $\iff\Delta_k>0$ pour tout $k$. $A$ est définie négative $\iff(-1)^k\Delta_k>0$ pour tout $k$.
	\end{theorem}
	\begin{remarque}
		Si certains $\Delta_k$ sont nuls, on ne peut pas conclure pour le caractère semi-défini : il faut tous les mineurs principaux, ou les valeurs propres ($A\succeq0\iff$ toutes les valeurs propres $\ge0$).
	\end{remarque}
	\begin{exemple}
		$A=\begin{pmatrix}-2&1\\1&-2\end{pmatrix}$ : $\Delta_1=-2<0$ et $\Delta_2=3>0$, donc $(-1)^k\Delta_k>0$ : $A$ est définie négative.
	\end{exemple}
	\begin{exemple}
		$f(x,y,z)=x^2+2y^2+3z^2+2xy+2yz=\frac12v\T Av$ avec $A=\begin{pmatrix}2&2&0\\2&4&2\\0&2&6\end{pmatrix}$ : $\Delta_1=2$, $\Delta_2=4$, $\Delta_3=16$, tous $>0$. $A$ est définie positive, $f$ est strictement convexe et son unique minimum global est $(0,0,0)$.
	\end{exemple}
	\chapter{Existence, unicité et optimisation sans contrainte}
	On considère $(P)\ \min_{x\in\R^n}f(x)$. Ce chapitre répond à trois questions : le minimum \emph{existe-t-il} (\S2.1) ? est-il \emph{unique} (\S2.1) ? comment le \emph{caractériser et le calculer} (\S2.2 et 2.3) ?
	
	\section{Conditions d'existence et d'unicité d'un minimum}
	
	\begin{definition}[Ensemble compact]
		Un sous-ensemble $K \subset \R^n$ est dit \textbf{compact} s'il est fermé et borné.
		\begin{itemize}
			\item \textbf{Fermé} signifie que $K$ contient la limite de toute suite convergente d'éléments de $K$.  
			\item \textbf{Borné} signifie qu'il existe un réel $M > 0$ tel que pour tout $x \in K$, $\norm{x} \le M$.
		\end{itemize}
	\end{definition}
	
	\begin{theorem}[Weierstrass]
		Si $K\subset\R^n$ est compact (fermé borné) non vide et $f:K\to\R$ continue, $f$ est bornée et atteint son minimum et son maximum sur $K$. C'est-à-dire qu'il existe au moins un point $x_{min} \in K$ (minimum global) et au moins un point $x_{max} \in K$ (maximum global) tels que :
		$$ \forall x \in K, \quad f(x_{min}) \le f(x) \le f(x_{max}) $$
	\end{theorem}
%	\begin{definition}
%		$f:\R^n\to\R$ est \textbf{coercive} si $f(x)\to+\infty$ quand $\norm x\to+\infty$.\\ \textbf{(Carractérisation de la Coercivité)} S'il existe $g:\R\to\R$ avec $f(x)\ge g(\norm x)$ et $g(t)\to+\infty$ ($t\to+\infty$), alors $f$ est coercive.
%	\end{definition}
%	\begin{theorem}[Existence]
%		Une fonction $f:\R^n\to\R$ continue et coercive admet au moins un minimum global.
%	\end{theorem}
%	\begin{proof}
%	\textcolor{red}{	\emph{Idée} : la coercivité permet de se ramener à une boule fermée bornée, où Weierstrass s'applique.
%		\emph{Étape 1 (choix d'un niveau).} Prenons un point quelconque $x_0\in\R^n$ et posons $A=f(x_0)$. Par définition de la coercivité (avec cette valeur $A$), il existe $R>0$, que l'on peut prendre $\ge\norm{x_0}$, tel que $\norm x>R\Rightarrow f(x)>A$.
%		\emph{Étape 2 (minimum sur la boule).} La boule fermée $\bar B(0,R)$ est fermée et bornée, donc compacte, et $f$ est continue : d'après Weierstrass, il existe $x^*\in\bar B(0,R)$ tel que $f(x^*)\le f(x)$ pour tout $x\in\bar B(0,R)$.
%		\emph{Étape 3 (comparaison).} Comme $x_0\in\bar B(0,R)$, $f(x^*)\le f(x_0)=A$.
%		\emph{Étape 4 (hors de la boule).} Si $\norm x>R$, alors $f(x)>A\ge f(x^*)$. Si $\norm x\le R$, $f(x)\ge f(x^*)$ par l'étape 2. Dans tous les cas $f(x)\ge f(x^*)$ : $x^*$ est un minimum global.}
%	\end{proof}
%	\begin{theorem}[Unicité, et existence-unicité]
%		Si $f$ est strictement convexe, elle a au plus un minimiseur : si $x_1^*\ne x_2^*$ étaient deux minimiseurs, $f(\frac{x_1^*+x_2^*}2)<\min f$, absurde. Si de plus $f$ est continue et coercive, le minimiseur existe : il est \textbf{unique}.
%	\end{theorem}
%	\begin{proof}[Preuve de l'unicité]
%		Supposons deux minimiseurs distincts $x_1^*\ne x_2^*$, de valeur commune $m=\min f$. La stricte convexité avec $\lambda=\frac12$ donne $f\bigl(\frac{x_1^*+x_2^*}2\bigr)<\frac12f(x_1^*)+\frac12f(x_2^*)=m$. Mais le milieu est un point de $\R^n$ où $f$ ne peut pas prendre une valeur inférieure à $m=\min f$ : contradiction.
%	\end{proof}
%	\begin{theorem}
%		Si $f\in\mathcal C^1$ et s'il existe $\alpha>0$ tel que $\scal{\grad f(x)-\grad f(y)}{x-y}\ge\alpha\norm{x-y}^2$ pour tous $x,y$, alors $f$ est strictement convexe et coercive ; $(P)$ a donc une unique solution.
%	\end{theorem}
%	\begin{remarque}
%		Les hypothèses sont nécessaires : $f(x)=x$ est continue mais n'a pas de minimum sur $\R$ (non coercive). Autre démonstration de l'existence, par \textbf{suite minimisante} : soit $m=\inf f$ et $(x_k)$ avec $f(x_k)\to m$. Si $(x_k)$ n'était pas bornée, une sous-suite vérifierait $\norm{x_{\psi(k)}}\to\infty$, donc $f(x_{\psi(k)})\to+\infty$ par coercivité, ce qui contredit $f(x_k)\to m<\infty$. La suite est donc bornée ; (Bolzano--Weierstrass) une sous-suite converge vers $x^*$ et, $f$ étant continue, $f(x^*)=m$. En particulier $m>-\infty$.
%	\end{remarque}

%=============================================================================================================
	\begin{definition} Une fonction $f: \Re^{n} \longrightarrow \Re$ est dite coercive si et seulement si :
	
	$$
	\begin{gathered}
		\lim _{\|x\| \longrightarrow+\infty} f(x)=+\infty . \\
		\left(c-\grave{a}-d: \forall A \in \Re_{+}^{*} \exists N>0 / \forall x \in \Re^{n},\|x\| \geq N \Longrightarrow f(x) \geq A\right) .
	\end{gathered}
	$$
\end{definition}
Pour montrer que $f$ est coercive on utilise souvent la proposition suivante :\\
\begin{proposition}[Caractérisations de la coercivité] Soit $f: \Re^{n} \longrightarrow \Re$ une fonction et $g: \Re \longrightarrow \Re$ vérifiant :
	
	$$
	f(x) \geq g(\|x\|) \text { avec } \lim _{t \longrightarrow+\infty} g(t)=+\infty
	$$
	
	Alors f est coercive.\\
\end{proposition}

\begin{theorem}[Existence] Soit $f: \Re^{n} \longrightarrow \Re$ une fonction propre, continue, et, coercive, alors il existe un point $x^{*} \in \Re^{n}$ qui réalise le minimum de $f$ sur $\Re^{n}$.
	
	Autrement dit $\exists x^{*} \in \Re^{n}$ tel que :
	
	$$
	f\left(x^{*}\right) \leq f(x) \quad, \forall x \in \Re^{n}
	$$
\end{theorem}
\begin{proof}
	Soit \(f:\mathbb{R}^{n}\to\mathbb{R}\) continue et coercive, c'est-à-dire
	\[
	\lim_{\|x\|\to+\infty} f(x)=+\infty.
	\]
	On veut montrer qu'il existe \(x^{*}\in\mathbb{R}^{n}\) tel que
	\[
	f(x^{*})\leq f(x)\qquad \forall x\in\mathbb{R}^{n}.
	\]
	
	Comme \(f\) est coercive, il existe \(R>0\) tel que
	\[
	\|x\|>R \quad\Longrightarrow\quad f(x)>f(0).
	\]
	En effet, il suffit d'appliquer la définition de la coercivité avec la valeur \(f(0)\in\mathbb{R}\).
	
	Considérons la boule fermée
	\[
	B_R=\{x\in\mathbb{R}^{n}:\|x\|\leq R\}.
	\]
	Dans \(\mathbb{R}^{n}\), \(B_R\) est compacte. Comme \(f\) est continue, par le théorème des bornes atteintes de Weierstrass, \(f\) atteint son minimum sur \(B_R\). Il existe donc \(x^{*}\in B_R\) tel que
	\[
	f(x^{*})\leq f(x)\qquad \forall x\in B_R.
	\]
	
	Montrons que \(x^{*}\) est en fait un minimiseur global. Soit \(x\in\mathbb{R}^{n}\).
	
	\begin{itemize}
		\item Si \(x\in B_R\), alors par définition de \(x^{*}\),
		\[
		f(x^{*})\leq f(x).
		\]
		
		\item Si \(x\notin B_R\), alors \(\|x\|>R\), donc
		\[
		f(x)>f(0).
		\]
		Or \(0\in B_R\), donc
		\[
		f(x^{*})\leq f(0).
		\]
		Ainsi
		\[
		f(x)>f(0)\geq f(x^{*}).
		\]
	\end{itemize}
	
	Dans tous les cas,
	\[
	f(x^{*})\leq f(x)\qquad \forall x\in\mathbb{R}^{n}.
	\]
	
	Donc \(x^{*}\) réalise le minimum global de \(f\) sur \(\mathbb{R}^{n}\). Le théorème est prouvé.
\end{proof}
\begin{theorem}[Unicité] Soit $f: \Re^{n} \longrightarrow \Re$ une fonction strictement convexe, alors il existe au plus $x^{*} \in \Re^{n}$ tel que : $f\left(x^{*}\right)=\displaystyle\min _{x \in \Re^{n}} f(x)$.
\end{theorem}
\begin{proof}
	Supposons, par l'absurde, qu'il existe deux points distincts \(x,y\in\mathbb{R}^{n}\) tels que
	\[
	f(x)=f(y)=\min_{z\in\mathbb{R}^{n}} f(z)=:m.
	\]
	Comme \(f\) est strictement convexe, pour tout \(t\in(0,1)\),
	\[
	f(tx+(1-t)y)<t f(x)+(1-t)f(y).
	\]
	En particulier, pour \(t=\frac12\), en posant \(z=\frac{x+y}{2}\), on obtient
	\[
	f(z)<\frac12 f(x)+\frac12 f(y)=\frac12 m+\frac12 m=m.
	\]
	Or \(m\) est le minimum global de \(f\), donc \(f(z)\geq m\), ce qui est contradictoire. Ainsi, il ne peut pas exister deux minimiseurs distincts. Le minimum, s'il existe, est donc unique.
\end{proof}

\begin{theorem}[Existence et unicité] 
	Soit $f: \Re^{n} \rightarrow \Re$, vérifiant :
	
	\begin{enumerate}
		\item $f$ continue;
		\item $f$ coercive;
		\item $f$ est strictement convexe.
	\end{enumerate}
	
	Alors il existe un unique $x^{*} \in \Re^{n}$, tel que $f\left(x^{*}\right)=\displaystyle\min _{x \in \Re^{n}} f(x)$.\\
\end{theorem}

\begin{theorem}
	Soit $f: \Re^{n} \rightarrow \Re$, une fonction de classe $C^{1}$. Supposons qu'il existe $\alpha>0$ tel que :
	
	$$
	\forall x, y \in \Re^{n} \quad, \quad\langle\nabla f(x)-\nabla f(y), x-y\rangle \geq \alpha\|x-y\|^{2}
	$$
	
	Alors $f$ est strictement convexe et coercive, en particulier le problème ( $P$ ) admet une solution unique .
\end{theorem}
%==============================================================================================================

%	\begin{retenir}
%		Existence : continuité + coercivité (ou compacité). Unicité : stricte convexité. Convexité : $\hess f\succeq0$ (la définie positive est \emph{suffisante} pour la stricte convexité, non nécessaire).
%	\end{retenir}
	
	\section{Conditions d'optimalité}
	Un point $x$ tel que $\grad f(x)=0$ est dit \textbf{critique} (ou stationnaire).
	\begin{theorem}[Conditions nécessaires]\label{thm:cn}
		Soit $f:\R^n\to\R$ et $x^*$ un minimum local. Si $f$ est différentiable en $x^*$ : (1) $\grad f(x^*)=0$ (ordre 1). Si $f$ est de classe $\mathcal C^2$ : (2) $\hess f(x^*)$ est semi-définie positive (ordre 2).
	\end{theorem}
%	\begin{proof}
%		\textbf{(1) Ordre 1.} Fixons une direction $d\in\R^n$. Comme $x^*$ est un minimum local, $f(x^*+td)\ge f(x^*)$ pour $\abs t$ assez petit. Considérons le quotient $q(t)=\frac{f(x^*+td)-f(x^*)}{t}$, dont la limite quand $t\to0$ est la dérivée directionnelle $d\T\grad f(x^*)$ (différentiabilité).
%		\begin{itemize}
%			\item Pour $t>0$ : numérateur $\ge0$ et dénominateur $>0$, donc $q(t)\ge0$ ; avec $t\to0^+$ : $d\T\grad f(x^*)\ge0$.
%			\item Pour $t<0$ : numérateur $\ge0$ mais dénominateur $<0$, donc $q(t)\le0$ ; avec $t\to0^-$ : $d\T\grad f(x^*)\le0$.
%		\end{itemize}
%		Donc $d\T\grad f(x^*)=0$ pour tout $d$. En choisissant $d=\grad f(x^*)$, on obtient $\norm{\grad f(x^*)}^2=0$, soit $\grad f(x^*)=0$.
%		
%		\textbf{(2) Ordre 2.} Comme $\grad f(x^*)=0$, Taylor--Young s'écrit $f(x^*+td)-f(x^*)=\frac{t^2}2d\T\hess f(x^*)d+o(t^2)$. Le membre de gauche est $\ge0$ pour $t$ petit. On divise par $t^2>0$ : $\frac12d\T\hess f(x^*)d+\frac{o(t^2)}{t^2}\ge0$. Quand $t\to0$, le dernier terme tend vers $0$ : $d\T\hess f(x^*)d\ge0$. Comme $d$ est arbitraire, $\hess f(x^*)$ est semi-définie positive.
%	\end{proof}
\begin{proof}
	\textbf{(1) Ordre 1.}
	Fixons une direction \(d\in\R^n\). Comme \(x^*\) est un minimum local,
	\[
	f(x^*+td)\ge f(x^*)
	\]
	pour \(\abs t\) assez petit. Considérons le quotient
	\[
	q(t)=\frac{f(x^*+td)-f(x^*)}{t},
	\]
	dont la limite quand \(t\to0\) est la dérivée directionnelle \(d\T\grad f(x^*)\) (différentiabilité).
	\begin{itemize}
		\item Pour \(t>0\) : le numérateur est \(\ge0\) et le dénominateur \(>0\), donc \(q(t)\ge0\) ; avec \(t\to0^+\) :
		\[
		d\T\grad f(x^*)\ge0.
		\]
		\item Pour \(t<0\) : le numérateur est \(\ge0\) mais le dénominateur \(<0\), donc \(q(t)\le0\) ; avec \(t\to0^-\) :
		\[
		d\T\grad f(x^*)\le0.
		\]
	\end{itemize}
	Donc \(d\T\grad f(x^*)=0\) pour tout \(d\). En choisissant \(d=\grad f(x^*)\), on obtient
	\[
	\norm{\grad f(x^*)}^2=0,
	\]
	soit \(\grad f(x^*)=0\).
	
	\textbf{(2) Ordre 2.}
	Comme \(\grad f(x^*)=0\), Taylor--Young s'écrit
	\[
	f(x^*+td)-f(x^*)=\frac{t^2}{2}d\T\hess f(x^*)d+o(t^2).
	\]
	Le membre de gauche est \(\ge0\) pour \(t\) petit. On divise par \(t^2>0\) :
	\[
	\frac12 d\T\hess f(x^*)d+\frac{o(t^2)}{t^2}\ge0.
	\]
	Quand \(t\to0\), le dernier terme tend vers \(0\) :
	\[
	d\T\hess f(x^*)d\ge0.
	\]
	Comme \(d\) est arbitraire, \(\hess f(x^*)\) est semi-définie positive.
\end{proof}
	\begin{theorem}[Condition suffisante]\label{thm:cs}
		Soit $f\in\mathcal C^2(\R^n)$ et $x^*$ tel que $\grad f(x^*)=0$ et $\hess f(x^*)$ définie positive. Alors $x^*$ est un minimum local strict.
	\end{theorem}
%	\begin{proof}
%		\emph{Étape 1 (minorer la forme quadratique).} $H=\hess f(x^*)$ est symétrique définie positive : ses valeurs propres sont $>0$. Soit $\lambda>0$ la plus petite. Dans une base orthonormée de vecteurs propres, $h=\sum\alpha_iv_i$ et $h\T Hh=\sum\lambda_i\alpha_i^2\ge\lambda\sum\alpha_i^2=\lambda\norm h^2$ \emph{pour tout} $h$.
%		\emph{Étape 2 (Taylor).} Comme $\grad f(x^*)=0$ : $f(x^*+h)-f(x^*)=\frac12h\T Hh+r(h)$ avec $r(h)/\norm h^2\to0$ quand $h\to0$.
%		\emph{Étape 3 (le reste est petit).} Il existe $\rho>0$ tel que $\abs{r(h)}\le\frac\lambda4\norm h^2$ pour $\norm h<\rho$.
%		\emph{Étape 4 (conclusion).} Pour $0<\norm h<\rho$ :
%		$f(x^*+h)-f(x^*)\ge\frac\lambda2\norm h^2-\frac\lambda4\norm h^2=\frac\lambda4\norm h^2>0$.
%		Donc $x^*$ est un minimum local strict.
%		\emph{Pourquoi cette uniformité ?} Argumenter direction par direction ($d$ fixé, $t\to0$) ne suffit pas : le rayon $\rho$ pourrait dépendre de $d$ et tendre vers $0$. La constante $\lambda$ de l'étape 1 est indépendante de $h$.
%	\end{proof}

\begin{proof}
	\emph{Étape 1 (minorer la forme quadratique).}
	\(H=\hess f(x^*)\) est symétrique définie positive : ses valeurs propres sont \(>0\). Soit \(\lambda>0\) la plus petite. Dans une base orthonormée de vecteurs propres, \(h=\sum\alpha_i v_i\) et
	\[
	h\T Hh=\sum\lambda_i\alpha_i^2\ge\lambda\sum\alpha_i^2=\lambda\norm h^2
	\]
	\emph{pour tout} \(h\).
	
	\emph{Étape 2 (Taylor).}
	Comme \(\grad f(x^*)=0\) :
	\[
	f(x^*+h)-f(x^*)=\frac12 h\T Hh+r(h)
	\]
	avec \(r(h)/\norm h^2\to0\) quand \(h\to0\).
	
	\emph{Étape 3 (le reste est petit).}
	Il existe \(\rho>0\) tel que
	\[
	\abs{r(h)}\le\frac{\lambda}{4}\norm h^2
	\qquad \text{pour } \norm h<\rho.
	\]
	
	\emph{Étape 4 (conclusion).}
	Pour \(0<\norm h<\rho\) :
	\[
	f(x^*+h)-f(x^*)
	\ge \frac{\lambda}{2}\norm h^2-\frac{\lambda}{4}\norm h^2
	= \frac{\lambda}{4}\norm h^2>0.
	\]
	Donc \(x^*\) est un minimum local strict.
	
	\emph{Pourquoi cette uniformité ?}
	Argumenter direction par direction (\(d\) fixé, \(t\to0\)) ne suffit pas : le rayon \(\rho\) pourrait dépendre de \(d\) et tendre vers \(0\). La constante \(\lambda\) de l'étape 1 est indépendante de \(h\).
\end{proof}
	\begin{remarque}[Nature d'un point critique $x$]\label{rem:nature}
		$\hess f(x)\succ0\Rightarrow$ minimum local strict ; $\hess f(x)\prec0\Rightarrow$ maximum local strict ; valeurs propres de signes opposés $\Rightarrow$ ni min ni max (\textbf{point selle}, ou col). Si $\hess f(x)$ est semi-définie avec une valeur propre nulle, le test est \emph{inconclusif} : il faut étudier $f$ directement. Un minimum local est critique ; la réciproque est fausse.
	\end{remarque}
	
%	\newpage
%		\begin{theorem}[Minimum d'une fonction convexe]\label{thm:minconv}
%		Soit $f$ convexe sur un convexe $U$. 
%		\begin{enumerate}
%			\item[(a)]Tout minimum local est global, et l'ensemble des minimiseurs globaux est convexe. 
%			\item[(b)] Si $f$ est strictement convexe, il y a au plus un minimiseur. 
%			\item[(c)] Si $f$ est différentiable sur un ouvert convexe, $x^*$ est un minimum global $\iff\grad f(x^*)=0$.
%		\end{enumerate}
%	\end{theorem}
%	\begin{proof}[Preuve de (c)]
%		$\Rightarrow$ est la condition nécessaire du premier ordre (chap.~2). $\Leftarrow$ : par le théorème~\ref{thm:conv}, $f(y)\ge f(x^*)+\grad f(x^*)\T(y-x^*)=f(x^*)$.
%	\end{proof}
%	\begin{proof}[Preuves de (a), (b)]
%		(a) \emph{Hypothèse} : il existe $r>0$ tel que $f(x)\le f(z)$ pour tout $z\in U$ avec $\norm{z-x}<r$. Soit $y\in U$ quelconque ($y\ne x$).
%		\emph{Étape 1.} Choisissons $\theta\in]0,1[$ assez petit pour que $\theta\norm{y-x}<r$. Le point $z=x+\theta(y-x)=(1-\theta)x+\theta y$ est dans $U$ (convexe) et dans la boule : donc $f(x)\le f(z)$.
%		\emph{Étape 2.} La convexité donne $f(z)\le(1-\theta)f(x)+\theta f(y)$. Ainsi $f(x)\le(1-\theta)f(x)+\theta f(y)$.
%		\emph{Étape 3.} On retranche $(1-\theta)f(x)$ : $\theta f(x)\le\theta f(y)$, puis on divise par $\theta>0$ : $f(x)\le f(y)$. Comme $y$ est arbitraire, $x$ est un minimum global.
%		\emph{Convexité de l'ensemble des minimiseurs} : si $f(x_1)=f(x_2)=m=\min f$, alors $f(\lambda x_1+(1-\lambda)x_2)\le\lambda m+(1-\lambda)m=m$, donc c'est aussi un minimiseur.
%		(b) Si $x_1\ne x_2$ étaient deux minimiseurs, la stricte convexité donnerait $f(\frac{x_1+x_2}2)<\frac12m+\frac12m=m$, ce qui contredit $m=\min f$.
%	\end{proof}
	
	
	
	\begin{theorem}[Cas convexe]
		Si $f$ est convexe et différentiable sur $\R^n$ : $x^*$ est un minimum global $\iff\grad f(x^*)=0$. Si $f$ est strictement convexe, ce minimum est unique.
	\end{theorem}
	\begin{proof}
		$\Rightarrow$ : condition nécessaire du premier ordre (théorème~\ref{thm:cn}). $\Leftarrow$ : $f$ étant convexe, par le théorème~\ref{thm:conv}, $f(x)\ge f(x^*)+(x-x^*)\T\grad f(x^*)=f(x^*)$ pour tout $x$. 
		
%	\end{proof}
	\textbf{Sur un convexe $U$.} Si $x$ est minimum local de $f$ différentiable sur $U$ convexe, alors $\scal{\grad f(x)}{y-x}\ge0$ pour tout $y\in U$ (considérer $x+\theta(y-x)$, $\theta\to0^+$) ; si $U=x+F$ est un sous-espace affine, on a $\scal{\grad f(x)}{y-x}=0$, c'est-à-dire $\grad f(x)\in F^\perp$. Si $f$ est convexe, cette condition est aussi suffisante.
	
%	\begin{proof}[Preuve de la condition sur un convexe]
		Soit $y\in U$ et $h=y-x$.
		\emph{Étape 1.} Pour $\theta\in]0,1]$, $x+\theta h=(1-\theta)x+\theta y\in U$ (convexité). Si $\theta$ est assez petit, ce point est dans la boule où $x$ est minimal : $f(x+\theta h)-f(x)\ge0$.
		\emph{Étape 2.} On divise par $\theta>0$ : $\frac{f(x+\theta h)-f(x)}\theta\ge0$, puis $\theta\to0^+$ : $\scal{\grad f(x)}{h}\ge0$.
		\emph{Étape 3 (sous-espace affine).} Si $U=x+F$ et $y=x+h\in U$, alors $h\in F$, donc $-h\in F$ et $x-h\in U$. L'étape 2 appliquée à $-h$ donne $\scal{\grad f(x)}{-h}\ge0$, c'est-à-dire $\scal{\grad f(x)}{h}\le0$. Avec l'étape 2 : $\scal{\grad f(x)}{h}=0$.
	\end{proof}
	\begin{remarque}[Interprétation géométrique]
		Si $x$ est un minimum local sur un convexe $U$ et $\grad f(x)\ne0$, alors $\grad f(x)$ forme un angle aigu avec tout vecteur $y-x$, $y\in U$ : $-\grad f(x)$ ne pointe vers aucune direction admissible. Sur un sous-espace affine $x+F$, $\grad f(x)\in F^\perp$.
	\end{remarque}
	\begin{center}\small
		\begin{tabular}{|p{5.2cm}|p{5.6cm}|p{4.2cm}|}\hline
			\textbf{Hessienne en un point critique} & \textbf{Conclusion} & \textbf{Théorème}\\\hline
			définie positive & minimum local strict & CS (théorème~\ref{thm:cs})\\
			définie négative & maximum local strict & CS appliqué à $-f$\\
			valeurs propres de signes opposés & point selle (ni min ni max) & contraposée de CN (théorème~\ref{thm:cn})\\
			semi-définie, une valeur propre nulle & \emph{inconclusif} : étudier $f$ directement & --\\\hline
		\end{tabular}
	\end{center}
	
	\begin{remarque}[Trois erreurs à éviter]
		(1) \emph{``$\grad f(x)=0$ suffit pour avoir un minimum''} : faux. $f(x)=x^3$ en $0$, ou $f(x,y)=x^2-y^2$ en $(0,0)$ (selle). (2) \emph{``$\hess f\succeq0$ suffit''} : faux. $f(x)=x^3$ a $f'(0)=f''(0)=0$ sans extremum ; la condition d'ordre 2 est seulement nécessaire. (3) \emph{``$\hess f(x^*)\succ0$ est nécessaire''} : faux. $f(x)=x^4$ a un minimum en $0$ avec $f''(0)=0$ ; la condition définie positive est seulement suffisante.
	\end{remarque}
	\begin{retenir}
		Nécessaire : $\grad f=0$ et $\hess f\succeq0$. Suffisant : $\grad f=0$ et $\hess f\succ0$. Entre les deux (cas $\hess f\succeq0$ singulière), on ne conclut pas sans étude directe. Pour $f$ convexe, $\grad f=0$ est nécessaire \emph{et} suffisant pour un minimum global.
	\end{retenir}
	
	
	\textbf{Énoncés (travaux dirigés).} \emph{Exercice 1} : pour $f_1=3x_1^2+3x_2^2+3x_3^2-2x_1-10x_3-2x_1x_3$, $f_2=x_1^3+x_2^3-3x_1-12x_2+20$, $f_3=x_1^4+x_2^4-2(x_1-x_2)^2$ : étudier l'existence d'extrema, déterminer les points critiques par la condition du premier ordre, préciser leur nature (min/max, local/global). \emph{Exercice 2} : pour $f(x,y)=x^4+y^4-2(x-y)^2$, trouver $(\alpha,\beta)$ avec $f\ge\alpha\norm{(x,y)}_2^2+\beta$, en déduire l'existence d'un minimum, étudier la convexité, déterminer les points critiques. \emph{Exercice 3} : pour $Q$ symétrique définie positive, montrer $\scal{Qx}{x}\ge\delta\norm x^2$, puis que $f(x)=\frac12\scal{Qx}{x}-\scal{b}{x}$ est coercive et que $\min f$ a une solution unique à déterminer.
	\textbf{Méthode type} : (a) existence (continuité + coercivité) ; (b) $\grad f=0$ ; (c) nature par $\hess f$, convexité ou étude directe.
	\section{Exemples et exercices corrigés}
	\begin{exemple}\label{ex:cubique}
		$f(x,y)=x^3+y^3+3xy$. $\grad f=(3x^2+3y,\,3y^2+3x)=0$. La première équation donne $y=-x^2$ ; en reportant dans la seconde, $3x^4+3x=3x(x^3+1)=0$, donc $x=0$ ou $x=-1$ : points critiques $(0,0)$ et $(-1,-1)$. Hessienne $\begin{pmatrix}6x&3\\3&6y\end{pmatrix}$.
		En $(0,0)$ : $\begin{pmatrix}0&3\\3&0\end{pmatrix}$, valeurs propres $\pm3$ : \textbf{point selle}.
		En $(-1,-1)$ : $\begin{pmatrix}-6&3\\3&-6\end{pmatrix}$, valeurs propres $-9,-3<0$ : \textbf{maximum local}, $f(-1,-1)=1$ (non global : $f(t,0)=t^3\to\pm\infty$).
	\end{exemple}
	
	\begin{exemple}[$f_1(x_1,x_2,x_3)=3x_1^2+3x_2^2+3x_3^2-2x_1-10x_3-2x_1x_3$]
		$f_1=\frac12x\T Ax-b\T x$ avec $A=\begin{pmatrix}6&0&-2\\0&6&0\\-2&0&6\end{pmatrix}$, $b=(2,0,10)\T$. Mineurs : $6,\ 36,\ \det A=6(36)-2\cdot12=192>0$ : $A\succ0$. (Coercivité aussi directe : $2|x_1x_3|\le x_1^2+x_3^2$ donne $f_1\ge2x_1^2+3x_2^2+2x_3^2-2|x_1|-10|x_3|\to+\infty$.) Donc $f_1$ est strictement convexe et coercive (car $f_1(x)\ge\frac{\lambda_{\min}}2\norm x^2-\norm b\norm x$). L'unique minimum global est le point critique $Ax=b$ : $6x_1-2x_3=2$, $x_2=0$, $-2x_1+6x_3=10$, soit $x^*=(1,0,2)$ et $f_1(x^*)=-11$.
	\end{exemple}
	
	\begin{exemple}[$f_2(x_1,x_2)=x_1^3+x_2^3-3x_1-12x_2+20$]
		Points critiques : $x_1=\pm1$, $x_2=\pm2$. Hessienne $\operatorname{diag}(6x_1,6x_2)$.
		$(1,2)$ : définie positive, minimum local, $f_2=2$. $(-1,-2)$ : définie négative, maximum local, $f_2=38$. $(1,-2)$ et $(-1,2)$ : valeurs propres de signes opposés, points selles ($f_2=34$ et $6$).
		Aucun n'est global : $f_2(t,t)=2t^3-15t+20\to\mp\infty$ quand $t\to\mp\infty$ ($f_2$ n'est pas coercive), donc $f_2$ n'a ni minimum ni maximum global. Les valeurs $f_2(1,2)=2<f_2(-1,2)=6<f_2(1,-2)=34<f_2(-1,-2)=38$ illustrent qu'un minimum \emph{local} peut être plus petit que les autres valeurs critiques sans être global.
	\end{exemple}
	
	\begin{exemple}[$f(x,y)=x^4+y^4-2(x-y)^2$]\label{ex:f3}
		\emph{Existence.} $(x^2-3)^2\ge0$ donne $x^4\ge6x^2-9$ et $(x-y)^2\le2(x^2+y^2)$, d'où
		\[f(x,y)\ \ge\ 2\norm{(x,y)}^2-18\qquad(\alpha=2,\ \beta=-18).\]
		Donc $f$ est coercive ; continue, elle admet un minimum global.
		\emph{Points critiques.} $\partial_xf=4(x^3-x+y)$ et $\partial_yf=4(y^3+x-y)$. En additionnant : $x^3+y^3=0$, donc $y=-x$, puis $x^3-2x=0$ : $(0,0)$, $(\sqrt2,-\sqrt2)$, $(-\sqrt2,\sqrt2)$.
		\emph{Nature.} $\hess f=\begin{pmatrix}12x^2-4&4\\4&12y^2-4\end{pmatrix}$. En $(\pm\sqrt2,\mp\sqrt2)$ : $\begin{pmatrix}20&4\\4&20\end{pmatrix}$ (valeurs propres $16,24$) : minima locaux, $f=-8$. En $(0,0)$ : valeurs propres $0$ et $-8$ : le test est inconclusif, mais $f(t,t)=2t^4>0$ et $f(t,-t)=2t^4-8t^2<0$ (petit $t\ne0$) : point selle. Comme $f$ atteint son minimum global et que les seuls candidats sont les deux points $(\pm\sqrt2,\mp\sqrt2)$, ces points sont les minima globaux, de valeur $-8$.
		\emph{Détail des valeurs propres.} Pour $\begin{pmatrix}20&4\\4&20\end{pmatrix}$ : $\det(H-\lambda I)=(20-\lambda)^2-16=0$ donne $\lambda=20\pm4$. Pour $\begin{pmatrix}-4&4\\4&-4\end{pmatrix}$ : $(-4-\lambda)^2-16=\lambda^2+8\lambda=\lambda(\lambda+8)=0$ donne $\lambda=0$ ou $-8$.
		\emph{Convexité.} $\hess f(0,0)$ a une valeur propre négative : $f$ n'est pas convexe.
	\end{exemple}
	
	\begin{exemple}[Quadratique générale]
		Soit $Q$ symétrique définie positive et $f(x)=\frac12\scal{Qx}{x}-\scal{b}{x}$. \emph{Question 1.} Dans une base orthonormée de vecteurs propres $(v_i)$, $x=\sum\alpha_iv_i$ et $\scal{Qx}{x}=\sum\lambda_i\alpha_i^2\ge\lambda_{\min}\sum\alpha_i^2=\lambda_{\min}\norm x^2$ : on prend $\delta=\lambda_{\min}(Q)>0$ (les valeurs propres d'une matrice symétrique définie positive sont strictement positives). \emph{Question 2.} Alors $f(x)\ge\frac\delta2\norm x^2-\norm b\norm x\to+\infty$ (Cauchy--Schwarz) : $f$ est coercive. Sa hessienne est $Q\succ0$ : $f$ est strictement convexe. Existence (continue + coercive) et unicité (stricte convexité). Le gradient est $Qx-b$ (car $Q=Q\T$), donc le point critique vérifie $Qx=b$ ; $Q$ est inversible (valeurs propres non nulles), d'où l'unique solution $x^*=Q^{-1}b$, avec $f(x^*)=-\frac12\scal{b}{Q^{-1}b}$.
	\end{exemple}
	\chapter{Optimisation sous contraintes d'égalité : multiplicateurs de Lagrange}
	On étudie $(P_E)$ : optimiser $f(x)$ sous $g_i(x)=0$, $i=1,\dots,m$, avec $f:\R^n\to\R$, $g=(g_1,\dots,g_m):\R^n\to\R^m$, $m<n$. L'ensemble admissible est $\mathcal C=\{x\mid g(x)=0\}$.
	
	\textbf{Intuition géométrique.} Une seule contrainte : en un optimum, on ne peut ni monter ni descendre en restant sur $\mathcal C$, donc $\grad f(x^*)$ est orthogonal à $\mathcal C$. Or $\mathcal C$ est la ligne de niveau $0$ de $g$, donc $\grad g(x^*)$ lui est aussi orthogonal : les deux gradients sont colinéaires, soit $\grad f(x^*)+\lambda^*\grad g(x^*)=0$.
	
	\section{Condition nécessaire et suffisante}
	\begin{definition}
		Un point $x^*\in\mathcal C$ est \textbf{régulier} (la contrainte est \emph{qualifiée} en $x^*$) si $\grad g_1(x^*),\dots,\grad g_m(x^*)$ sont linéairement indépendants. Le \textbf{lagrangien} est $\M(x,\lambda)=f(x)+\sum_{i=1}^m\lambda_ig_i(x)$ ; les $\lambda_i$ sont les \textbf{multiplicateurs de Lagrange}.
	\end{definition}
	\begin{theorem}[Condition nécessaire, ordre 1]
		Soient $f,g_i$ de classe $\mathcal C^1$. Si $f|_{\mathcal C}$ présente un extremum local en $x^*$ et si $x^*$ est régulier, il existe $\lambda^*\in\R^m$ tel que $\grad\M(x^*,\lambda^*)=0$, c'est-à-dire
		\[\grad f(x^*)+\sum_i\lambda_i^*\grad g_i(x^*)=0,\qquad g(x^*)=0.\]
	\end{theorem}
%	\begin{proof}[Idée de la preuve \emph{(complément, absent du cours)}]
%		\emph{Étape 1 (espace tangent).} Posons $T=\ker D\T=\{h\mid\scal{\grad g_i(x^*)}{h}=0\ \forall i\}$ : ce sont les directions qui ne font pas varier $g$ au premier ordre.
%		\emph{Étape 2 (courbes sur $\mathcal C$).} Comme les $\grad g_i(x^*)$ sont indépendants (régularité), le théorème des fonctions implicites montre que $\mathcal C$ est, près de $x^*$, une surface de dimension $n-m$ dont $T$ est le plan tangent : pour tout $h\in T$ il existe une courbe $\mathcal C^1$ $\gamma:]-\eta,\eta[\to\mathcal C$ avec $\gamma(0)=x^*$, $\gamma'(0)=h$.
%		\emph{Étape 3 (dérivée nulle).} La fonction $\varphi(t)=f(\gamma(t))$ a un extremum local en $t=0$ (car $\gamma(t)\in\mathcal C$), donc $\varphi'(0)=\scal{\grad f(x^*)}{h}=0$.
%		\emph{Étape 4 (conclusion algébrique).} Ainsi $\grad f(x^*)\perp T$. Or $T=(\operatorname{Im}D)^\perp$ et donc $T^\perp=\operatorname{Im}D=\operatorname{Vect}(\grad g_1(x^*),\dots,\grad g_m(x^*))$. Il existe donc des réels $\lambda_i^*$ avec $\grad f(x^*)=-\sum\lambda_i^*\grad g_i(x^*)$.
%	\end{proof}
\begin{proof}
Notons \(\mathcal C=\{x\in\R^n : g(x)=0\}\) avec \(g=(g_1,\dots,g_m)\). Comme \(x^*\) est régulier, les vecteurs \(\grad g_1(x^*),\dots,\grad g_m(x^*)\) sont linéairement indépendants. L'espace tangent à \(\mathcal C\) en \(x^*\) est alors
\[
T_{x^*}\mathcal C=\ker Dg(x^*)
=\bigl\{d\in\R^n : \grad g_i(x^*)\T d=0,\ i=1,\dots,m\bigr\}.
\]

Soit \(d\in T_{x^*}\mathcal C\). Par régularité et le théorème des fonctions implicites, il existe une courbe \(\gamma:(-\varepsilon,\varepsilon)\to\mathcal C\) de classe \(\mathcal C^1\) telle que \(\gamma(0)=x^*\) et \(\gamma'(0)=d\). Comme \(f|_{\mathcal C}\) admet un extremum local en \(x^*\), la fonction \(t\mapsto f(\gamma(t))\) admet un extremum local en \(t=0\), donc
\[
0=\left.\frac{d}{dt}f(\gamma(t))\right|_{t=0}
=\grad f(x^*)\T\gamma'(0)
=\grad f(x^*)\T d.
\]
Ainsi \(\grad f(x^*)\) est orthogonal à \(T_{x^*}\mathcal C=\ker Dg(x^*)\). Par conséquent
\[
\grad f(x^*)\in\bigl(\ker Dg(x^*)\bigr)^{\perp}
=\operatorname{Im}\bigl(Dg(x^*)\T\bigr)
=\operatorname{Vect}\{\grad g_1(x^*),\dots,\grad g_m(x^*)\}.
\]
Il existe donc \(\lambda^*\in\R^m\) tel que
\[
\grad f(x^*)+\sum_{i=1}^m\lambda_i^*\grad g_i(x^*)=0.
\]
Enfin \(x^*\in\mathcal C\), donc \(g(x^*)=0\). C'est exactement \(\grad\M(x^*,\lambda^*)=0\) avec
\[
\M(x,\lambda)=f(x)+\sum_i\lambda_i g_i(x).
\]
\end{proof}
	\begin{remarque}
		Si $x^*$ n'est pas régulier, l'extremum peut exister sans qu'aucun multiplicateur n'existe (exemple~\ref{ex:nonqual}). Il faut alors traiter séparément les points non réguliers.
	\end{remarque}
	\begin{theorem}[Condition suffisante, ordre 2]\label{thm:l2}
		Soient $f,g_i$ de classe $\mathcal C^2$, $(x^*,\lambda^*)$ un point critique du lagrangien avec $x^*$ régulier. Si, pour tout $h\ne0$ tel que $\scal{\grad g_i(x^*)}{h}=0$ ($i=1,\dots,m$),
		\[h\T\hess_x\M(x^*,\lambda^*)h>0\quad(\text{resp.}<0),\]
		alors $x^*$ est un minimum (resp.\ maximum) local \emph{strict} du problème contraint. (Version faible : si $\hess_x\M(x^*,\lambda^*)\succ0$ sur tout $\R^n$, $x^*$ est un minimum.)
	\end{theorem}
	\textbf{Idée de la condition d'ordre 2.} Sur $\mathcal C$, on a $g=0$ donc $f=\M(\cdot,\lambda^*)$. Comme $\grad_x\M(x^*,\lambda^*)=0$, Taylor donne, pour $x\in\mathcal C$ proche de $x^*$, $f(x)-f(x^*)\approx\frac12(x-x^*)\T\hess_x\M(x^*,\lambda^*)(x-x^*)$, et le déplacement $x-x^*$ est (au premier ordre) \emph{tangent} à $\mathcal C$, c'est-à-dire dans $T$. C'est donc le signe de $\hess_x\M$ \emph{sur $T$} qui décide.
	
	
	\begin{proof}
		Notons \(\M(x,\lambda)=f(x)+\sum_{i=1}^m\lambda_i g_i(x)\) le lagrangien, et \(H:=\hess_x\M(x^*,\lambda^*)\). Comme \((x^*,\lambda^*)\) est un point critique,
		\[
		\grad_x\M(x^*,\lambda^*)=\grad f(x^*)+\sum_{i=1}^m\lambda_i^*\grad g_i(x^*)=0,
		\qquad g(x^*)=0.
		\]
		Rappelons que sur \(\mathcal C=\{g=0\}\), on a \(\M(x,\lambda^*)=f(x)\).
		
		\textbf{Espace tangent et constante de positivité.}
		L'espace tangent en \(x^*\) est
		\[
		T:=T_{x^*}\mathcal C=\bigl\{h\in\R^n : \grad g_i(x^*)\T h=0,\ i=1,\dots,m\bigr\}.
		\]
		C'est un sous-espace vectoriel fermé de \(\R^n\). Par hypothèse, la forme quadratique \(h\mapsto h\T Hh\) est strictement positive sur \(T\setminus\{0\}\). Comme la sphère unité de \(T\) est compacte (fermée bornée dans \(\R^n\)), il existe \(\mu>0\) tel que
		\[
		h\T Hh\ \ge\ \mu\,\norm h^2\qquad \forall h\in T.
		\]
		
		\textbf{Paramétrage local de la contrainte.}
		Comme \(x^*\) est régulier, le théorème des fonctions implicites (ou du rang) fournit un voisinage de \(0\) dans \(T\) et une application \(\varphi:T\cap B(0,\rho)\to T^\perp\) de classe \(\mathcal C^2\) tels que \(\varphi(0)=0\), \(D\varphi(0)=0\), et
		\[
		\mathcal C\cap B(x^*,r)=\bigl\{x^*+d+\varphi(d) : d\in T,\ \norm d<\rho\bigr\}
		\]
		pour \(r>0\) assez petit. En particulier \(\varphi(d)=O(\norm d^2)\) au voisinage de \(0\).
		
		\textbf{Développement de Taylor du lagrangien.}
		Puisque \(\M(\cdot,\lambda^*)\) est de classe \(\mathcal C^2\) et \(\grad_x\M(x^*,\lambda^*)=0\), le développement de Taylor--Young donne, pour \(x=x^*+d+\varphi(d)\in\mathcal C\) avec \(\norm d\) petit,
		\[
		\M(x,\lambda^*)-\M(x^*,\lambda^*)
		=\tfrac12\,(x-x^*)\T H(x-x^*)+o(\norm{x-x^*}^2).
		\]
		Or \(x-x^*=d+\varphi(d)\) avec \(\varphi(d)=O(\norm d^2)\), donc
		\[
		(x-x^*)\T H(x-x^*)=d\T Hd+o(\norm d^2),
		\qquad \norm{x-x^*}=\norm d+o(\norm d).
		\]
		Il vient
		\[
		\M(x,\lambda^*)-\M(x^*,\lambda^*)
		=\tfrac12\,d\T Hd+o(\norm d^2).
		\]
		Comme \(d\in T\) et \(d\T Hd\ge\mu\norm d^2\), on obtient, pour \(\norm d\) assez petit et \(d\neq 0\),
		\[
		\M(x,\lambda^*)-\M(x^*,\lambda^*)
		\ \ge\ \tfrac{\mu}{4}\norm d^2>0.
		\]
		
		\textbf{Conclusion.}
		Pour tout \(x\in\mathcal C\) suffisamment proche de \(x^*\) avec \(x\neq x^*\), on a \(\M(x,\lambda^*)=f(x)\) et \(\M(x^*,\lambda^*)=f(x^*)\), d'où
		\[
		f(x)-f(x^*)=\M(x,\lambda^*)-\M(x^*,\lambda^*)>0.
		\]
		Ainsi \(x^*\) est un minimum local strict de \(f|_{\mathcal C}\).
		
		Le cas du maximum s'obtient en appliquant le résultat à \(-f\) (ou à \(-H\)), ce qui donne la condition \(h\T Hh<0\) sur \(T\setminus\{0\}\).
		
		\textbf{Version faible.}
		Si \(H\succ0\) sur \(\R^n\) tout entier, alors en particulier \(h\T Hh>0\) pour tout \(h\in T\setminus\{0\}\), donc l'hypothèse est satisfaite et la conclusion tient.
	\end{proof}
	
	\textbf{Règle pratique.} Soit $A=\hess_x\M(x^*,\lambda^*)$ et $D=(\grad g_1(x^*)\ \cdots\ \grad g_m(x^*))\in\R^{n\times m}$. Les signes de $h\T Ah$ sur $\ker D\T$ sont ceux des racines du polynôme
	\[p(t)=\det\begin{pmatrix}A-tI_n&D\\ D\T&0\end{pmatrix}:\]
	toutes $>0\Rightarrow$ minimum ; toutes $<0\Rightarrow$ maximum ; sinon on ne conclut pas par ce critère.
	
	\begin{remarque}
		Le théorème faible demande $\hess_x\M\succ0$ sur \emph{tout} $\R^n$ (condition plus forte, donc moins souvent applicable) ; le théorème raffiné ne la demande que sur l'\textbf{espace tangent} $T=\ker D\T$, ce qui suffit puisque seules les directions tangentes à $\mathcal C$ sont admissibles au premier ordre.
	\end{remarque}
	\textbf{Méthode.} (1) Résoudre $\grad\M=0$. (2) Pour chaque point critique : vérifier la régularité, puis étudier $h\T Ah$ sur l'espace tangent (ou $p(t)$). (3) Examiner à part les points non réguliers.
	
	\section{Exemples corrigés}
	\begin{exemple}\label{ex:lag1}
		$f(x,y)=x^3y-x^2+3x+y^2+2y$ sous $g(x,y)=2x+y-1=0$. $\M=f+\lambda g$ :
		\[\partial_x\M=3x^2y-2x+3+2\lambda,\quad\partial_y\M=x^3+2y+2+\lambda,\quad\partial_\lambda\M=2x+y-1.\]
		Avec $y=1-2x$ et $\lambda=-x^3-4+4x$, la première équation devient $-8x^3+3x^2+6x-5=(x+1)(-8x^2+11x-5)=0$ ; le second facteur a un discriminant négatif. Unique point critique : $(x,y,\lambda)=(-1,3,-7)$. Contrainte régulière ($\grad g=(2,1)\ne0$).
		$A=\hess_x\M=\begin{pmatrix}6xy-2&3x^2\\3x^2&2\end{pmatrix}=\begin{pmatrix}-20&3\\3&2\end{pmatrix}$, $D=(2,1)\T$. Espace tangent : $h=(1,-2)$, $h\T Ah=-20-12+8=-24<0$. Via la règle pratique,
		$p(t)=\det\begin{pmatrix}-20-t&3&2\\3&2-t&1\\2&1&0\end{pmatrix}=24+5t$, racine $-\frac{24}5<0$.
		Donc $(-1,3)$ est un \textbf{maximum local} et $f(-1,3)=8$. (Ici $A$ a une valeur propre de chaque signe : le théorème faible ne permettait pas de conclure ; le raffinement de l'espace tangent le permet.)
	\end{exemple}
	\begin{exemple}\label{ex:sphere}
		Optimiser $f=x+y+z$ sous $x^2+y^2+z^2=r^2$. $\grad\M=0$ : $1+2\lambda x=1+2\lambda y=1+2\lambda z=0$, donc $x=y=z=-\frac1{2\lambda}$ et $3x^2=r^2$ : deux points $x=\pm\frac r{\sqrt3}(1,1,1)$, $\lambda=\mp\frac{\sqrt3}{2r}$, réguliers car $\grad g=2x\ne0$. Ici $\hess_x\M=2\lambda I_3$ :
		$A$ (en $+\frac r{\sqrt3}$) a pour valeurs propres $-\frac{\sqrt3}r<0$ : \textbf{maximum}, $f=\sqrt3\,r$ ;
		$B$ (en $-\frac r{\sqrt3}$) a pour valeurs propres $+\frac{\sqrt3}r>0$ : \textbf{minimum}, $f=-\sqrt3\,r$.
	\end{exemple}
	\begin{exemple}
		Extrema de $f=xy$ sur $x^2+y^2=1$. $\M=xy+\lambda(x^2+y^2-1)$ : $y+2\lambda x=0$, $x+2\lambda y=0$ donne $x(1-4\lambda^2)=0$ ; $x=0$ est exclu (alors $y=0$), donc $\lambda=\pm\frac12$.
		$\lambda=-\frac12$ : $y=x=\pm\frac1{\sqrt2}$, $f=\frac12$ ; $\hess_x\M=\begin{pmatrix}-1&1\\1&-1\end{pmatrix}$, tangent $h=(1,-1)$ : $h\T Ah=-4<0$, \textbf{maxima}.
		$\lambda=\frac12$ : $y=-x=\mp\frac1{\sqrt2}$, $f=-\frac12$ ; $\hess_x\M=\begin{pmatrix}1&1\\1&1\end{pmatrix}$, tangent $h=(1,1)$ : $h\T Ah=4>0$, \textbf{minima}.
	\end{exemple}
	\begin{exemple}[Cas simple]
		$\min x^2+y^2$ sous $x+y=1$. $\M=x^2+y^2+\lambda(x+y-1)$ : $2x+\lambda=0$, $2y+\lambda=0$, $x+y=1$ donnent $x=y=\frac12$, $\lambda=-1$ ; $\grad g=(1,1)\ne0$ (régulier) ; $\hess_x\M=2I_2\succ0$ : minimum, $f=\frac12$ (la distance de l'origine à la droite est $\frac1{\sqrt2}$ et $f$ est son carré). Le minimum est global car $f$ est coercive sur la droite.
	\end{exemple}
	\begin{exemple}[Contraintes multiples]
		$\min x^2+y^2+z^2$ sous $x+y+z=1$, $x-y=0$. Gradients des contraintes $(1,1,1)$ et $(1,-1,0)$ : indépendants. Stationnarité : $2x+\lambda+\mu=0$, $2y+\lambda-\mu=0$, $2z+\lambda=0$. On obtient $\mu=0$, $x=y=z=\frac13$. Comme $\hess_x\M=2I_3\succ0$, c'est le minimum, de valeur $\frac13$.
	\end{exemple}
	Vérifications des points de l'exemple précédent. Pour $f=xy$, la contrainte $\grad g=2(x,y)\ne0$ sur le cercle : tous les points sont réguliers, et comme le cercle est compact, $f$ y atteint un maximum et un minimum (Weierstrass) : ce sont nécessairement les points critiques trouvés ; la comparaison des valeurs $\pm\frac12$ confirme que les maxima valent $\frac12$ et les minima $-\frac12$ (globaux).
	\begin{exemple}[Contrainte non qualifiée]\label{ex:nonqual}
		$f(x,y)=x$ sur $\mathcal C=\{x^2+y^2=0\}=\{(0,0)\}$. En $(0,0)$ : $\grad g=0$ (point non régulier) et $\partial_x\M=1+2\lambda x=1\ne0$ pour tout $\lambda$ : aucun multiplicateur n'existe, bien que $(0,0)$ soit trivialement l'optimum. La méthode échoue : on raisonne directement sur $\mathcal C$.
	\end{exemple}
	\begin{exemple}[Contraintes redondantes]
		Pour $x+y+z=1$ et $2x+2y+2z=2$ les gradients sont colinéaires : non régulier. Seule la combinaison $\lambda+2\mu$ intervient dans $\M$ : les multiplicateurs ne sont pas uniques. On supprime la contrainte redondante.
	\end{exemple}
	\begin{retenir}
		Lagrange : $\grad f+\sum\lambda_i\grad g_i=0$ \emph{si} les $\grad g_i$ sont indépendants. Le test de second ordre se fait \emph{sur l'espace tangent} $\ker D\T$, pas sur tout $\R^n$. Un point critique non régulier doit être traité à part.
	\end{retenir}
	\chapter{Conditions de Karush--Kuhn--Tucker}
	Cadre : $(P)$ minimiser $f(x)$ sous $h_i(x)=0$ ($i=1,\dots,p$) et $g_j(x)\le0$ ($j=1,\dots,q$), avec $f,h_i,g_j$ de classe $\mathcal C^1$. L'ensemble admissible est $U=\{x\mid h(x)=0,\ g(x)\le0\}$. Pour $x\in U$, $I(x)=\{j\mid g_j(x)=0\}$ est l'ensemble des contraintes d'inégalité \textbf{actives} (saturées) ; les autres sont \emph{inactives}.
	
	\section{Conditions nécessaires non qualifiées (Fritz John)}
	\begin{theorem}[Fritz John]\label{thm:fj}
		Si $x^*\in U$ est un minimum local de $(P)$, il existe $\mu_0^*\ge0$, $\lambda^*\in\R^p$, $\mu^*\in\R^q_+$, \emph{non tous nuls}, tels que
		\[\mu_0^*\grad f(x^*)+\sum_{i=1}^p\lambda_i^*\grad h_i(x^*)+\sum_{j=1}^q\mu_j^*\grad g_j(x^*)=0,\qquad\mu_j^*g_j(x^*)=0\ \ (\forall j).\]
	\end{theorem}
	\textbf{Interprétation.}\\
% Les directions admissibles en $x^*$ vérifient $\scal{\grad h_i}{d}=0$ et $\scal{\grad g_j}{d}\le0$ ($j\in I(x^*)$). Au minimum local, $\scal{\grad f(x^*)}{d}\ge0$ pour toute direction admissible : aucune descente n'est possible. Par un argument de séparation (lemme de Farkas, théorème des alternatives de Motzkin), $\grad f(x^*)$ est alors dans le cône engendré par les opposés des normales aux contraintes actives (avec coefficient $\mu_0^*$ potentiellement nul), d'où la combinaison linéaire. La preuve rigoureuse repose sur ces lemmes ; elle est admise ici.
%	
%La relation $\mu_j^*g_j(x^*)=0$ est la \textbf{complémentarité} : $\mu_j^*=0$ dès que $g_j(x^*)<0$ ; seules les contraintes actives interviennent. Les $\lambda_i^*$ sont de signe quelconque, les $\mu_j^*$ sont $\ge0$ (une inégalité ne « pousse » que dans un sens).

Les directions admissibles en \(x^*\) vérifient
\(\scal{\grad h_i}{d}=0\) et \(\scal{\grad g_j}{d}\le0\) (\(j\in I(x^*)\)).
Au minimum local, \(\scal{\grad f(x^*)}{d}\ge0\) pour toute direction admissible :
aucune descente n'est possible. Par un argument de séparation
(lemme de Farkas, théorème des alternatives de Motzkin), \(\grad f(x^*)\)
est alors dans le cône engendré par les opposés des normales aux contraintes
actives (avec coefficient \(\mu_0^*\) potentiellement nul), d'où la combinaison
linéaire. La preuve rigoureuse repose sur ces lemmes.

La relation \(\mu_j^*g_j(x^*)=0\) est la \textbf{complémentarité} :
\(\mu_j^*=0\) dès que \(g_j(x^*)<0\) ; seules les contraintes actives
interviennent. Les \(\lambda_i^*\) sont de signe quelconque, les \(\mu_j^*\)
sont \(\ge0\) (une inégalité ne « pousse » que dans un sens).
	
	\begin{remarque}[]\quad
	\begin{enumerate}
		\item[(i)] Les multiplicateurs ne sont pas uniques (on peut les multiplier par un réel $>0$). 
		\item[(ii)] Si $\mu_0^*>0$ on divise par $\mu_0^*$ et on obtient les conditions KKT usuelles. 
		\item[(iii)] Fritz John caractérise des minima \emph{locaux} ; le global exige des hypothèses supplémentaires (convexité).
	\end{enumerate}
	\end{remarque}
	\textbf{Limite.} Si $\mu_0^*=0$, la relation ne contient plus $\grad f$ : elle ne dit rien sur l'objectif. Exemple : $\min x$ sur $x^2+y^2=0$ (donc $x^*=(0,0)$) : $\grad h(0,0)=0$ et $(\mu_0,\lambda)=(0,1)$ convient. On cherche donc des \emph{qualifications} garantissant $\mu_0^*>0$.
	
	\section{Qualifications des contraintes}
	Une mauvaise description de $U$ peut créer des dépendances entre gradients sans changer $U$ (par exemple remplacer $h=0$ par $(h,-h)=0$, ou doubler une contrainte), et rend Fritz John inutile.
	\begin{definition}[QCMF -- Mangasarian--Fromovitz]
		$x^*\in U$ vérifie la QCMF si (i) les $\grad h_i(x^*)$ sont linéairement indépendants, et (ii) il existe $d\in\R^n$ tel que $\scal{\grad h_i(x^*)}{d}=0$ pour tout $i$ et $\scal{\grad g_j(x^*)}{d}<0$ pour tout $j\in I(x^*)$.
	\end{definition}
	\begin{definition}[QCKT / LICQ]
		$x^*\in U$ vérifie la LICQ si les vecteurs $\grad h_i(x^*)$ ($i\le p$) et $\grad g_j(x^*)$ ($j\in I(x^*)$) sont linéairement indépendants.
	\end{definition}
	\begin{definition}[Slater -- problème convexe]
		$f,g_j$ convexes, $h_i$ affines (à gradients indépendants) : il existe $\bar x$ avec $h(\bar x)=0$ et $g_j(\bar x)<0$ pour tout $j$.
	\end{definition}
	\begin{proposition}\label{prop:qual}
		(a) LICQ $\Rightarrow$ QCMF. (b) Slater $\Rightarrow$ QCMF en tout point admissible d'un problème convexe. En général : LICQ $\Rightarrow$ MFCQ $\Rightarrow$ CPLD $\Rightarrow$ ACQ, et CRCQ $\Rightarrow$ CPLD.
	\end{proposition}
	\begin{proof}
		\textbf{(a) LICQ $\Rightarrow$ QCMF.}
		La condition (i) de la QCMF est déjà contenue dans la LICQ.
		Pour la condition (ii), notons $M$ la matrice dont les lignes sont les gradients actifs
		\[
		\nabla h_i(x^*)^\top \quad (i=1,\dots,p)
		\quad\text{et}\quad
		\nabla g_j(x^*)^\top \quad (j\in I(x^*)).
		\]
		Cette matrice possède $r=p+|I(x^*)|$ lignes linéairement indépendantes, donc elle est de rang $r$.
		L'application linéaire $d\mapsto Md$ de $\mathbb{R}^n$ dans $\mathbb{R}^r$ est par conséquent \emph{surjective}.
		On peut donc choisir un vecteur $d$ tel que
		\[
		Md=(0,\dots,0,-1,\dots,-1)^\top,
		\]
		c'est-à-dire des zéros sur les $p$ lignes correspondant aux égalités et des $-1$ sur les lignes correspondant aux inégalités actives.
		On a alors
		\[
		\langle\nabla h_i(x^*),d\rangle=0
		\quad\text{et}\quad
		\langle\nabla g_j(x^*),d\rangle=-1<0
		\quad\text{pour tout }j\in I(x^*),
		\]
		ce qui est exactement la condition (ii) de la QCMF.
		
		\textbf{(b) Slater $\Rightarrow$ QCMF (cas convexe).}
		Soit $x^*$ un point admissible et $\bar{x}$ un point strictement admissible. Posons $d=\bar{x}-x^*$.
		
		\begin{itemize}
			\item Pour toute contrainte active $j\in I(x^*)$, la convexité de $g_j$ (caractérisation du premier ordre) donne
			\[
			g_j(\bar{x})\ge g_j(x^*)+\langle\nabla g_j(x^*),d\rangle.
			\]
			Comme $g_j(x^*)=0$ et $g_j(\bar{x})<0$, on en déduit
			\[
			\langle\nabla g_j(x^*),d\rangle\le g_j(\bar{x})<0.
			\]
			
			\item Pour toute contrainte d'égalité (affine), on a
			\[
			\langle\nabla h_i(x^*),d\rangle=h_i(\bar{x})-h_i(x^*)=0.
			\]
		\end{itemize}
		
		La condition (ii) est donc satisfaite. La condition (i) est quant à elle déjà incluse dans l'hypothèse de Slater.
	\end{proof}
	
	\begin{remarque}
		La CRCQ (rang constant de la matrice des gradients actifs au voisinage de $x^*$) est plus faible que la LICQ. En général elle \textbf{n'implique pas} la QCMF, mais elle implique la CPLD.
		En pratique, la LICQ se vérifie simplement par un calcul de rang de la matrice des gradients actifs.
	\end{remarque}
	
	
%	\begin{exemple}[Vérifier une qualification]
%		(i) $\min x^2+y^2$ sous $x+y-1\le0$, au point actif $(\frac12,\frac12)$ : $\grad g=(1,1)$ seul, indépendant : LICQ. Si l'on ajoute la contrainte redondante $2x+2y-2\le0$, les deux gradients actifs $(1,1)$ et $(2,2)$ sont dépendants : la LICQ échoue alors que l'ensemble $U$ est inchangé ; le rang de la matrice des gradients actifs reste constant égal à $1$ (CRCQ vérifiée).
%		(ii) Deux égalités redondantes $x_1+x_2=0$ et $2x_1+2x_2=0$ : gradients dépendants, multiplicateurs non uniques et $\mu_0^*=0$ possible même en un point optimal.
%		(iii) $\min\norm x^2$ sous $Ax-b\le0$ : si $\{x\mid Ax<b\}\ne\emptyset$, Slater est vérifiée ; si les inégalités déguisent une égalité (ensemble strict vide), elle échoue.
%	\end{exemple}
	\begin{exemple}[Vérifier une qualification de contraintes]
		\begin{enumerate}
			\item[(i)] Considérons le problème
			\[
			\min x^2+y^2 \quad\text{sous}\quad x+y-1\le 0.
			\]
			Au point actif $\bigl(\frac12,\frac12\bigr)$, le seul gradient actif est $\nabla g=(1,1)$, qui est non nul : la \textbf{LICQ} est donc vérifiée.
			
			Si l'on ajoute la contrainte redondante $2x+2y-2\le 0$, les deux gradients actifs $(1,1)$ et $(2,2)$ deviennent linéairement dépendants. La LICQ échoue, alors que l'ensemble admissible $U$ reste inchangé. En revanche, le rang de la matrice des gradients actifs reste constant égal à $1$ : la \textbf{CRCQ} est vérifiée.
			
			\item[(ii)] Deux égalités redondantes
			\[
			x_1+x_2=0,\qquad 2x_1+2x_2=0
			\]
			donnent des gradients linéairement dépendants. Les multiplicateurs ne sont plus uniques et il est possible d'avoir $\mu_0^*=0$ même en un point optimal.
			
			\item[(iii)] Pour le problème
			\[
			\min\|x\|^2 \quad\text{sous}\quad Ax-b\le 0,
			\]
			si l'ensemble $\{x\mid Ax<b\}$ est non vide, la condition de \textbf{Slater} est vérifiée.
			En revanche, si les inégalités déguisent en réalité une égalité (l'intérieur relatif de l'ensemble admissible est vide), Slater échoue.
		\end{enumerate}
	\end{exemple}
	\section{Théorème de KKT}
	\begin{definition}
		Le \textbf{lagrangien} de $(P)$ est $\M(x,\lambda,\mu)=f(x)+\sum_i\lambda_ih_i(x)+\sum_j\mu_jg_j(x)$.
	\end{definition}
%	\begin{theorem}[Karush--Kuhn--Tucker]\label{thm:kkt}
%		Soit $x^*$ un minimum local de $(P)$ vérifiant la QCMF (en particulier la LICQ). Il existe $\lambda^*\in\R^p$, $\mu^*\in\R^q$ tels que :
%		\begin{enumerate}
%			\item \textbf{Stationnarité} : $\grad_x\M(x^*,\lambda^*,\mu^*)=\grad f(x^*)+\sum\lambda_i^*\grad h_i(x^*)+\sum\mu_j^*\grad g_j(x^*)=0$ ;
%			\item \textbf{Admissibilité primale} : $h(x^*)=0$, $g(x^*)\le0$ ;
%			\item \textbf{Admissibilité duale} : $\mu^*\ge0$ ;
%			\item \textbf{Complémentarité} : $\mu_j^*g_j(x^*)=0$ pour tout $j$.
%		\end{enumerate}
%	\end{theorem}
%	\begin{proof}
%		\textbf{Stratégie.} Fritz John fournit $(\mu_0^*,\lambda^*,\mu^*)\ne0$ avec $\mu_0^*\ge0$. Si l'on prouve $\mu_0^*>0$, on divise toute l'égalité par $\mu_0^*$ : $\lambda^*/\mu_0^*$ et $\mu^*/\mu_0^*$ sont les multiplicateurs KKT (la complémentarité et $\mu\ge0$ sont conservées). Il reste à montrer que $\mu_0^*=0$ est impossible.
%		
%		\textbf{Raisonnement par l'absurde} : supposons $\mu_0^*=0$. Comme $\mu_j^*g_j(x^*)=0$, on a $\mu_j^*=0$ pour les contraintes inactives, et l'égalité de Fritz John devient
%		\[(E)\qquad\sum_{i=1}^p\lambda_i^*\grad h_i(x^*)+\sum_{j\in I(x^*)}\mu_j^*\grad g_j(x^*)=0,\]
%		avec $(\lambda^*,\mu^*)\ne0$ (puisque $\mu_0^*=0$ et que le vecteur total est non nul).
%		\emph{Cas A : tous les $\mu_j^*$ sont nuls.} Alors $\lambda^*\ne0$ et (E) devient $\sum\lambda_i^*\grad h_i(x^*)=0$ : les $\grad h_i(x^*)$ sont liés, ce qui contredit (i) de la QCMF.
%		\emph{Cas B : un $\mu_{j_0}^*>0$.} Soit $d$ donné par (ii). On fait le produit scalaire de (E) avec $d$ : les termes d'égalité s'annulent ($\scal{\grad h_i}{d}=0$), il reste $0=\sum_{j\in I}\mu_j^*\scal{\grad g_j(x^*)}{d}$. Or chaque terme est $\le0$ ($\mu_j^*\ge0$ et $\scal{\grad g_j}{d}<0$) et le terme $j_0$ est $<0$ : la somme est $<0$, contradiction.
%		Dans les deux cas on aboutit à une contradiction : $\mu_0^*>0$.
%	\end{proof}
	
	\begin{theorem}[Karush--Kuhn--Tucker]\label{thm:kkt}
		Soit \(x^*\) un minimum local de \((P)\) vérifiant la QCMF (en particulier la
		LICQ). Il existe \(\lambda^*\in\R^p\), \(\mu^*\in\R^q_+\) tels que :
		\begin{enumerate}
			\item \textbf{Stationnarité} :
			\[
			\grad_x\M(x^*,\lambda^*,\mu^*)
			=\grad f(x^*)+\sum_{i=1}^p\lambda_i^*\grad h_i(x^*)
			+\sum_{j=1}^q\mu_j^*\grad g_j(x^*)=0\,;
			\]
			\item \textbf{Admissibilité primale} : \(h(x^*)=0\), \(g(x^*)\le0\) ;
			\item \textbf{Admissibilité duale} : \(\mu^*\ge0\) ;
			\item \textbf{Complémentarité} : \(\mu_j^*g_j(x^*)=0\) pour tout \(j\).
		\end{enumerate}
	\end{theorem}

	\begin{proof}
		\textbf{Étape 1 -- Condition de Fritz John.}
		Sous la seule hypothèse de différentiabilité, le théorème de Fritz John assure
		l'existence de multiplicateurs
		\[
		(\mu_0^*,\lambda^*,\mu^*)\neq0,\qquad
		\mu_0^*\ge0,\quad \mu^*\ge0,
		\]
		tels que
		\[
		\mu_0^*\grad f(x^*)+\sum_{i=1}^p\lambda_i^*\grad h_i(x^*)
		+\sum_{j=1}^q\mu_j^*\grad g_j(x^*)=0,
		\]
		avec la complémentarité \(\mu_j^*g_j(x^*)=0\) pour tout \(j\).
		
		\textbf{Étape 2 -- Il suffit de montrer \(\mu_0^*>0\).}
		Si \(\mu_0^*>0\), on divise toute l'égalité par \(\mu_0^*\) : les multiplicateurs
		\(\lambda^*/\mu_0^*\) et \(\mu^*/\mu_0^*\) vérifient alors exactement les
		conditions KKT (stationnarité, \(\mu\ge0\), et complémentarité, préservée par
		homogénéité). Le travail consiste donc à exclure le cas \(\mu_0^*=0\).
		
		\textbf{Étape 3 -- Raisonnement par l'absurde : supposons \(\mu_0^*=0\).}
		L'égalité de Fritz John devient alors
		\[
		(E)\qquad
		\sum_{i=1}^p\lambda_i^*\grad h_i(x^*)
		+\sum_{j\in I(x^*)}\mu_j^*\grad g_j(x^*)=0,
		\]
		et le vecteur \((\lambda^*,\mu^*_{I(x^*)})\) est non nul (puisque le
		multiplicateur total l'est, avec \(\mu_0^*=0\)).
		
		\emph{Cas A : tous les \(\mu_j^*\) sont nuls.}
		Alors \(\lambda^*\neq0\), et \((E)\) s'écrit
		\(\sum_i\lambda_i^*\grad h_i(x^*)=0\). Les gradients \(\grad h_i(x^*)\) sont
		donc linéairement dépendants, ce qui contredit directement la \textbf{LICQ}
		(indépendance linéaire des gradients des contraintes actives et des égalités).
		
		\emph{Cas B : il existe \(j_0\in I(x^*)\) tel que \(\mu_{j_0}^*>0\).}
		On utilise la seconde partie de la LICQ (ou plus généralement la QCMF) : il
		existe une direction \(d\) telle que
		\[
		\grad h_i(x^*)\T d=0\quad(i=1,\dots,p),\qquad
		\grad g_j(x^*)\T d<0\quad(j\in I(x^*)).
		\]
		En faisant le produit scalaire de \((E)\) avec \(d\), les termes d'égalité
		s'annulent, et il reste
		\[
		0=\sum_{j\in I(x^*)}\mu_j^*\underbrace{\grad g_j(x^*)\T d}_{<0}.
		\]
		Chaque terme de la somme est \(\le0\) (car \(\mu_j^*\ge0\)), et le terme \(j_0\)
		est strictement négatif (car \(\mu_{j_0}^*>0\)). La somme est donc strictement
		négative, ce qui contredit \(0\).
		
		Dans les deux cas, on aboutit à une contradiction. Donc \(\mu_0^*\) ne peut pas
		être nul, et l'on a bien \(\mu_0^*>0\).
		
		\textbf{Conclusion.}
		En divisant par \(\mu_0^*>0\), on obtient des multiplicateurs
		\(\lambda^*\in\R^p\) et \(\mu^*\in\R^q_+\) vérifiant les quatre conditions KKT :
		\[
		\grad f(x^*)+\sum_{i=1}^p\lambda_i^*\grad h_i(x^*)
		+\sum_{j=1}^q\mu_j^*\grad g_j(x^*)=0,
		\]
		\[
		h(x^*)=0,\quad g(x^*)\le0,\quad \mu^*\ge0,\quad
		\mu_j^*g_j(x^*)=0\ \ (\forall j).
		\]
	\end{proof}
%	\begin{theorem}[KKT suffisant, cas convexe]\label{thm:kktconv}
%		Si $(P)$ est convexe ($f,g_j$ convexes, $h_i$ affines) et si $(x^*,\lambda^*,\mu^*)$ vérifie les quatre conditions KKT, alors $x^*$ est un minimum \emph{global}. Aucune qualification n'est nécessaire pour cette implication.
%	\end{theorem}
%	\begin{proof}
%		\emph{Étape 1 : le lagrangien est convexe en $x$.} $\ell(x)=f(x)+\sum\lambda_i^*h_i(x)+\sum\mu_j^*g_j(x)$ est somme de $f$ (convexe), de fonctions affines $\lambda_i^*h_i$ (convexes) et de $\mu_j^*g_j$ (convexes car $\mu_j^*\ge0$ et $g_j$ convexe).
%		\emph{Étape 2 : $x^*$ minimise $\ell$.} La stationnarité dit $\grad\ell(x^*)=0$ ; pour une fonction convexe différentiable, un point critique est un minimum global (chapitre~1). Donc $\ell(x)\ge\ell(x^*)$ pour tout $x$.
%		\emph{Étape 3 : valeur de $\ell$ en $x^*$.} $h(x^*)=0$ et $\mu_j^*g_j(x^*)=0$ (complémentarité), donc $\ell(x^*)=f(x^*)$.
%		\emph{Étape 4 : comparaison avec un point admissible $x$.} $h(x)=0$ et, comme $\mu_j^*\ge0$, $g_j(x)\le0$, on a $\sum\mu_j^*g_j(x)\le0$ ; donc $\ell(x)=f(x)+\sum\mu_j^*g_j(x)\le f(x)$.
%		\emph{Conclusion.} Pour tout $x$ admissible : $f(x)\ge\ell(x)\ge\ell(x^*)=f(x^*)$.
%	\end{proof}
	
	\begin{theorem}[KKT suffisant, cas convexe]\label{thm:kktconv}
		Si \((P)\) est convexe (\(f,g_j\) convexes, \(h_i\) affines) et si
		\((x^*,\lambda^*,\mu^*)\) vérifie les quatre conditions KKT, alors \(x^*\)
		est un minimum \emph{global}. Aucune qualification n'est nécessaire pour
		cette implication.
	\end{theorem}
	
	\begin{proof}
		\emph{Étape 1 : le lagrangien est convexe en \(x\).}
		\[
		\ell(x)=f(x)+\sum_{i=1}^p\lambda_i^*h_i(x)+\sum_{j=1}^q\mu_j^*g_j(x)
		\]
		est somme de \(f\) (convexe), de fonctions affines \(\lambda_i^*h_i\)
		(convexes) et de \(\mu_j^*g_j\) (convexes car \(\mu_j^*\ge0\) et \(g_j\)
		convexe).
		
		\emph{Étape 2 : \(x^*\) minimise \(\ell\).}
		La stationnarité dit \(\grad\ell(x^*)=0\) ; pour une fonction convexe
		différentiable, un point critique est un minimum global (chapitre~1). Donc
		\(\ell(x)\ge\ell(x^*)\) pour tout \(x\).
		
		\emph{Étape 3 : valeur de \(\ell\) en \(x^*\).}
		\(h(x^*)=0\) et \(\mu_j^*g_j(x^*)=0\) (complémentarité), donc
		\(\ell(x^*)=f(x^*)\).
		
		\emph{Étape 4 : comparaison avec un point admissible \(x\).}
		Comme \(h(x)=0\) et, \(\mu_j^*\ge0\), \(g_j(x)\le0\), on a
		\(\sum_j\mu_j^*g_j(x)\le0\) ; donc
		\[
		\ell(x)=f(x)+\sum_{j=1}^q\mu_j^*g_j(x)\le f(x).
		\]
		
		\emph{Conclusion.}
		Pour tout \(x\) admissible :
		\[
		f(x)\ge\ell(x)\ge\ell(x^*)=f(x^*).
		\]
	\end{proof}
	Avec Slater, les conditions KKT sont donc \textbf{nécessaires et suffisantes} pour un problème convexe. Hors cas convexe, KKT n'est que nécessaire (sous qualification) : on identifie les candidats puis on conclut par compacité/coercivité ou par une condition de second ordre.
	
	\begin{retenir}
		\textbf{Résoudre un problème avec KKT : mode d'emploi.} (1) Écrire les contraintes sous la forme $h_i=0$, $g_j\le0$ et le lagrangien. (2) Vérifier une qualification (LICQ, ou Slater si convexe). (3) Écrire les 4 conditions. (4) Discuter les cas ``$\mu_j=0$ ou $g_j=0$'' pour chaque $j$ ($2^q$ cas au plus). (5) Garder les points admissibles avec $\mu\ge0$. (6) Conclure : convexe $\Rightarrow$ ces points sont des minima globaux ; sinon comparer les valeurs de $f$ (si un minimum existe, par compacité ou coercivité).
	\end{retenir}
	
	\section{Exemples et exercices corrigés}
	\begin{exemple}[Une contrainte]
		$\min(x-2)^2+(y-1)^2$ sous $g=x+y-2\le0$. Problème convexe, Slater vérifiée ($(0,0)$), $\grad g=(1,1)\ne0$ : LICQ. KKT : $2(x-2)+\mu=0$, $2(y-1)+\mu=0$, $\mu(x+y-2)=0$, $\mu\ge0$.
		$\mu=0$ donne $(2,1)$, non admissible ($g=1>0$). $\mu>0$ : $x+y=2$ et $x-2=y-1$, soit $(x,y)=(\frac32,\frac12)$, $\mu=1\ge0$. Par le théorème~\ref{thm:kktconv}, c'est l'unique minimum global ; $f=\frac12$ (point de tangence de la droite avec un cercle centré en $(2,1)$).
	\end{exemple}
	\begin{center}
		\begin{tikzpicture}[scale=1.5]
			\fill[blue!15] (-0.3,-0.3)--(2.3,-0.3)--(-0.3,2.3)--cycle;
			\draw[->](-.3,0)--(3,0) node[right]{$x$};\draw[->](0,-.3)--(0,2.8) node[above]{$y$};
			\draw[thick](-.3,2.3)--(2.3,-.3) node[right]{\small$x+y=2$};
			\foreach \r in {0.5,0.7071,1}\draw[red!50,dashed](2,1) circle (\r);
			\fill(2,1) circle (1pt) node[above right]{\small$(2,1)$};\fill[red](1.5,.5) circle (1.5pt) node[below left]{\small$(\frac32,\frac12)$};
		\end{tikzpicture}
	\end{center}
	\begin{exemple}[Un exercice KKT simple]
		$\min x^2+y^2$ sous $x\ge1$, $y\ge2$, soit $g_1=1-x\le0$, $g_2=2-y\le0$ : $\M=x^2+y^2+\mu_1(1-x)+\mu_2(2-y)$. Convexe (contraintes affines), Slater en $(2,3)$ ; $\grad g_1=(-1,0)$, $\grad g_2=(0,-1)$ indépendants : LICQ en tout point. Stationnarité : $\mu_1=2x$, $\mu_2=2y$. Cas $\mu_1=0$ : $x=0$, violant $x\ge1$ ; cas $\mu_2=0$ : $y=0$, violant $y\ge2$ ; reste $\mu_1,\mu_2>0$ : $x=1$, $y=2$, $\mu=(2,4)\ge0$. Le seul point KKT est $(1,2)$, minimum global (convexité), $f=5$ : le point admissible le plus proche de l'origine.
	\end{exemple}
	\begin{exemple}[Triangle]
		$\min(x_1-1)^2+(x_2-2)^2$ sous $-x_1\le0$, $-x_2\le0$, $x_1+x_2-2\le0$. Les trois contraintes ne sont jamais actives ensemble ; deux contraintes actives ont des gradients indépendants (par exemple $(-1,0)$ et $(1,1)$) : la LICQ est vérifiée en tout point admissible. Cas $\mu_3>0$, $\mu_1=\mu_2=0$ : $x_1+x_2=2$ et $2(x_1-1)=2(x_2-2)=-\mu_3$ donnent $x^*=(\frac12,\frac32)$, $\mu^*=(0,0,1)$, $f=\frac12$. Le problème étant convexe, c'est l'unique minimum global ($(1,2)$, minimum sans contrainte, est inadmissible).
	\end{exemple}
	\begin{exemple}[Disque et demi-plan]\label{ex:disque}
		$\min(x-3)^2+(y-5)^2$ sous $g_1=x^2+y^2-20\le0$, $g_2=x+y-4\le0$. Convexe ; Slater en $(0,0)$. KKT : $2(x-3)+2\mu_1x+\mu_2=0$, $2(y-5)+2\mu_1y+\mu_2=0$, complémentarité, $\mu\ge0$.
		\begin{itemize}
			\item $\mu=0$ : $(3,5)$, inadmissible.
			\item $\mu_1=0<\mu_2$ : $x+y=4$ et $x-3=y-5$ donnent $(1,3)$, $\mu_2=4>0$, $g_1(1,3)=-10<0$ : \textbf{point KKT}, $f=8$.
			\item $\mu_1>0$, $\mu_2=0$ : les deux équations donnent $\mu_1=-\frac{x-3}{x}=-\frac{y-5}{y}$, d'où $y(x-3)=x(y-5)$, soit $y=\frac53x$. Avec $x^2+y^2=20$ : $\frac{34}9x^2=20$, $x^2=\frac{90}{17}$ ($\abs x\approx2{,}30$). Pour $x\approx+2{,}30$, $y\approx3{,}83$ et $x+y\approx6{,}1>4$ : viole $g_2$. Pour $x\approx-2{,}30$ : $\mu_1=-\frac{x-3}x\approx-2{,}3<0$ : rejeté.
			\item $\mu_1,\mu_2>0$ : $x^2+y^2=20$, $x+y=4$ donnent $(2\pm\sqrt6,\,2\mp\sqrt6)$ ; en soustrayant les deux équations de stationnarité, $\mu_1=-\frac{x-y+2}{x-y}$, qui vaut $-\frac{1+\sqrt6}{\sqrt6}<0$ en $(2+\sqrt6,2-\sqrt6)$ et $\approx-0{,}59<0$ en l'autre point : rejetés.
		\end{itemize}
		Le seul point KKT est $(1,3)$, donc le minimum global (théorème~\ref{thm:kktconv}), de valeur $8$.
		\begin{center}
			\begin{tikzpicture}[scale=0.5]
				\clip (-5.2,-4.9) rectangle (7.5,8.2);
				\begin{scope}\clip (0,0) circle ({sqrt(20)});\fill[blue!15] (-6,10)--(10,-6)--(-6,-6)--cycle;\end{scope}
				\draw[->] (-5,0)--(7.2,0) node[right]{$x$};\draw[->] (0,-4.8)--(0,8) node[above]{$y$};
				\draw[thick,blue] (0,0) circle ({sqrt(20)});
				\draw[thick,orange] (-2,6)--(6,-2) node[right]{\small$x+y=4$};
				\draw[red!60,dashed] (3,5) circle ({sqrt(8)});
				\fill (3,5) circle (2pt) node[above right]{\small$(3,5)$};
				\fill[red] (1,3) circle (3pt) node[left=3pt]{\small$(1,3)$};
			\end{tikzpicture}
		\end{center}
	\end{exemple}
	\begin{exemple}[Problème non convexe : $x_1^2+x_2^2\le4e^{x_1+x_2-2}$ sur $\R_+^2$]
		Soit $f(x)=-(x_1^2+x_2^2)e^{-x_1-x_2}$, à minimiser sur $x_1,x_2\ge0$ ($g_i=-x_i\le0$ ; gradients indépendants : LICQ). Avec $S=x_1^2+x_2^2$ et $t=x_1+x_2$, $\partial_if=(S-2x_i)e^{-t}$. KKT : $(S-2x_i)e^{-t}=\mu_i\ge0$, $\mu_ix_i=0$.
		Aucune contrainte active : $S=2x_1=2x_2$ donne $(1,1)$, $f=-2e^{-2}$. Une active ($x_1=0<x_2$) : $x_2^2-2x_2=0$ donc $x_2=2$, $\mu_1=4e^{-2}$ : $(0,2)$ et $(2,0)$, $f=-4e^{-2}$. Deux actives : $(0,0)$, $f=0$.
		$f$ n'est pas coercive mais $-\norm x^2e^{-\norm x}\le f\le0$ (car $\norm x\le x_1+x_2$) donc $f\to0$ à l'infini ; comme $\min f<0$, le minimum est atteint sur un compact et c'est un point KKT : $\min f=-4e^{-2}$, atteint en $(0,2)$ et $(2,0)$. D'où $(x_1^2+x_2^2)e^{-x_1-x_2}\le4e^{-2}$ pour tout $x\ge0$, avec égalité seulement en ces points.
	\end{exemple}
	\begin{exemple}[Égalité et inégalité]
		Maximiser $y$ sous $x^2+y^2=1$, $x+y\ge0$, soit $\min -y$ avec $h=x^2+y^2-1=0$, $g=-x-y\le0$. LICQ : en un point où les deux contraintes sont actives, $\grad h=(2x,2y)$ et $\grad g=(-1,-1)$ sont indépendants ; si seule $h$ est active, $\grad h\ne0$. KKT : $2\lambda x-\mu=0$, $-1+2\lambda y-\mu=0$, $\mu(-x-y)=0$, $\mu\ge0$.
		$\mu=0$ : $x=0$, $y=\pm1$ avec $y\ge0$ : $(0,1)$ ($\lambda=\frac12$). $\mu>0$ : $x+y=0$ et $x^2+y^2=1$ donnent $\pm(\frac{\sqrt2}2,-\frac{\sqrt2}2)$ et le calcul donne $\mu=-\frac12<0$ dans les deux cas : rejetés.
		L'ensemble admissible (arc de cercle) est compact : le minimum existe et est un point KKT, donc c'est $(0,1)$ ; $\max y=1$. (Attention : $h=0$ n'est pas une contrainte convexe, le théorème~\ref{thm:kktconv} ne s'applique pas ; c'est la compacité qui conclut.)
	\end{exemple}
	\begin{retenir}
		KKT = stationnarité + admissibilité primale + $\mu\ge0$ + complémentarité. Nécessaire sous qualification (LICQ $\Rightarrow$ QCMF), suffisant si le problème est convexe. Méthode : discuter les cas selon les contraintes actives, rejeter les multiplicateurs de mauvais signe, conclure par convexité ou compacité.
	\end{retenir}
	
%	\chapter{Programmation linéaire : modélisation et résolution graphique}
	
	\chapter{Généralités sur la programmation linéaire}
	\label{chap:pl_generalites}
	
	\section{Introduction}
	
	En mathématiques appliquées et en recherche opérationnelle, les problèmes de \textbf{programmation linéaire (PL)} sont des problèmes d'optimisation dans lesquels la \textbf{fonction objectif} et les \textbf{contraintes} sont toutes deux linéaires.  
	Autrement dit, la programmation linéaire vise à \textit{trouver la meilleure décision possible} (au sens de maximiser un profit ou minimiser un coût) parmi un ensemble de solutions réalisables, défini par des contraintes linéaires.  
	\section{Les conditions de formulation}
	
	La programmation linéaire, comme étant un modèle, admet certaines hypothèses (ou conditions) que le décideur doit vérifier avant de pouvoir l'utiliser pour représenter correctement son problème réel. Ces hypothèses garantissent que la modélisation est cohérente et que les outils mathématiques de la PL s'appliquent.  
	
	Ces hypothèses sont :
	\begin{itemize}
		\item[$\diamond $] les variables de décision du problème sont positives ;
		\item[$\diamond $] le critère de sélection de la meilleure décision est décrit par une fonction linéaire de ces variables;
		\item[$\diamond $] les restrictions relatives aux variables de décision (exemple : limitation des ressources) peuvent être exprimées par un ensemble d'équations linéaires. Ces équations forment l'ensemble des contraintes;
		\item[$\diamond $]  les paramètres du problème en dehors des variables de décisions ont une valeur connue avec certitude.
	\end{itemize}
	\section{Les étapes de formulation de PL}
	Généralement, il y a trois étapes pour pouvoir construire le modèle d'un programme linéaire :
	\begin{enumerate}
		\item identifier les variables du problème à valeur non connue (variables de décision) et leur nature ou type (discret, continu, binaire)
		\item identifier les restrictions (les contraintes) du problème et les exprimer par un système d'équation linéaire ou un système d'inéquation linéaire ;
		\item  identifier l'objectif ou le critère de sélection qui se traduit par la fonction économique ou la fonction "objectif". Spécifier si la fonction objectif est à maximiser ou à minimiser.
	\end{enumerate}
	
	\section{Formes générales d'un programme linéaire}
	
	La forme la plus générale d'un problème de PL est la suivante :
	\begin{equation*}\label{key}
		\begin{array}{lc}
			opt\:\: z =\displaystyle\sum_{j=1}^{n}c_{j}x_{j}\\
			\\
			sc \left\lbrace \begin{array}{llc}
				\displaystyle\sum_{j=1}^{n}a_{ij}x_{j}\leq (\geq) b_{i}&i=\overline{1,p}& (2a)\\
				\displaystyle\sum_{j=1}^{n}a_{ij}x_{j}= b_{i}&i=\overline{1,q}& (2b)\\
				x_{j}\geq 0&j=\overline{1,q}& (3)\\
				x_{j}\in \mathbb{R}& j\in q+1,\ldots,n& (4).
			\end{array}\right.
		\end{array}
	\end{equation*}
	Dans cette formulation générale :
	\begin{itemize}
		\item[$\diamond $] $ opt $ pour optimisation peut-être une maximisation ou une minimisation ;
		\item[$\diamond $] les $ x_{j} $ représentent les variables de décision au nombre de $ n $ ;
		\item[$\diamond $] $ z $ est la fonction objectif ou économique ;
		\item[$\diamond $] $ a_{ij} $, $ c_{j} $ et $ b_{i} $ sont des constantes, c'est-à-dire des données du problème.
		\item[$\diamond $] les équations $ (2a) $ et $ (2b) $ sont les contraintes réelles du problème. Avec $ (2a) $ les contraintes d'inégalités et $ (2b) $ les contraintes d'égalités,
		\item[$\diamond $] $ (3) $ sont des contraintes de non-négativité ou de signes,
		\item[$\diamond $] $ (4) $ implique que les variables sont sans restriction de signes.
	\end{itemize}
	
	\subsection{Forme canonique mixte}
	
	Il s'agit d'un problème de programmation linéaire (encore appelé programme linéaire) écrit sous la forme suivante.
	
	\[
	\displaystyle\max_{(x_1, \ldots, x_n)} f(x_1, \ldots, x_n) = c_1 x_1 + c_2 x_2 + \cdots + c_n x_n = \sum_{j=1}^n c_j x_j
	\]
	sous contraintes :
	\[
	\begin{cases}
		\text{contraintes d'inégalités : } & \forall i \in I_1, \quad \sum_{j=1}^n a_{ij} x_j = a_{i1} x_1 + a_{i2} x_2 + \cdots + a_{in} x_n \leq b_i \\
		\text{contraintes d'égalités : } & \forall i \in I_2, \quad \sum_{j=1}^n a_{ij} x_j = b_i \\
		\text{contraintes de signes : } & \forall j \in J_1, \quad x_j \geq 0 \\
		& \forall j \in J_2, \quad x_j \text{ de signe quelconque}
	\end{cases}
	\]
	
	L'ensemble $I=I_{1} \cup I_{2}$ est l'ensemble des indices de contraintes avec $\operatorname{card}(I)=m$. Autrement dit, il y a $m$ contraintes. L'ensemble $J=J_{1} \cup J_{2}$ est l'ensemble des indices des variables avec $\operatorname{card}(J)=n$. Il y a $n$ variables.
	
	\subsection{Forme canonique pure}
	
	Sous cette forme, il n'y a pas de contraintes d'égalité c'est-à-dire $I_{2}=\varnothing$ et $J_{2}=\varnothing$.\\
	On note
	
	$$
	\begin{aligned}
		x & =(x_{1}, \ldots, x_{n})^{t} \in \mathbb{R}^{n} \\
		c & =(c_{1}, \ldots, c_{n})^{t} \in \mathbb{R}^{n} \\
		b & =(b_{1}, \ldots, b_{m})^{t} \in \mathbb{R}^{m}
	\end{aligned}
	$$
	
	et la matrice $A$ de taille $m \times n$ :
	
	$$
	A=\left(\begin{array}{cccc}
		a_{11} & a_{12} & \cdots & a_{1 n} \\
		\vdots & & & \vdots \\
		a_{m 1} & a_{m 2} & \cdots & a_{m n}
	\end{array}\right)
	$$
	
	Un programme linéaire (PL) est dit sous forme canonique pure s'il s'écrit :
	
	\begin{align*}
		& \displaystyle\max_{x}\left[f\left(x_{1}, \ldots, x_{n}\right)=c^{t} x=c_{1} x_{1}+c_{2} x_{2}+c_{3} x_{3}+\ldots +c_{n} x_{n}\right] \\
		& \left\{\begin{aligned}
			\text { s.t. } & A x \leq b \\
			& x \geq 0
		\end{aligned}\right. \tag{2.2}
	\end{align*}
	
	\begin{remarque}
		Il est toujours possible de ramener :
		
		\begin{itemize}
			\item[$\diamond $] L'optimisation par maximisation équivaut à une minimisation, et vice versa. En effet, maximiser une fonction \( z \) revient à minimiser \(-z\), et minimiser \( z \) revient à maximiser \(-z\) :
			\begin{equation}
				\displaystyle\max(z) = -\displaystyle\min(-z) \quad \text{et} \quad \displaystyle\min(z) = -\displaystyle\max(-z).
			\end{equation}
			\item[$\diamond $] toutes les inégalités à des inégalités de même type en multipliant les deux membres par \(-1\);
			\item[$\diamond $] toutes les variables à être non-négatives.
			si $ x_{j}\in\mathbb{R} $, on remplace $ x_{j} $ par $ x_{j}^{+}-x_{j}^{-} $ avec $ x_{j}^{+}=\displaystyle\max(x_{j},0)\geq0 $ et $ x_{j}^{-}=\displaystyle\max(-x_{j},0)\leq0 $,
			\item[$\diamond $] les contraintes d'égalités en contraintes d'inégalités sous la forme
			\begin{equation}\label{key}
				\sum_{j=1}^{n}a_{ij}x_{j}=b_{i}\Leftrightarrow\left\lbrace \begin{array}{cc}
					\displaystyle\sum_{j=1}^{n}a_{ij}x_{j}\geq b_{i}&\\
					&i=\overline{1,p}\\
					\displaystyle\sum_{j=1}^{n}a_{ij}x_{j}\leq b_{i}&
				\end{array}\right. 
			\end{equation}
		\end{itemize}
	\end{remarque}
	
	\subsection{Forme standard}
	
	Sous cette forme, $I_{1}=\varnothing$ et $J_{2}=\varnothing$. Un programme linéaire (PL) est dit sous forme standard s'il s'écrit :
	
	\begin{align*}
		& \displaystyle\max_{x}\Big[f\left(x_{1}, \ldots, x_{n}\right)=c^{t} x=c_{1} x_{1}+c_{2} x_{2}+c_{3} x_{3}+\ldots +c_{n} x_{n}\Big] \\
		& \left\{\begin{aligned}
			\text { s.t. } & A x=b \\
			& x \geq 0
		\end{aligned}\right. \tag{2.3}
	\end{align*}
	
	
	On dit de plus que le PL est sous forme standard simpliciale si $A$ de taille $m \times n$ avec $m \leq n$, se décompose en :
	
	$$
	A=\left(I_{m} \mid H\right)
	$$
	
	où $I_{m}$ désigne la matrice identité de taille $m \times m$ et $H$ est une matrice de taille $m \times(n-m)$.
	\section{Exemple de formulation}
	%	\subsection{Exemple 1}
	\begin{exemple}
		Un agriculteur veut allouer $ 150 $ hectares de surface irrigable entre culture de tomates et celles de piments. Il dispose de $ 480 $ heures de main d’œuvre et de $ 440 \:m^{3}$ d’eau. Un hectare de tomates demande 1 heure de main d’œuvre, $ 4 \:m^{3}$ d’eau et donne un bénéfice net de $ 1000000 $ $ FCFA $. Un hectare de piments demande $ 4 $ heures de main d’œuvre, $ 2 \:m^{3}$ d’eau et donne un bénéfice net de $ 2000000 $ $ FCFA $.
		
		Le bureau du périmètre irrigué veut protéger le prix des tomates et ne lui permet pas de cultiver plus de $ 90 $ hectares de tomates. Quelle est la meilleure allocation de ses ressources ?
		\textbf{Formulation}\\
		
		\textbf{Étape 1} Identification des variables de décision.\\
		\begin{itemize}
			\item[$\diamond $] $ x_{1}:$ la surface allouée à la culture des tomates
			\item[$\diamond $] $ x_{2}: $ la surface allouée à la culture des piments.
		\end{itemize}
		On vérifie bien que $ x_{1}\geq0 $ et $ x_{2}\geq0 $.\\
		
		\textbf{Étape 2} Identification des contraintes\\
		\begin{itemize}
			\item[$ \diamond $] Terrain : $ x_{1}+x_{2}\leq150 $
			\item[$\diamond $] Eau : $ 4x_{1}+2x_{2}\leq440 $
			\item[$\diamond $] Main d'œuvre : $ x_{1}+4x_{2}\leq480 $
			\item[$\diamond $] La limitation du bureau du périmètre irrigué $ x_{1}\leq 90 $.
		\end{itemize}
		\textbf{Étape 3} Identification de la fonction objectif\\
		$ z=10^{6}x_{1}+2.10^{6}x_{2} $.\\
		
		Le programme linéaire qui modélise le problème d'agriculture est :
		\begin{equation}\label{key}
			\begin{array}{lc}
				\displaystyle\max\:\:10^{6}x_{1}+2.10^{6}x_{2}  \\
				sc \left\lbrace \begin{array}{llc}
					x_{1}+x_{2}\leq150 \\
					4x_{1}+2x_{2}\leq440\\
					x_{1}+4x_{2}\leq480 \\
					x_{1}\leq 90 \\
					x_{1}\geq0,\: x_{2}\geq0.
				\end{array}\right.
			\end{array}
		\end{equation}
	\end{exemple}
	
	
%		\chapter{Programmation linéaire : modélisation et résolution graphique}
%	
%	\section{Hypothèses et étapes de formulation}
%	La programmation linéaire (PL) suppose : variables de décision positives ; critère de choix \emph{linéaire} en ces variables ; restrictions exprimables par des équations ou inéquations linéaires ; paramètres connus avec certitude. Formuler un PL, c'est (1) identifier les variables de décision et leur nature (continue, discrète, binaire) ; (2) écrire les contraintes ; (3) écrire la fonction objectif (économique) et préciser max ou min.
%	
%	\section{Formes d'un programme linéaire}
%	\textbf{Forme générale.} $\operatorname{opt}z=\sum_jc_jx_j$ sous des contraintes $\sum_ja_{ij}x_j\le(\ge)\,b_i$ ou $=b_i$, avec $x_j\ge0$ pour certaines variables et $x_j$ de signe quelconque pour les autres.
%	
%	\textbf{Forme canonique mixte.} $\max\sum_jc_jx_j$ avec : pour $i\in I_1$, $\sum_ja_{ij}x_j\le b_i$ ; pour $i\in I_2$, $\sum_ja_{ij}x_j=b_i$ ; $x_j\ge0$ pour $j\in J_1$ et $x_j$ de signe quelconque pour $j\in J_2$. On note $I=I_1\cup I_2$ ($m$ contraintes) et $J=J_1\cup J_2$ ($n$ variables).
%	
%	\textbf{Forme canonique pure.}\textbf{Forme canonique pure.} $\max\,c\T x$ sous $Ax\le b$, $x\ge0$, avec $A\in\R^{m\times n}$.\\
%	\textbf{Forme standard.} $\max\,c\T x$ sous $Ax=b$, $x\ge0$.\\
%	Elle est \emph{simpliciale} si $A=(I_m\mid H)$.
%	
%	\textbf{Passages entre formes.}
%	\begin{itemize}
%		\item $\max z=-\min(-z)$ ; une inégalité $\ge$ devient $\le$ en multipliant par $-1$.
%		\item Variable de signe quelconque : $x_j=x_j^+-x_j^-$ avec $x_j^+=\max(x_j,0)\ge0$ et $x_j^-=\max(-x_j,0)\ge0$.
%		\item Égalité : $a\T x=b\iff a\T x\le b$ et $-a\T x\le-b$.
%		\item Inégalité $\le$ $\to$ égalité : on ajoute une \textbf{variable d'écart} $e_i\ge0$ ; pour $\ge$ on soustrait une \textbf{variable d'excès} $e_i\ge0$ ; pour $\ge$ ou $=$ on ajoute en outre une \textbf{variable artificielle} $t_i\ge0$, outil temporaire pour démarrer le simplexe (phase I, \S\ref{sec:deux}).
%	\end{itemize}
%	\begin{proposition}
%		Tout PL sous forme canonique pure est équivalent à un PL sous forme standard (variables d'écart : $(A\mid I_m)\binom xe=b$) et réciproquement ($Ax=b\iff\binom{A}{-A}x\le\binom b{-b}$).
%	\end{proposition}
%	\begin{exemple}
%		$\min z=x_1-2x_2$ sous $-x_1+4x_2\ge5$, $2x_1+3x_2=4$, $x_1\ge0$, $x_2$ libre. Avec $x_2=x_2^+-x_2^-$ et $\max(-z)=-x_1+2x_2^+-2x_2^-$ :
%		\[x_1-4x_2^++4x_2^-\le-5,\quad 2x_1+3x_2^+-3x_2^-\le4,\quad-2x_1-3x_2^++3x_2^-\le-4,\quad x_1,x_2^\pm\ge0.\]
%	\end{exemple}
%	
%	\section{Exemples de modélisation}
%	\begin{exemple}[Production]\label{ex:prod}
%		Une usine fabrique $P_1,P_2$ avec équipement (3 ; 9 ; disponible 81), main d'œuvre (4 ; 5 ; 55), matière première (2 ; 1 ; 20) ; bénéfices unitaires $6$ et $4$. Avec $x_j$ la quantité de $P_j$ :
%		\[\max f=6x_1+4x_2\quad\text{s.c.}\quad3x_1+9x_2\le81,\ 4x_1+5x_2\le55,\ 2x_1+x_2\le20,\ x_1,x_2\ge0.\]
%	\end{exemple}
%	\begin{exemple}[Mélange minimal]
%		Une pilule petite contient (aspirine, bicarbonate, codéine) $=(2,5,1)$ grains, une grande $(1,8,6)$ ; un patient doit recevoir au moins $12$, $74$, $24$ grains. Avec $x_1,x_2$ le nombre de petites et grandes pilules :
%		$\min x_1+x_2$ s.c. $2x_1+x_2\ge12$, $5x_1+8x_2\ge74$, $x_1+6x_2\ge24$, $x_1,x_2\ge0$.
%	\end{exemple}
%	\begin{exemple}[Allocation de ressources]
%		$150$ ha, $480$ h de main d'œuvre, $440\,\mathrm{m}^3$ d'eau ; tomates : $1$ h, $4\,\mathrm{m}^3$, bénéfice $10^6$ FCFA/ha ; piments : $4$ h, $2\,\mathrm{m}^3$, $2\cdot10^6$ FCFA/ha ; au plus $90$ ha de tomates :
%		$\max10^6x_1+2\cdot10^6x_2$ s.c. $x_1+x_2\le150$, $4x_1+2x_2\le440$, $x_1+4x_2\le480$, $x_1\le90$, $x_1,x_2\ge0$.
%	\end{exemple}
%	
%	\section{Résolution graphique (deux variables)}
%	Chaque contrainte $\le$ définit un demi-plan ; leur intersection est le domaine admissible $D$ (polygone convexe). La fonction objectif définit une famille de droites parallèles $c_1x_1+c_2x_2=z$ (isomarges, ou isocoûts), de vecteur normal $c=(c_1,c_2)$. Étapes : (1) tracer $D$ ; (2) tracer une droite de niveau et le vecteur $c$ ; (3) la déplacer dans le sens de $c$ (max) ou opposé (min) jusqu'à la dernière position touchant $D$ : le contact est un sommet (ou une arête).
%	\begin{exemple}[Production, suite]
%		Les sommets de $D$ sont $(0,0)$, $(10,0)$, $(\frac{15}2,5)$, $(\frac{30}7,\frac{53}7)$, $(0,9)$. Le maximum est atteint en $(\frac{15}2,5)$ (intersection de $4x_1+5x_2=55$ et $2x_1+x_2=20$) : $f=6\cdot\frac{15}2+4\cdot5=65$.
%		\begin{center}
%			\begin{tikzpicture}[scale=0.5]
%				\fill[blue!15] (0,0)--(10,0)--(7.5,5)--(4.2857,7.5714)--(0,9)--cycle;
%				\draw[->] (0,0)--(12,0) node[right]{$x_1$};\draw[->] (0,0)--(0,11) node[above]{$x_2$};
%				\draw[thick] (0,0)--(10,0)--(7.5,5)--(4.2857,7.5714)--(0,9);
%				\draw[red,dashed] (4.5,9.5)--(10.8,0.2) node[right]{\small$6x_1+4x_2=65$};
%				\draw[->,thick,red] (7.5,5)--(8.7,5.8) node[right]{\small$c$};
%				\fill (7.5,5) circle (3pt) node[below left]{\small$(\frac{15}{2},5)$};
%				\foreach \x in {5,10}\draw (\x,.1)--(\x,-.1) node[below]{\small\x};
%				\foreach \y in {5,9}\draw (.1,\y)--(-.1,\y) node[left]{\small\y};
%			\end{tikzpicture}
%		\end{center}
%	\end{exemple}
%	\begin{exemple}[Allocation agricole, résolue]
%		On maximise $x_1+2x_2$ (en millions de FCFA) sous $x_1+x_2\le150$, $4x_1+2x_2\le440$, $x_1+4x_2\le480$, $x_1\le90$, $x\ge0$. Sommets du domaine : $(0,0)$, $(90,0)$, $(90,40)$ [$x_1=90$ et $4x_1+2x_2=440$], $(70,80)$ [$x_1+x_2=150$ et $4x_1+2x_2=440$], $(40,110)$ [$x_1+x_2=150$ et $x_1+4x_2=480$], $(0,120)$. Valeurs de $x_1+2x_2$ : $0,\,90,\,170,\,230,\,260,\,240$. Le maximum est en $(40,110)$ : $40$ ha de tomates, $110$ ha de piments, bénéfice $260$ millions de FCFA (il reste $4\cdot40+2\cdot110=380<440\,\mathrm{m}^3$ d'eau et $40<90$ ha de tomates : ces deux contraintes ne sont pas saturées).
%	\end{exemple}
%	\begin{exemple}[Minimisation]
%		$\min6x_1+9x_2$ s.c. $3x_1+2x_2\ge18$, $x_1+3x_2\ge12$, $x_1+x_2\ge8$, $x\ge0$. Sommets du domaine (non borné) : $(0,9)$, $(2,6)$, $(6,2)$, $(12,0)$, de coût $81,66,54,72$. On déplace l'isocoût (pente $-\frac23$) vers l'origine : minimum $54$ en $(6,2)$.
%	\end{exemple}
%	\begin{center}
%		\begin{tikzpicture}[scale=0.42]
%			\fill[blue!15] (0,9)--(2,6)--(6,2)--(12,0)--(14,0)--(14,11)--(0,11)--cycle;
%			\draw[->] (0,0)--(14.5,0) node[right]{$x_1$};\draw[->] (0,0)--(0,11.5) node[above]{$x_2$};
%			\draw[thick] (0,9)--(2,6)--(6,2)--(12,0)--(14,0);\draw[thick](0,11)--(0,9);
%			\draw[red,dashed] (0,6)--(9,0) node[right]{\small$6x_1+9x_2=54$};
%			\fill[red] (6,2) circle (3.5pt) node[above right]{\small$(6,2)$};
%			\foreach \p in {(0,9),(2,6),(12,0)}\fill \p circle (2.5pt);
%		\end{tikzpicture}
%	\end{center}
%	\begin{exemple}[Solutions multiples]\label{ex:multiple}
%		$\min-50x_1-30x_2$ s.c. $10x_1+6x_2\le45$, $4x_1+6x_2\le36$, $2x_1+6x_2\le27$, $x\ge0$. L'isocoût est parallèle à la première contrainte ($50:30=10:6$) : le minimum $-225$ est atteint en tout point du segment $[(\frac92,0),(\frac94,\frac{15}4)]$ : infinité de solutions optimales.
%	\end{exemple}
%	
%	\chapter{Résolution graphique}
%	
%	Dans le cas d'un (PL) à deux variables, on peut envisager une résolution graphique. Les contraintes où apparaissent des inégalités correspondent géométriquement à des demi-plans. L'intersection de ces demi-plans forme l'ensemble des variables satisfaisant à toutes les contraintes.
%	
%	\section{Technique de résolution graphique dans $\mathbb{R}^{2}$}
%	
%	La résolution graphique d'un programme linéaire ne peut se faire que pour les dimensions 2 et 3. Elle est donc peu efficace, mais permet de bien comprendre le rôle de PL.
%	Dans $\mathbb{R}^{2}$, la  représentation graphique des éléments du problème s'effectue dans le cadre d'un repère où les valeurs de $x_{1}$ figurent sur l'axe des abscisses et celle de $x_{2}$ sur l'axe des ordonnées.\\
%	Les étapes de résolution sont les suivantes :
%	\begin{itemize}
%		\item[$ \diamond $]\textbf{ Étape 1} \\
%		Tracer le domaine des solutions admissibles, cette étape consiste à représenter graphiquement les différentes contraintes du système dans le repère. Et cette représentation permet de localiser l'ensemble $ D$.
%		\item[$ \diamond $] \textbf{Étape 2} \\
%		Tracer les droites isomarge  (resp isocoût si on est dans un problème de minimisation ) ie les droites $z=c_{1}x_{1}+c_{2}x_{2}$.
%		Dans le cas d'une minimisation pour diminuer la valeur de $z=c_{1}x_{1}+c_{2}x_{2}$, il suffit de déplacer la droite d'isocoût en suivant le sens inverse d'un vecteur $\overrightarrow{u}\begin{pmatrix}
%			c_{1}\\
%			c_{2}
%		\end{pmatrix}$ et pour augmenter la valeur de $z$, il faut suivre le sens du vecteur $\overrightarrow{u}$. Il est donc nécessaire de représenter le sens du vecteur $\overrightarrow{u}$.
%		\item[$ \diamond $] \textbf{Étape 3} \\
%		L'optimum est alors la solution réalisable située sur la droite isomarge la plus haute pour le cas de maximisation et respectivement sur la droite isocoût la plus basse le cas de minimisation.
%	\end{itemize}
%	
%	Dans cette section, nous allons étudier deux types de problèmes :
%	
%	\subsection{Résolution graphique d'un problème de maximisation}
%	
%	Reprenons l'exemple de production.
%	
%	$$
%	\left\{\begin{aligned}
%		\displaystyle\max_{\left(x_{1}, x_{2}\right)}\Big[f\left(x_{1}, x_{2}\right)=6 x_{1}+4 x_{2}\Big] &  \\
%		\text { s.t. } \quad 3 x_{1}+9 x_{2} & \leq 81 \\
%		4 x_{1}+5 x_{2} & \leq 55 \\
%		2 x_{1}+x_{2} & \leq 20 \\
%		x_{1}, x_{2} & \geq 0
%	\end{aligned}\right.
%	$$
%	
%	A la fonction objectif $f$ correspond une droite $f\left(x_{1}, x_{2}\right)=6 x_{1}+4 x_{2}=$ constante, de coefficient directeur $\left(-\dfrac{6}{4}, 1\right)$. La constante précédente qui définie la droite doit être la plus grande possible (maximisation) et rencontrer l'ensemble des variables qui satisfont les contraintes. Pour déterminer cette valeur maximale, on fait donc "glisser" la droite (translation parallèle à la direction de la droite) du haut vers le bas jusqu'à rencontrer l'ensemble des variables satisfaisant les contraintes. Le maximum de $f$ sur cet ensemble des contraintes est alors atteint. On obtient ainsi la solution optimale $\left(x_{1}, x_{2}\right)=\left(\dfrac{15}{2},5\right)$ et ce qui donne une valeur maximale $\displaystyle\max (f)=65$.
%	
%	On remarque que l'ensemble des contraintes (la partie hachurée de la figure \ref{fig2.1}) est un polygone convexe et que le maximum de $f$ est atteint en un sommet de ce polygone. Cette observation est, en fait, un résultat général que l'on établira dans les sections suivantes.
%	
%	\begin{figure}[h]
%		\begin{center}
%			\includegraphics[scale=0.35]{2025_10_05_27e91bcb17e07a381837g-03}
%			%			\captionsetup{labelformat=empty}
%			\caption{Résolution graphique du problème de production}
%			\label{fig2.1}
%		\end{center}
%	\end{figure}
%	
%	\subsection{Résolution graphique d'un problème de minimisation}
%	
%	Pour résoudre le problème linéaire dans le cas d'une minimisation, nous allons appliquer une démarche équivalente, mais dans ce cas, nous cherchons la droite (représentant la fonction économique) la plus rapprochée de l'origine.
%	
%	Résoudre le problème linéaire suivant :
%	
%	$$
%	\left\{\begin{aligned}
%		\displaystyle\min \quad & 6 x_{1}+9 x_{2}  \\
%		\text { s.t } \quad & 3 x_{1}+2 x_{2}  \geq 18 \\
%		& x_{1}+3 x_{2}  \geq 12 \\
%		& x_{1}+x_{2}  \geq 8 \\
%		& x  \geq 0
%	\end{aligned}\right.
%	$$
%	
%	Représentons graphiquement l'ensemble-solution puis représentons la famille de droites correspondante à :
%	
%	$$
%	\begin{aligned}
%		6 x_{1}+9 x_{2}=p & \Longleftrightarrow x_{2}=-\dfrac{2}{3} x_{1}+\dfrac{p}{9} \\
%		& \Longleftrightarrow \text { droite de pente } -\dfrac{2}{3}
%	\end{aligned}
%	$$
%	
%	\begin{figure}[h]
%		\begin{center}
%			\includegraphics[scale=0.35]{2025_10_05_27e91bcb17e07a381837g-04}
%			%			\captionsetup{labelformat=empty}
%			\caption{Résolution graphique d'un problème de minimisation}
%		\end{center}
%	\end{figure}
%	
%	\textbf{Exercice 2.1} Résoudre graphiquement le problème suivant
%	
%	$$
%	\left\{\begin{aligned}
%		\displaystyle\min \quad & -50 x_{1}-30 x_{2}  \\
%		\text { s.t } \quad & 10 x_{1}+6 x_{2}  \leq 45 \\
%		& 4 x_{1}+6 x_{2}  \leq 36 \\
%		& 2 x_{1}+6 x_{2}  \leq 27 \\
%		& x  \geq 0
%	\end{aligned}\right.
%	$$
%	
%	\chapter{Solutions de base et méthode algébrique}
%	
%	\section{Solutions de base réalisables}
%	On travaille en forme standard $Ax=b$, $x\ge0$, avec $A\in\R^{m\times n}$ et l'\textbf{hypothèse $\rg A=m\le n$}. Elle n'est pas restrictive : si $\rg A<m$, soit $b\notin\operatorname{Im}A$ et le système n'a pas de solution, soit des équations sont redondantes et on les supprime. Si $m=n$, la solution $A^{-1}b$ est unique ; si $m<n$, il y en a une infinité.
%	\begin{definition}
%		Une \textbf{base} est un ensemble $B\subset\{1,\dots,n\}$ de $m$ indices tel que $A_B$ (colonnes $A_j$, $j\in B$) soit inversible. Les $x_B$ sont les variables de base, les $x_H$ ($j\notin B$) les variables hors base. La \textbf{solution de base} associée est $x_H=0$, $x_B=A_B^{-1}b$ ; elle est \textbf{réalisable} si $x_B\ge0$. Il y a au plus $\binom nm$ solutions de base.
%	\end{definition}
%	\begin{exemple}\label{ex:base}
%		Le PL de l'exemple~\ref{ex:prod} en forme standard (variables d'écart $e_1,e_2,e_3$) a $m=3$, $n=5$. La base $B=\{e_1,e_2,e_3\}$ ($A_B=I_3$) donne $(x_1,x_2,e_1,e_2,e_3)=(0,0,81,55,20)$ : le sommet $(0,0)$, $f=0$.
%	\end{exemple}
%	
%	\begin{exemple}
%		Base $B=\{x_1,e_1,e_2\}$ : $A_B=\begin{pmatrix}3&1&0\\4&0&1\\2&0&0\end{pmatrix}$ (colonnes de $x_1,e_1,e_2$), $x_B=A_B^{-1}b=(10,51,15)$, avec $x_2=e_3=0$ : le sommet $(10,0)$, $f=60$. Les bases à $m=3$ colonnes parmi $n=5$ fournissent au plus $\binom53=10$ solutions de base, dont certaines non réalisables.
%	\end{exemple}
%	
%	\section{Géométrie}
%	$\mathbf D_R=\{x\mid Ax=b,x\ge0\}$ est un \textbf{polyèdre} convexe fermé ; un polyèdre borné est un \textbf{polytope}. Un point $x\in\mathbf D_R$ est un \textbf{sommet} (point extrême) s'il n'existe pas $y\ne z$ dans $\mathbf D_R$ et $\lambda\in]0,1[$ tels que $x=\lambda y+(1-\lambda)z$.
%	\begin{theorem}\label{thm:sommet}
%		$x$ est une solution de base réalisable $\iff$ $x$ est un sommet de $\mathbf D_R$.
%	\end{theorem}
%	\begin{proof}[Preuve \emph{(complément)}]
%		\textbf{($\Rightarrow$) Une solution de base réalisable est un sommet.} Soit $x$ de base, $x_H=0$, et supposons $x=\lambda y+(1-\lambda)z$ avec $y,z\in\mathbf D_R$, $\lambda\in]0,1[$. Pour $j\in H$ : $0=\lambda y_j+(1-\lambda)z_j$ avec $y_j,z_j\ge0$ et $\lambda,1-\lambda>0$, donc $y_j=z_j=0$. Alors $A_By_B=b=A_Bz_B$ ; $A_B$ étant inversible, $y_B=z_B=A_B^{-1}b=x_B$. Donc $y=z=x$ : $x$ est extrême.
%		\textbf{($\Leftarrow$) Un sommet est une solution de base.} Soit $x$ un sommet et $S=\{j\mid x_j>0\}$. Montrons que les colonnes $A_j$, $j\in S$, sont indépendantes. Si elles étaient liées, il existerait $w\ne0$, nul hors de $S$, avec $Aw=0$. Pour $\varepsilon>0$ assez petit, $y=x+\varepsilon w$ et $z=x-\varepsilon w$ restent $\ge0$ (car $x_j>0$ sur $S$ et $w_j=0$ ailleurs) et $Ay=Az=b$ : $y,z\in\mathbf D_R$, $y\ne z$ et $x=\frac12(y+z)$, ce qui contredit le caractère extrême. Les colonnes de $S$ étant indépendantes (donc $\abs S\le m$), on les complète en $m$ colonnes indépendantes (possible car $\rg A=m$) : cela définit une base $B\supseteq S$, et $x_H=0$.
%	\end{proof}
%	\begin{center}
%		\begin{tikzpicture}[x={(-0.62cm,-0.45cm)}, y={(1cm,0cm)}, z={(0cm,1cm)}, scale=1.5]
%			\coordinate (O) at (0,0,0);\coordinate (A) at (1.5,0,0);\coordinate (B) at (0,1,0);\coordinate (C) at (0,0,3);
%			\draw[->] (O)--(2.2,0,0) node[below left]{$x_1$};\draw[->] (O)--(0,1.6,0) node[right]{$x_2$};\draw[->] (O)--(0,0,3.6) node[above]{$x_3$};
%			\fill[pattern=north east lines,pattern color=gray!60] (A)--(B)--(C)--cycle;\draw[thick] (A)--(B)--(C)--cycle;
%			\node[below left] at (A){\small$(\frac32,0,0)$};\node[right] at (B){\small$(0,1,0)$};\node[left] at (C){\small$(0,0,3)$};
%		\end{tikzpicture}\\
%		{\small Ensemble $\mathbf D_R=\{2x_1+3x_2+x_3=3,\ x\ge0\}\subset\R^3$ : un triangle, de sommets $(\frac32,0,0)$, $(0,1,0)$, $(0,0,3)$ (solutions de base : $m=1$, $n=3$).}
%	\end{center}
%	\begin{theorem}
%		Si $\mathbf D_R\ne\emptyset$ (et ne contient pas de droite), $\mathbf D_R$ a au moins un sommet. Si le PL admet un optimum, il est atteint en au moins un sommet de $\mathbf D_R$.
%	\end{theorem}
%	Trois situations seulement : (1) $\mathbf D_R=\emptyset$ ; (2) $\mathbf D_R\ne\emptyset$ mais $f$ non majorée (optimum infini ; exclu si $\mathbf D_R$ est borné) ; (3) il existe une solution optimale (pas nécessairement unique). Il suffit donc de parcourir les solutions de base réalisables ; le simplexe le fait sans les énumérer toutes (une énumération coûterait de l'ordre de $m^3\binom nm$ opérations).
%	
%	\begin{definition}
%		Un PL (ou un domaine) est \textbf{non borné} si au moins une variable peut tendre vers l'infini sans quitter le domaine réalisable. Les ensembles $\{Ax\le b\}$ sont des polyèdres (convexes) ; un polyèdre borné est un polytope.
%	\end{definition}
%	
%	\section{Méthode algébrique sur l'exemple de production}
%	On part de la base $\{e_1,e_2,e_3\}$, $f=0$, et on exprime les variables de base en fonction des variables hors base.
%	\begin{enumerate}
%		\item \emph{Variable entrante.} $f=6x_1+4x_2$ : $x_1$ (coefficient $6$, le plus grand). \emph{Sortante.} Valeurs maximales de $x_1$ dans chaque équation : $27$, $\frac{55}4$, $10$ ; le minimum $10$ est atteint pour $e_3$, qui sort. Nouveau sommet $(10,0)$, $f=60$, base $\{x_1,e_1,e_2\}$ :
%		\[x_1=10-\tfrac12x_2-\tfrac12e_3,\quad e_1=51-\tfrac{15}2x_2+\tfrac32e_3,\quad e_2=15-3x_2+2e_3,\quad f=60+x_2-3e_3.\]
%		\item $x_2$ a un coefficient positif dans $f$ : il entre. Ratios : $\frac{51}{15/2}=\frac{34}5$, $\frac{15}3=5$, $\frac{10}{1/2}=20$ ; $e_2$ sort. Sommet $(\frac{15}2,5)$, base $\{x_1,x_2,e_1\}$ :
%		\[x_1=\tfrac{15}2+\tfrac16e_2-\tfrac56e_3,\quad x_2=5-\tfrac13e_2+\tfrac23e_3,\quad e_1=\tfrac{27}2+\tfrac52e_2-\tfrac72e_3,\quad f=65-\tfrac13e_2-\tfrac73e_3.\]
%	\end{enumerate}
%	Tous les coefficients de $f$ sur les variables hors base sont $\le0$ : on ne peut plus améliorer : \textbf{optimum} $f=65$, en accord avec la résolution graphique.
%	
%	\begin{retenir}
%		Un programme de base est optimal (max) si tous les coefficients des variables hors base dans l'expression de $f$ sont $\le0$. Marche à suivre : forme standard $\to$ base de départ $\to$ exprimer $f$ en variables hors base $\to$ entrer la variable de plus grand coefficient positif $\to$ sortir celle du plus petit ratio positif $\to$ recommencer.
%	\end{retenir}
%	\chapter{La méthode du simplexe}
%	La méthode (Dantzig, 1947) comporte deux phases : \textbf{phase I} (trouver une solution de base réalisable ou détecter que $\mathbf D_R=\emptyset$) et \textbf{phase II} (passer de sommet en sommet voisin en améliorant $f$, ou détecter que $f$ est non majorée).
%	
%	\section{Principe de la phase II}
%	On suppose une solution de base réalisable $\bar x$ de base $B$ : $A=(A_B\mid A_H)$, $\bar x_B=A_B^{-1}b\ge0$ (pour $Ax\le b$, $b\ge0$, la base des variables d'écart convient). Comme $x_B=A_B^{-1}(b-A_Hx_H)$,
%	\[f(x)=c_B\T A_B^{-1}b+\bigl(c_H\T-c_B\T A_B^{-1}A_H\bigr)x_H=f(\bar x)+L_H\T x_H,\]
%	où $L_H\T=c_H\T-c_B\T A_B^{-1}A_H$ est le vecteur des \textbf{coûts réduits}.
%	\begin{itemize}
%		\item \textbf{Optimalité (max).} Si $L_H\le0$, on ne peut pas augmenter $f$ : $\bar x$ est optimale.
%		\item \textbf{Variable entrante.} Sinon, on choisit $e=\operatorname{argmax}_j\{(L_H)_j\mid(L_H)_j>0\}$. Pour une \emph{minimisation}, le critère est inversé : optimal si $L_H\ge0$, entrante $=\operatorname{argmin}\{(L_H)_j<0\}$.
%		\item \textbf{Variable sortante.} On pose $z=A_B^{-1}A_e$ ; alors $x_B=\bar x_B-zx_e$. Si $z\le0$, $x_e$ peut croître indéfiniment : $f$ est \textbf{non majorée}, arrêt. Sinon $x_e=\min_i\{\frac{(\bar x_B)_i}{z_i}\mid z_i>0\}$, et la variable de base réalisant ce minimum sort.
%	\end{itemize}
%	On remplace $A_s$ par $A_e$ dans la base, on recalcule, et on recommence.
%	
%	\section{Représentation par tableaux}
%	Le PL en forme standard est écrit avec la ligne $c\T x-z=0$. Le tableau a une ligne par variable de base, une dernière ligne des coûts réduits $c_j$ (nuls pour les variables de base), et en bas à droite la valeur $-z$ courante. \textbf{Pivot} : élément de la colonne entrante et de la ligne du plus petit ratio $b_i/a_{ie}$ ($a_{ie}>0$). On divise la ligne pivot par le pivot, puis on annule le reste de la colonne par opérations sur les lignes (\emph{règle du rectangle} : $a'=a-\frac{b\,c}{\text{pivot}}$, où $b$ est dans la colonne du pivot et $c$ dans la ligne du pivot ; une ligne dont l'élément de la colonne pivot est nul, ou une colonne dont l'élément de la ligne pivot est nul, reste inchangée).
%	
%	\begin{retenir}
%		\textbf{Marche à suivre (maximisation, $b\ge0$, contraintes $\le$).}
%		1.~Forme standard avec variables d'écart ; base initiale = les écarts ; tableau initial.
%		2.~Si tous les coûts de la ligne $z$ sont $\le0$ : fin (optimum).
%		3.~Colonne pivot : le plus grand coût $>0$.
%		4.~Si tous les éléments de cette colonne sont $\le0$ : problème non borné, fin.
%		5.~Ligne pivot : plus petit ratio $b_i/a_{ie}$ parmi les $a_{ie}>0$.
%		6.~Diviser la ligne pivot par le pivot, annuler le reste de la colonne par opérations sur les lignes, changer le nom de la variable de base, retour en 2.
%	\end{retenir}
%	\begin{exemple}[Une itération détaillée, à partir de $T_0$]
%		(1) \emph{Lire la solution} : variables de base $e_1,e_2,e_3$ valent le second membre $(81,55,20)$, variables hors base $x_1=x_2=0$ ; la valeur courante est $z=0=-(\text{coin droit})$. (2) \emph{Choisir la colonne} : le plus grand coût de la ligne $z$ est $6$ (colonne $x_1$) : augmenter $x_1$ améliore $z$ de $6$ par unité. (3) \emph{Choisir la ligne} : on ne regarde que les éléments \emph{strictement positifs} de la colonne $x_1$ : $3,4,2$ ; ratios $81/3=27$, $55/4=13{,}75$, $20/2=10$ ; le plus petit est $10$ (ligne $e_3$) : on ne peut pas dépasser $x_1=10$ sans rendre $e_3$ négative. (4) \emph{Pivoter} : $L_3\leftarrow\frac12L_3$, puis $L_1\leftarrow L_1-3L_3$, $L_2\leftarrow L_2-4L_3$, $L_z\leftarrow L_z-6L_3$. (5) \emph{Relire} : $x_1=10$ remplace $e_3$ dans la base ; $-z=-60$ signifie $z=60$.
%	\end{exemple}
%	\begin{exemple}[Maximisation : production]\label{ex:simpl}
%		Forme standard avec $e_1,e_2,e_3$. Tableau initial $T_0$ ; $x_1$ entre (coefficient $6$) ; ratios $27,\frac{55}4,10$ : le pivot est $2$ (ligne $e_3$) ; puis $x_2$ entre (coefficient $1$) ; ratios $\frac{34}5,5,20$ : pivot $3$ (ligne $e_2$).
%		\begin{center}\small
%			\begin{tabular}{c|ccccc|c}
%				$T_0$ & $x_1$&$x_2$&$e_1$&$e_2$&$e_3$&$b$\\\hline
%				$e_1$&3&9&1&0&0&81\\ $e_2$&4&5&0&1&0&55\\ $e_3$&2&1&0&0&1&20\\\hline $z$&6&4&0&0&0&0
%			\end{tabular}\quad
%			\begin{tabular}{c|ccccc|c}
%				$T_1$ & $x_1$&$x_2$&$e_1$&$e_2$&$e_3$&$b$\\\hline
%				$e_1$&0&$\frac{15}2$&1&0&$-\frac32$&51\\ $e_2$&0&3&0&1&$-2$&15\\ $x_1$&1&$\frac12$&0&0&$\frac12$&10\\\hline $z$&0&1&0&0&$-3$&$-60$
%			\end{tabular}
%			\medskip
%			
%			\begin{tabular}{c|ccccc|c}
%				$T_2$ & $x_1$&$x_2$&$e_1$&$e_2$&$e_3$&$b$\\\hline
%				$e_1$&0&0&1&$-\frac52$&$\frac72$&$\frac{27}2$\\ $x_2$&0&1&0&$\frac13$&$-\frac23$&5\\ $x_1$&1&0&0&$-\frac16$&$\frac56$&$\frac{15}2$\\\hline $z$&0&0&0&$-\frac13$&$-\frac73$&$-65$
%			\end{tabular}
%		\end{center}
%		Tous les coûts réduits de $T_2$ sont $\le0$ : optimum $(x_1,x_2,e_1,e_2,e_3)=(\frac{15}2,5,\frac{27}2,0,0)$, $z=65$. $T_1$ et $T_2$ reproduisent exactement les systèmes de la méthode algébrique.
%	\end{exemple}
%	
%	\begin{exemple}[Minimisation]
%		$\min-x_1-2x_2$ s.c. $-3x_1+2x_2\le2$, $-x_1+2x_2\le4$, $x_1+x_2\le5$, $x\ge0$. Ligne des coûts $(-1,-2,0,0,0)$, entrante = coût le plus négatif.
%		
%		\begin{center}\scriptsize
%			\begin{tabular}{c|ccccc|c}
%				$T_0$&$x_1$&$x_2$&$e_1$&$e_2$&$e_3$&$b$\\\hline
%				$e_1$&$-3$&\textbf{2}&1&0&0&2\\$e_2$&$-1$&2&0&1&0&4\\$e_3$&1&1&0&0&1&5\\\hline $z$&$-1$&$-2$&0&0&0&0
%			\end{tabular}\quad
%			\begin{tabular}{c|ccccc|c}
%				$T_1$&$x_1$&$x_2$&$e_1$&$e_2$&$e_3$&$b$\\\hline
%				$x_2$&$-\frac32$&1&$\frac12$&0&0&1\\$e_2$&\textbf{2}&0&$-1$&1&0&2\\$e_3$&$\frac52$&0&$-\frac12$&0&1&4\\\hline $z$&$-4$&0&1&0&0&2
%			\end{tabular}\quad
%			\begin{tabular}{c|ccccc|c}
%				$T_2$&$x_1$&$x_2$&$e_1$&$e_2$&$e_3$&$b$\\\hline
%				$x_2$&0&1&$-\frac14$&$\frac34$&0&$\frac52$\\$x_1$&1&0&$-\frac12$&$\frac12$&0&1\\$e_3$&0&0&\textbf{$\frac34$}&$-\frac54$&1&$\frac32$\\\hline $z$&0&0&$-1$&2&0&6
%			\end{tabular}
%		\end{center}
%		(Le pivot est en gras ; le coin inférieur droit vaut $-z$.) \emph{Règle du rectangle, passage $T_1\to T_2$} (pivot $2$ en ligne $e_2$, colonne $x_1$) : cellule $(x_2,e_1)$, $a=\frac12$, $b=-\frac32$ (même ligne, colonne du pivot), $c=-1$ (même colonne, ligne du pivot) : $a'=\frac12-\frac{(-\frac32)(-1)}{2}=-\frac14$, valeur lue dans $T_2$. Les cellules de la colonne $e_3$ de la ligne $e_2$ valent $0$ : la colonne $e_3$ est inchangée sauf pour la ligne pivot.
%		\emph{Itération 1} : $x_2$ entre, ratios $1,2,5$ : $e_1$ sort (pivot $2$) ; $z=-2$. \emph{Itération 2} : $x_1$ entre (coût $-4$), seuls $e_2$ (ratio $1$) et $e_3$ (ratio $\frac85$) ont un coefficient positif : $e_2$ sort ; $z=-6$. \emph{Itération 3} : $e_1$ entre (coût $-1$), seule la ligne $e_3$ est positive : $e_3$ sort. Tableau final :
%		\begin{center}\small
%			\begin{tabular}{c|ccccc|c}
%				&$x_1$&$x_2$&$e_1$&$e_2$&$e_3$&$b$\\\hline
%				$x_2$&0&1&0&$\frac13$&$\frac13$&3\\ $x_1$&1&0&0&$-\frac13$&$\frac23$&2\\ $e_1$&0&0&1&$-\frac53$&$\frac43$&2\\\hline $z$&0&0&0&$\frac13$&$\frac43$&8
%			\end{tabular}
%		\end{center}
%		Coûts réduits $\ge0$ : optimum $(x_1,x_2)=(2,3)$, $z=-8$.
%	\end{exemple}
%	
%	\section{Méthode des deux phases}\label{sec:deux}
%	Si le PL contient des contraintes $\ge$ ou $=$, les variables d'écart ne forment plus une base de départ. On construit le \textbf{programme auxiliaire (PA)} : (1) rendre tous les $b_i\ge0$ ; (2) ajouter $+e_i$ pour $\le$ et $-e_i$ pour $\ge$ ; (3) ajouter une variable artificielle $t_i\ge0$ aux contraintes $\ge$ et $=$ ; (4) minimiser $w=\sum t_i$. La base initiale (écarts $\le$ et artificielles) est évidente. À l'optimum de la phase I : si $w^*>0$, le PL est \textbf{irréalisable} ; si $w^*=0$ et toutes les artificielles sont hors base, on les retire et on obtient une base réalisable pour la phase II ; si une artificielle est en base à la valeur $0$, on la remplace par une variable hors base non artificielle (valeur $0$).
%	\begin{exemple}
%		$\max2x_1+3x_2$ s.c. $x_1+x_2\le10$, $5x_1+4x_2\ge20$, $x\ge0$. Forme : $x_1+x_2+e_1=10$, $5x_1+4x_2-e_2+t_1=20$. PA : $\min t_1$. Base $\{e_1,t_1\}$ ; on annule le coût de $t_1$ dans la ligne $w$ :
%		\begin{center}\small
%			\begin{tabular}{c|ccccc|c}
%				&$x_1$&$x_2$&$e_1$&$e_2$&$t_1$&$b$\\\hline
%				$e_1$&1&1&1&0&0&10\\ $t_1$&\textbf{5}&4&0&$-1$&1&20\\\hline $w$&$-5$&$-4$&0&1&0&$-20$
%			\end{tabular}\ $\to$\ 
%			\begin{tabular}{c|ccccc|c}
%				&$x_1$&$x_2$&$e_1$&$e_2$&$t_1$&$b$\\\hline
%				$e_1$&0&$\frac15$&1&$\frac15$&$-\frac15$&6\\ $x_1$&1&$\frac45$&0&$-\frac15$&$\frac15$&4\\\hline $w$&0&0&0&0&1&0
%			\end{tabular}
%		\end{center}
%		$w^*=0$ avec base $\{e_1,x_1\}$ : $(4,0,6,0)$ est une solution de base réalisable ; on supprime $t_1$. \emph{Phase II} avec $c=(2,3,0,0)$ : on annule le coût de $x_1$ ($L_3\leftarrow L_3-2L_2$) : ligne $z$ $=(0,\frac75,0,\frac25\mid-8)$. $x_2$ entre (ratios $30$ et $5$), $x_1$ sort ; puis $e_2$ entre ($z$ vaut $15$), $e_1$ sort. Tableaux de la phase II :
%		\begin{center}\scriptsize
%			\begin{tabular}{c|cccc|c}
%				$II_0$&$x_1$&$x_2$&$e_1$&$e_2$&$b$\\\hline $e_1$&0&$\frac15$&1&$\frac15$&6\\ $x_1$&1&\textbf{$\frac45$}&0&$-\frac15$&4\\\hline $z$&0&$\frac75$&0&$\frac25$&$-8$
%			\end{tabular}\quad
%			\begin{tabular}{c|cccc|c}
%				$II_1$&$x_1$&$x_2$&$e_1$&$e_2$&$b$\\\hline $e_1$&$-\frac14$&0&1&\textbf{$\frac14$}&5\\ $x_2$&$\frac54$&1&0&$-\frac14$&5\\\hline $z$&$-\frac74$&0&0&$\frac34$&$-15$
%			\end{tabular}\quad
%			\begin{tabular}{c|cccc|c}
%				$II_2$&$x_1$&$x_2$&$e_1$&$e_2$&$b$\\\hline $e_2$&$-1$&0&4&1&20\\ $x_2$&1&1&1&0&10\\\hline $z$&$-1$&0&$-3$&0&$-30$
%			\end{tabular}
%		\end{center}
%		Optimum : $(x_1,x_2)=(0,10)$, $z=30$.
%	\end{exemple}
%	\textbf{Méthode du grand $M$.} On garde un seul programme : on pénalise les artificielles dans l'objectif par $-M$ ($M$ très grand) pour un max, $+M$ pour un min. Ici $\max2x_1+3x_2-Mt_1$ ; on élimine $t_1$ de la ligne $z$ (coûts $2+5M$, $3+4M$, $-M$, ... , second membre $20M$ en convention $-z$). Après le pivot $x_1$/$t_1$ le coût réduit de $t_1$ vaut $-\frac25-M<0$ : $t_1$ ne rentre plus et on retrouve les tableaux de la phase II.
%	\begin{center}\scriptsize
%		\begin{tabular}{c|ccccc|c}
%			&$x_1$&$x_2$&$e_1$&$e_2$&$t_1$&$b$\\\hline $e_1$&1&1&1&0&0&10\\ $t_1$&\textbf{5}&4&0&$-1$&1&20\\\hline $z$&$2+5M$&$3+4M$&0&$-M$&0&$20M$
%		\end{tabular}\ $\to$\ 
%		\begin{tabular}{c|ccccc|c}
%			&$x_1$&$x_2$&$e_1$&$e_2$&$t_1$&$b$\\\hline $e_1$&0&$\frac15$&1&$\frac15$&$-\frac15$&6\\ $x_1$&1&$\frac45$&0&$-\frac15$&$\frac15$&4\\\hline $z$&0&$\frac75$&0&$\frac25$&$-\frac25-M$&$-8$
%		\end{tabular}
%	\end{center} Si à l'optimum une artificielle reste en base avec une valeur \emph{strictement positive}, le PL est irréalisable.
%	
%	\section{Cas particuliers}
%	\begin{center}\small
%		\begin{tabular}{|p{4.6cm}|p{9.2cm}|}\hline
%			\textbf{Ce qu'on observe dans le tableau} & \textbf{Conclusion}\\\hline
%			Coûts réduits tous $\le0$ (max), variables artificielles hors base & optimum atteint (phase II terminée)\\
%			Coût réduit $>0$ et colonne entièrement $\le0$ & problème \emph{non borné}\\
%			Optimum de phase I avec artificielle $>0$ en base & problème \emph{irréalisable}\\
%			Coût réduit nul d'une variable hors base à l'optimum & \emph{solutions optimales multiples}\\
%			Second membre nul d'une variable de base & solution \emph{dégénérée} (risque de cyclage : règle de Bland)\\\hline
%		\end{tabular}
%	\end{center}
%	\textbf{Optimum infini.} Dès qu'une variable entrante a tous ses coefficients $a_{ie}\le0$ dans la colonne, le ratio ne s'applique plus : la variable peut croître indéfiniment et $f$ est non bornée. Exemple : $\min-2x_1-3x_2+x_3$ s.c. $-x_1-x_2-x_3\le3$, $x_1-x_2+x_3\le4$, $-x_1+x_2+2x_3\le1$. Après l'entrée de $x_2$ (ligne 3), les colonnes de $x_1$ valent $(-2,0,-1)$, la ligne $z$ a le coût $-5$ pour $x_1$ : $x_1$ entre mais aucun ratio n'existe ; le rayon $(t,1+t,0)$ est admissible et $f=-3-5t\to-\infty$.
%	\begin{center}\scriptsize
%		\begin{tabular}{c|cccccc|c}
%			$U_0$&$x_1$&$x_2$&$x_3$&$e_1$&$e_2$&$e_3$&$b$\\\hline $e_1$&$-1$&$-1$&$-1$&1&0&0&3\\ $e_2$&1&$-1$&1&0&1&0&4\\ $e_3$&$-1$&\textbf{1}&2&0&0&1&1\\\hline $z$&$-2$&$-3$&1&0&0&0&0
%		\end{tabular}\ $\to$\ 
%		\begin{tabular}{c|cccccc|c}
%			$U_1$&$x_1$&$x_2$&$x_3$&$e_1$&$e_2$&$e_3$&$b$\\\hline $e_1$&$-2$&0&1&1&0&1&4\\ $e_2$&0&0&3&0&1&1&5\\ $x_2$&$-1$&1&2&0&0&1&1\\\hline $z$&$-5$&0&7&0&0&3&3
%		\end{tabular}
%	\end{center}
%	La colonne de $x_1$ dans $U_1$ est $(-2,0,-1)\le0$.
%	
%	\textbf{Problème impossible.} $\max4x_1+3x_2$ s.c. $x_1+x_2\le2$, $3x_1+x_2\ge10$. Avec $\max4x_1+3x_2-Mt$, après le pivot $x_1$ ($e_1$ sort), la ligne de $t$ est $(0,-2,-3,-1,1\mid4)$ et les coûts réduits $(-1-2M,\,-4-3M,\,-M)$ sont tous négatifs : le tableau est optimal \emph{avec $t=4>0$ en base} : $\mathbf D_R=\emptyset$.
%	\begin{center}\scriptsize
%		\begin{tabular}{c|ccccc|c}
%			&$x_1$&$x_2$&$e_1$&$e_2$&$t$&$b$\\\hline $e_1$&1&1&1&0&0&2\\ $t$&3&1&0&$-1$&1&10\\\hline $z$&$4+3M$&$3+M$&0&$-M$&0&$10M$
%		\end{tabular}\ $\to$\ 
%		\begin{tabular}{c|ccccc|c}
%			&$x_1$&$x_2$&$e_1$&$e_2$&$t$&$b$\\\hline $x_1$&1&1&1&0&0&2\\ $t$&0&$-2$&$-3$&$-1$&1&4\\\hline $z$&0&$-1-2M$&$-4-3M$&$-M$&0&$4M-8$
%		\end{tabular}
%	\end{center}
%	
%	\textbf{Solutions multiples.} Si à l'optimum un coût réduit d'une variable hors base est \emph{nul}, on peut la faire entrer sans changer $z$ : autre sommet optimal, donc tout le segment qui les joint. Exemple (exemple~\ref{ex:multiple}) : après l'entrée de $x_1$, le coût de $x_2$ est $0$ avec $-z=225$ ; faire entrer $x_2$ (ratios $\frac{15}2,5,\frac{15}4$ : $e_3$ sort) donne le second sommet optimal $(\frac94,\frac{15}4)$.
%	\begin{center}\scriptsize
%		\begin{tabular}{c|ccccc|c}
%			&$x_1$&$x_2$&$e_1$&$e_2$&$e_3$&$b$\\\hline $x_1$&1&$\frac35$&$\frac1{10}$&0&0&$\frac92$\\ $e_2$&0&$\frac{18}5$&$-\frac25$&1&0&18\\ $e_3$&0&\textbf{$\frac{24}5$}&$-\frac15$&0&1&18\\\hline $z$&0&0&5&0&0&$225$
%		\end{tabular}\ $\to$\ 
%		\begin{tabular}{c|ccccc|c}
%			&$x_1$&$x_2$&$e_1$&$e_2$&$e_3$&$b$\\\hline $x_1$&1&0&$\frac18$&0&$-\frac18$&$\frac94$\\ $e_2$&0&0&$-\frac14$&1&$-\frac34$&$\frac92$\\ $x_2$&0&1&$-\frac1{24}$&0&$\frac5{24}$&$\frac{15}4$\\\hline $z$&0&0&5&0&0&$225$
%		\end{tabular}
%	\end{center}
%	
%	\textbf{Dégénérescence et cyclage.} Si le plus petit ratio est atteint à plusieurs lignes, la solution suivante a une variable de base nulle (\emph{dégénérée}) ; plusieurs pivots peuvent laisser $z$ inchangé, et on peut même retrouver une base déjà visitée (\emph{cyclage}). Exemple : $\min-4x_1-3x_3$ s.c. $x_1-x_2\le2$, $2x_1+x_3\le4$, $x_1+x_2+x_3\le3$. L'entrée de $x_1$ a deux ratios égaux à $2$ ; en faisant sortir $e_1$, on obtient $e_2=0$ ($z=-8$) ; les deux pivots suivants ($x_2$ entre/$e_2$ sort, puis $x_3$ entre/$x_2$ sort) laissent $z=-8$ ; le pivot suivant ($e_1$ entre, $e_3$ sort) donne l'optimum $x=(1,0,2)$, $z=-10$.
%	\begin{center}\scriptsize
%		\begin{tabular}{c|cccccc|c}
%			$D_0$&$x_1$&$x_2$&$x_3$&$e_1$&$e_2$&$e_3$&$b$\\\hline $e_1$&\textbf{1}&$-1$&0&1&0&0&2\\ $e_2$&2&0&1&0&1&0&4\\ $e_3$&1&1&1&0&0&1&3\\\hline $z$&$-4$&0&$-3$&0&0&0&0
%		\end{tabular}\ 
%		\begin{tabular}{c|cccccc|c}
%			$D_1$&$x_1$&$x_2$&$x_3$&$e_1$&$e_2$&$e_3$&$b$\\\hline $x_1$&1&$-1$&0&1&0&0&2\\ $e_2$&0&\textbf{2}&1&$-2$&1&0&\textcolor{red}{0}\\ $e_3$&0&2&1&$-1$&0&1&1\\\hline $z$&0&$-4$&$-3$&4&0&0&8
%		\end{tabular}
%		\medskip
%		
%		\begin{tabular}{c|cccccc|c}
%			$D_2$&$x_1$&$x_2$&$x_3$&$e_1$&$e_2$&$e_3$&$b$\\\hline $x_1$&1&0&$\frac12$&0&$\frac12$&0&2\\ $x_2$&0&1&\textbf{$\frac12$}&$-1$&$\frac12$&0&\textcolor{red}{0}\\ $e_3$&0&0&0&1&$-1$&1&1\\\hline $z$&0&0&$-1$&0&2&0&8
%		\end{tabular}\ 
%		\begin{tabular}{c|cccccc|c}
%			$D_3$&$x_1$&$x_2$&$x_3$&$e_1$&$e_2$&$e_3$&$b$\\\hline $x_1$&1&$-1$&0&1&0&0&2\\ $x_3$&0&2&1&$-2$&1&0&\textcolor{red}{0}\\ $e_3$&0&0&0&\textbf{1}&$-1$&1&1\\\hline $z$&0&2&0&$-2$&3&0&8
%		\end{tabular}
%		\medskip
%		
%		\begin{tabular}{c|cccccc|c}
%			$D_4$&$x_1$&$x_2$&$x_3$&$e_1$&$e_2$&$e_3$&$b$\\\hline $x_1$&1&$-1$&0&0&1&$-1$&1\\ $x_3$&0&2&1&0&$-1$&2&2\\ $e_1$&0&0&0&1&$-1$&1&1\\\hline $z$&0&2&0&0&1&2&10
%		\end{tabular}
%	\end{center}
%	($D_1$, $D_2$, $D_3$ : bases dégénérées (second membre nul en rouge), $z$ reste égal à $-8$ ; $D_4$ est optimal.)
%	\begin{remarque}[Règle de Bland]
%		Pour éviter le cyclage : parmi les candidats, entrer la variable de plus petit indice, et sortir la variable de plus petit indice parmi les ex æquo du ratio minimal.
%	\end{remarque}
%	\begin{retenir}
%		Max : entrer le coût réduit $>0$ maximal, optimal si tous $\le0$. Min : entrer le coût $<0$ minimal, optimal si tous $\ge0$. Sortir le plus petit ratio $b_i/a_{ie}$ avec $a_{ie}>0$. Colonne entrante $\le0$ : non borné. Artificielle $>0$ à l'optimum : irréalisable. Coût réduit nul à l'optimum : solutions multiples.
%	\end{retenir}

	\chapter{Résolution algébrique}

Pour résoudre des problèmes de programmation linéaire ayant un nombre de variables (principales) supérieur à trois, nous devons obligatoirement faire appel à une autre méthode de résolution que celle déjà étudiée. Toutefois, avant d'aborder la méthode du simplexe comme technique de résolution d'un problème de programmation linéaire, nous allons découvrir une version modifiée "la méthode algébrique" qui permettra d'étudier les différentes étapes à franchir pour arriver à une solution optimale.

\section{Variable d'écarts}

La méthode de résolution que nous étudions dans ce chapitre nécessite que les contraintes du modèle soient exprimées sous forme d'équations linéaires au lieu d'inéquations.

On peut facilement transformer une inéquations linéaire ayant un signe $\leq$ en une équation linéaire en additionnant une variable non négative dite variable d'écart.

\begin{proposition}
	Tout PL sous forme standard s'écrit de façon équivalente en un PL sous forme canonique pure et inversement.
\end{proposition}
\label{prop:2.1}

%	\section{Démonstration}
\begin{proof}
	\begin{enumerate}
		\item Soit un PL sous forme canonique pure. On a
		$$
		\begin{aligned}
			A x \leq b \Longleftrightarrow & \displaystyle\sum_{j=1}^{n} a_{i j} x_{j}+e_{i}=b_{i}, \quad \forall i=1, \ldots, m \\
			& \text { où } e_{i}=b_{i}-\displaystyle\sum_{j=1}^{n} a_{i j} x_{j} \geq 0
		\end{aligned}
		$$
		
		Ainsi,
		
		$$
		\left\{\begin{array}{r} 
			A x \leq b , \\
			x \geq 0 .
		\end{array} \Leftrightarrow \left\{\begin{array}{r} 
			(A \mid I_{m}) \begin{pmatrix} x \\ e \end{pmatrix} = b , \\
			\begin{pmatrix} x \\ e \end{pmatrix} \geq 0 .
		\end{array} \Leftrightarrow \left\{\begin{array}{r}
			\tilde{A} \tilde{x}=b \\
			\tilde{x} \geq 0 .
		\end{array}\right.\right.\right.
		$$
		
		où $e=(e_{1}, \ldots, e_{m})^{t}$ sont appelées variables d'écart et $\tilde{A}$ est une matrice de taille $m \times n+m$.\\
		\item  Soit un PL sous forme standard. On a
		
		$$
		A x=b \Longleftrightarrow\left\{\begin{array}{c} 
			A x \leq b , \\
			A x \geq b .
		\end{array} \Longleftrightarrow \left\{\begin{array}{c}
			A x \leq b, \\
			-A x \leq-b .
		\end{array} \Longleftrightarrow\begin{pmatrix} A \\ -A \end{pmatrix} x \leq\begin{pmatrix} b \\ -b \end{pmatrix} \Longleftrightarrow \tilde{\tilde{A}} x \leq \tilde{\tilde{b}}\right.\right.
		$$
		
		où $\tilde{\tilde{A}}$ est une matrice de taille $2 m \times n$ et $\tilde{\tilde{b}} \in \mathbb{R}^{2 m}$.
	\end{enumerate}
\end{proof}

\begin{remarque}
	Pour transformer un programme linéaire sous forme canonique en sa forme standard (où toutes les contraintes sont des égalités avec un second membre positif), on introduit des variables supplémentaires non-négatives. La méthode varie selon le type d'inégalité :
	\begin{itemize}
		\item \textbf{Pour une contrainte de type "inférieur ou égal" ($\leq$)} : On ajoute une \textbf{variable d'écart} ($e_i \geq 0$) pour transformer l'inégalité en une égalité.
		\[ a_{i,1}x_1 + \dots + a_{i,n}x_n \leq b_i \quad \Rightarrow \quad a_{i,1}x_1 + \dots + a_{i,n}x_n + e_i = b_i \]
		
		\item \textbf{Pour une contrainte de type "supérieur ou égal" ($\geq$)} : On soustrait une \textbf{variable d'excès} ($e_i \geq 0$) et on ajoute une \textbf{variable artificielle} ($r_i \geq 0$).
		\[ a_{i,1}x_1 + \dots + a_{i,n}x_n \geq b_i \quad \Rightarrow \quad a_{i,1}x_1 + \dots + a_{i,n}x_n - e_i + r_i = b_i \]
		
		\item \textbf{Pour une contrainte de type "égal" (=)} : On ajoute simplement une \textbf{variable artificielle} ($r_i \geq 0$).
		\[ a_{i,1}x_1 + \dots + a_{i,n}x_n = b_i \quad \Rightarrow \quad a_{i,1}x_1 + \dots + a_{i,n}x_n + r_i = b_i \]
	\end{itemize}
	Les \textbf{variables d'écart et d'excès} représentent la différence entre les deux membres de l'inéquation. Les \textbf{variables artificielles} sont des outils de calcul temporaires, indispensables pour démarrer l'algorithme du simplexe. Elles doivent impérativement être nulles dans la solution finale, ce qui est généralement réalisé à l'aide de la méthode des deux phases.
\end{remarque}

\section{Solutions de base réalisables}

On considère désormais (sauf mention contraire) un PL toujours sous forme standard.

\begin{definition}
	On appelle solution réalisable tout vecteur $x$ qui satisfait les contraintes du PL i.e. tel que $A x=b$ et $x \geq 0$.
\end{definition}
\label{def:2.1}

\begin{hypothese}
	On suppose désormais que la matrice $A$ est de taille $m \times n$ avec $\operatorname{rang}(A)=m \leq n$.\\
	On rappelle que le rang de $A$ est le nombre maximal de lignes de $A$ linéairement indépendantes. C'est aussi le nombre de colonnes de $A$ linéairement indépendantes. L'hypothèse de rang plein n'est pas restrictive car :
\end{hypothese}

\begin{itemize}
	\item si $\operatorname{rang}(A)<m$ le système $A x=b$ n'a pas de solution en général.
	\item Si $\operatorname{rang}(A)<m$ et $b \in \operatorname{Im}(A)$, il y a des équations redondantes dans le système $A x=b$, qu'on peut donc supprimer pour obtenir un nouveau système de rang plein. Sous l'hypothèse de rang plein, le système $A x=b$ admet toujours des solutions.
	\item Si $m<n$, il y en a une infinité.
	\item Si $m=n$, la solution est unique et vaut $x=A^{-1} b$, dans ce cas, il n'y a rien à maximiser...
\end{itemize}

\begin{remarque}
	Pour la forme canonique pure avec $A x \leq b$, l'hypothèse de rang plein n'est pas non plus restrictive car si $\operatorname{rang}(A)<m$, il y a des contraintes d'inégalités redondantes qu'on peut supprimer pour obtenir un système de rang plein.
\end{remarque}
\label{rem:2.1}

\begin{definition}
	Soit $B \subseteq\{1, \ldots, n\}$ un ensemble d'indices avec $\operatorname{card}(B)=m$ tel que les colonnes $A_{j}, j \in B$, de $A$ sont linéairement indépendantes. Autrement dit, la matrice carrée $A_{B}$ formée des colonnes $A_{j}$, $j \in B$, est inversible. On dit que l'ensemble $B$ des indices est une base.
\end{definition}
\label{def:2.2}

\begin{itemize}
	\item Les variables $x_{B}=\left(x_{j}, j \in B\right)$ sont appelées variables de base.
	\item Les variables $x_{H}=\left(x_{j}, j \notin B\right)$ sont appelées variables hors-base.
\end{itemize}

\begin{remarque}
	\begin{itemize}
		\item Sous l'hypothèse de rang plein, il existe toujours une base non vide.
		\item Quitte à re-numéroter les indices, on peut toujours écrire les décompositions par blocs :
	\end{itemize}
\end{remarque}
\label{rem:2.2}

$$
\begin{aligned}
	A & =\left(A_{B} \mid A_{H}\right) \text { où } A_{H} \text { est la matrice formée des colonnes } A_{j}, j \notin B \\
	x & =\left(x_{B} \mid x_{H}\right)^{t}
\end{aligned}
$$

Le système $A x=b$ est équivalent à $A_{B} x_{B}+A_{H} x_{H}=b$.

Par la relation précédente et du fait que la matrice $A_{B}$ est inversible, on peut fixer les variables hors-base et les variables de base sont alors complètement déterminées.

\begin{definition}
	On dit que $x=\left(x_{B} \mid x_{H}\right)^{t}$ est solution de base associée à la base $B$ si $$ x_{H}=0 $$.
\end{definition}
\label{def:2.3}

\textbf{Propriétés des solutions de base réalisables :} Si $x=\left(x_{B} \mid x_{H}\right)^{t}$ est une solution de base réalisable alors $x_{H}=0$ et $x_{B}=A_{B}^{-1} b$.

%	\textbf{Exemple 2.1} 
%\begin{exemple}
%	Reprenons l'exemple du problème de production de l'introduction. Sous forme standard (en introduisant des variables d'écart), le PL s'écrit
%	
%	$$
%	\left\{\begin{array}{rl} 
%		\displaystyle\max_{( x_{1} , x_{2} )} & f ( x_{1} , x_{2} ) = 6 x_{1} + 4 x_{2}  \\
%		\text { s.t. }  & 3 x_{1} + 9 x_{2} \leq 81  \\
%		& 4 x_{1} + 5 x_{2} \leq 55  \\
%		& 2 x_{1} + x_{2} \leq 20  \\
%		& x_{1} , x_{2} \geq 0 
%	\end{array} \right. \Longrightarrow \left\{\begin{array}{rl} 
%		\displaystyle\max_{( x_{1} , x_{2} , e_{1} , e_{2} , e_{3} )} & f ( x_{1} , x_{2} , e_{1} , e_{2} , e_{3} ) = 6 x_{1} + 4 x_{2} + 0 e_{1} + 0 e_{2} + 0 e_{3}  \\
%		\text { s.t. } & 3 x_{1} + 9 x_{2} + e_{1} = 81  \\
%		& 4 x_{1} + 5 x_{2} + e_{2} = 55  \\
%		& 2 x_{1} + x_{2} + e_{3} = 20  \\
%		& x_{1} , x_{2} , e_{1} , e_{2} , e_{3} \geq 0 
%	\end{array} \right.
%	$$
%	
%	On a $m=3, n=5$ et $\operatorname{rang}(A)=m=3$.
%	
%	Nous sommes alors en présence d'un système de plus d'inconnues que d'équations.
%	
%	Rappelons qu'il est possible de résoudre un système d'équations ayant plus d'inconnues que d'équations ( $m \leq n$ ) en fixant arbitrairement (par une valeur ou un paramètre) les valeurs de $n-m$ variables.
%	
%	\begin{description}
%		\item[Programme initial] \quad\\	
%		Pour déterminer le programme initial, on pose habituellement à zéro les variables principales du modèle, ceci correspond à ne fabriquer aucun produit :
%		
%		$$
%		x_{1}=0 \text { et } x_{2}=0 .
%		$$
%		
%		Notre système de 3 équations à 5 inconnues devient alors un système de 3 équations à 3 inconnues que l'on va pouvoir résoudre. Une base est donnée par $B=\{3,4,5\}$ avec $A_{B}=\left(\begin{array}{lll}1 & 0 & 0 \\ 0 & 1 & 0 \\ 0 & 0 & 1\end{array}\right)$. La première solution de base réalisable correspondante est $x=\left(x_{1}, x_{2}, e_{1}, e_{2}, e_{3}\right)^{t}=(\underbrace{0,0}_{x_{H}}, \underbrace{81,55,20}_{x_{B}=A_{B}^{-1} b})^{t}$, la valeur de la fonction économique est : 0. La première solution de base réalisable correspond
%		
%		\begin{figure}[H]
%			\begin{center}
%				\includegraphics[scale=0.35]{2025_10_05_27e91bcb17e07a381837g-07}
%				%			\captionsetup{labelformat=empty}
%				\caption{Résolution graphique du problème de production}
%			\end{center}
%		\end{figure}
%		
%		au point extrême $(0,0)^{t}$ de la résolution graphique. On peut évidemment améliorer la valeur de la fonction économique en fabriquant quelques unités du produit $P_{1}$ ou du produit $P_{2}$.
%		
%		\item[Révision du programme initial] \quad\\
%		On doit maintenant examiner la possibilité d'améliorer la valeur de la fonction économique en introduisant l'une ou l'autre des variables principales (soit $x_{1}$, soit $x_{2}$ mais non les deux simultanément) dans le programme de base.\\
%		Le choix de la variable (actuellement hors base) à introduire dans le programme s'effectue à l'aide d'un critère qui permet d'améliorer le plus rapidement possible la valeur de la fonction\\
%		économique, c-à-d on choisira la variable qui donne la meilleure augmentation de la valeur de la fonction économique. Dans notre exemple la fonction économique est
%		
%		$$
%		f\left(x_{1}, x_{2}\right)=6 x_{1}+4 x_{2}+0 e_{1}+0 e_{2}+0 e_{3}
%		$$
%		
%		l'augmentation sera la plus élevée si nous introduisons la variable $x_{1}$ dans le programme, Donc : $x_{1}$ est une variable entrante de base\\
%		Il nous reste maintenant deux choses à déterminer. D'une part, calculer la quantité du $P_{1}$ que l'on doit fabriquer et d'autre part, quelle variable actuellement dans le programme doit devenir une variable hors programme c.-à-d. déterminer la variable sortante.
%		
%		\item[Détermination de la valeur entrante] \quad\\			
%		Déterminons la plus grande valeur pouvant être prise par $x_{1}$ dans chaque équation.
%		
%		\begin{align}
%			x_{1} & =27-3 x_{2}-\dfrac{1}{3} e_{1}  \label{2.44}\\
%			x_{1} & =\dfrac{55}{4}-\dfrac{5}{4} x_{2}-\dfrac{1}{4} e_{2}  \label{2.5}\\
%			x_{1} & =10-\dfrac{1}{2} x_{2}-\dfrac{1}{2} e_{3} \label{2.6}
%		\end{align}
%		
%		
%		Pour s'assurer que toutes les contraintes soient respectées en même temps, on choisit parmi ces valeurs la plus petite quantité positive :
%		
%		$$
%		x_{1}=10 .
%		$$
%		
%		Ce minimum s'obtient dans l'équation \eqref{2.6}, ce qui se traduit par $e_{3}=0$ et donc c'est la variable $e_{3}$ qui sort du programme :
%		
%		\begin{equation}
%			e_{3}\text{ est une variable hors base.}
%		\end{equation}
%		
%		\item[Détermination de la nouvelle solution : solution de base № 2]\quad\\
%		Une deuxième base est donnée par $B=\{1,3,4\}$ avec $A_{B}=\left(\begin{array}{lll}3 & 1 & 0 \\ 4 & 0 & 1 \\ 2 & 0 & 0\end{array}\right)$. La première solution de base réalisable correspondante est $x=\left(x_{1}, x_{2}, e_{1}, e_{2}, e_{3}\right)^{t}=(\underbrace{10}_{x_{B}}, \underbrace{0}_{x_{H}}, \underbrace{51,15}_{x_{B}}, \underbrace{0}_{x_{H}})^{t}$, la valeur de la fonction économique est : 60.\\
%		Il correspond au point-sommet $(10,0)^{t}$ de la résolution graphique. S'agit-il d'une solution optimal? Toutefois, avant de tester si le programme est optimal, il faut d'abord transformer les contraintes de mon problème pour exprimer les variables de base en fonction des variables hors-base:
%		
%		\begin{itemize}
%			\item les variables de base sont : $x_{1}, e_{1}, e_{2}$.
%			\item les variables hors-base sont : $x_{2}, e_{3}$.
%		\end{itemize}
%		
%		\begin{align}
%			\left\{\begin{array}{l}e_{1}=81-3 x_{1}-9 x_{2} \\ e_{2}=55-4 x_{1}-5 x_{2} \\ e_{3}=20-2 x_{1}-x_{2}\end{array}\right.& \Rightarrow\left\{\begin{array}{l}e_{1}=81-3\left(10-\dfrac{1}{2} x_{2}-\dfrac{1}{2} e_{3}\right)-9 x_{2} \\ e_{2}=55-4\left(10-\dfrac{1}{2} x_{2}-\dfrac{1}{2} e_{3}\right)-5 x_{2} \\ x_{1}=10-\dfrac{1}{2} x_{2}-\dfrac{1}{2} e_{3}\end{array}\right.\nonumber\\
%			&\Rightarrow \begin{cases}e_{1} & =51-\dfrac{15}{2} x_{2}+\dfrac{3}{2} e_{3} \\ e_{2} & =15-3 x_{2}-2 e_{3} \\ x_{1} & =10-\dfrac{1}{2} x_{2}-\dfrac{1}{2} e_{3}\end{cases}\label{equ2.7}
%		\end{align}
%		%			\[
%		%			
%		%			\]
%		%			
%		%			\[ \]
%		
%		Il nous reste à exprimer la valeur de la fonction économique, en fonction des variables hors-base $x_{2}$ et $e_{3}$ à l'aide de la substitution :
%		
%		$$
%		f_{o b j}=6 x_{1}+4 x_{2}=60+x_{2}-3 e_{3} .
%		$$
%		
%		Cette dernière expression nous indique qu'on peut améliorer encore la valeur de la fonction économique puisque le coefficient de la variable $x_{2}$ est positif.
%		
%		\item[Détermination de la troisième base : variable entrante est $x_{2}$]\quad\\
%		Déterminons la valeur de $x_{2}$, par le même principe, on obtient que la variable $x_{2}$ peut prendre la valeur $20, \dfrac{102}{15}$ ou 5 , Ainsi $x_{2}$ devant respecter les différentes contraintes, nous choisirons $x_{2}=5$ qui se réalise dans l'équation $e_{2}=15-3 x_{2}+2 e_{3}$ et qui impose que $e_{2}$ devienne nulle, ainsi : $e_{2}$ devient variable hors base.
%		
%		\begin{itemize}
%			\item les variables de base sont : $x_{1}, e_{1}, x_{2}$.
%			\item les variables hors-base sont : $e_{2}, e_{3}$.
%		\end{itemize}
%		
%		Exprimons à nouveau les variables dans le programme en fonction des variables hors-base.
%		
%		$$
%		\left\{\begin{aligned}
%			x_{1} & =\dfrac{15}{2}+\dfrac{1}{6} e_{2}-\dfrac{1}{3} e_{3} \\
%			e_{1} & =\dfrac{27}{2}+\dfrac{5}{2} e_{2}-\dfrac{7}{2} e_{3} \\
%			x_{2} & =5-\dfrac{1}{3} e_{2}+\dfrac{2}{3} e_{3}
%		\end{aligned}\right.
%		$$
%		
%		La valeur de la fonction économique est alors :
%		
%		$$
%		f_{o b j}=6 \times \dfrac{15}{2}+4 \times 5=65
%		$$
%		
%		Puis on exprime la valeur de la fonction économique, en fonction des variables hors-base $e_{2}$ et $e_{3}$ à l'aide de la substitution :
%		
%		$$
%		\begin{aligned}
%			f_{o b j} & =6 \times\left(\dfrac{15}{2}+\dfrac{1}{6} e_{2}-\dfrac{1}{3} e_{3}\right)+4 \times\left(5-\dfrac{1}{3} e_{2}+\dfrac{2}{3} e_{3}\right), \\
%			& =65-\dfrac{1}{3} e_{2}-\dfrac{7}{3} e_{3} .
%		\end{aligned}
%		$$
%		
%		Pouvons-nous améliorer encore la valeur de la fonction économique?\\
%		La réponse est non puisque les coefficients des variables hors-base $e_{2}$ et $e_{3}$ sont négatifs.
%	\end{description}
%\end{exemple}
\begin{remarque}
	Un programme de base est optimal si tous les coefficients des variables hors-base, dans l'expression de la fonction économique, sont $\leq 0$.
\end{remarque}
\label{rem:2.3}

On constate qu'elle correspond bien à la solution optimale obtenue lors de la résolution graphique.

\begin{remarque}
	Il y a au plus $C_{n}^{m}$ solutions de base (toutes ne sont pas réalisables).
\end{remarque}
\label{rem:2.4}

\section{Propriétés géométriques des solutions de base réalisables}

On note

\begin{equation*}
	\mathbf{D}_{R}=\left\{x \in \mathbb{R}^{n} \mid A x=b, x \geq 0\right\}, \tag{2.8}
\end{equation*}

l'ensemble des solutions réalisables d'un PL sous forme standard.\\
Commençons par rappeler les notions d'ensemble convexe et de polyèdre :\\
\begin{minipage}[t]{0.45\textwidth}
	\centering
	\includegraphics[scale=0.35]{2025_10_05_27e91bcb17e07a381837g-10(1)}
\end{minipage}
\begin{minipage}[t]{0.45\textwidth}
	\centering
	\includegraphics[scale=0.35]{2025_10_05_27e91bcb17e07a381837g-10}
\end{minipage}

\begin{definition}
	Un ensemble $E$ est dit convexe si $\forall x, y \in E, \lambda x+(1-\lambda) y \in E$ pour tout $0 \leq \lambda \leq 1$.
\end{definition}
\label{def:2.4}

\begin{definition}
	Un polyèdre $P \subseteq \mathbb{R}^{n}$ est l'ensemble de solutions d'un système fini d'inégalités linéaires, i.e.,
	
	$$
	P=\left\{x \in \mathbb{R}^{n} \mid A x \leq b\right\}
	$$
	
	où $A$ est une matrice $m \times n$ et $b \in \mathbb{R}^{m}$, et $m$ et $n$ sont deux entiers positifs.
	
	Un polyèdre borné est appelé polytope. En d'autres termes, un polyèdre $P \subseteq \mathbb{R}^{n}$ est un polytope s'il existe $x, y \in \mathbb{R}^{n}$ tel que $x \leq z \leq y$ pour tout $z \in P$.
	
	Noter que tout polyèdre est convexe, c'est-à-dire si $x$ et $y$ sont deux points quelconques de $P$, alors $\lambda x+(1-\lambda) y \in P$ pour tout $0 \leq \lambda \leq 1$.
\end{definition}
\label{def:2.5}

\begin{remarque}
	L'ensemble $\mathbf{D}_{R}$ des solutions réalisables est un polyèdre convexe, fermé.
\end{remarque}
\label{rem:2.5}

%	\textbf{Exemple 2.2} 
\begin{exemple}
	L'ensemble $\mathbf{D}_{R}=\left\{x \in \mathbb{R}^{3} \mid 2 x_{1}+3 x_{2}+x_{3}=3, x_{1}, x_{2}, x_{3} \geq 0\right\}$ est représenté sur la figure \ref{2.4}.
\end{exemple}

\begin{definition}
	Un modèle mathématique est dit non borné si son domaine réalisable est non borné, ç-à-d si au moins l'une de ses variables peut tendre vers l'infini sans quitter le domaine réalisable.
\end{definition}
\label{def:2.6}

\begin{definition}
	Un point $x \in \mathbf{D}_{R}$ est un sommet (ou point extrême) si et seulement s'il n'existe pas $y, z \in \mathbf{D}_{R}, y \neq z$ tels que $x=\lambda y+(1-\lambda) z$ avec $0<\lambda<1$.
\end{definition}

\begin{theorem}
	$x$ est une solution de base réalisable si et seulement si $x$ est un sommet de $\mathbf{D}_{R}$.
\end{theorem}
\label{thm:2.1}

On a vu sur l'exemple de production que la solution était atteinte sur un sommet de $\mathbf{D}_{R}$ et correspond donc à une solution de base.

\begin{figure}[h]
	\centering
	\begin{tikzpicture}[scale=1.2,
		x={(-0.707cm,-0.707cm)}, y={(1cm,0cm)}, z={(0cm,1cm)}]
		
		% Définir les sommets
		\coordinate (O) at (0,0,0);
		\coordinate (A) at (2,0,0);
		\coordinate (B) at (0,2,0);
		\coordinate (C) at (0,0,3);
		
		% Axes
		\draw[->, thick] (O) -- (3,0,0) node[anchor=north east]{$x_1$};
		\draw[->, thick] (O) -- (0,3,0) node[anchor=north west]{$x_2$};
		\draw[->, thick] (O) -- (0,0,4) node[anchor=south]{$x_3$};
		
		% Triangle avec hachures diagonales
		\fill[pattern=north east lines, pattern color=gray!40] (A) -- (B) -- (C) -- cycle;
		\draw[thick] (A) -- (B) -- (C) -- cycle;
		
		% Graduations sur les axes
		\foreach \x in {1, 2} {
			\draw (\x,0,0.1) -- (\x,0,-0.1) node[below] {\small \x};
		}
		\foreach \y in {1, 2} {
			\draw (0.1,\y,0) -- (-0.1,\y,0) node[below left] {\small \y};
		}
		\foreach \z in {1, 2, 3} {
			\draw (0,0.1,\z) -- (0,-0.1,\z) node[left] {\small \z};
		}
		
	\end{tikzpicture}
	\caption{Polyèdre des solutions réalisables}
	\label{2.4}
\end{figure}

\begin{proposition}
	Si le domaine réalisable $\mathbf{D}_{R}$ n'est pas vide, alors ce polytope comprend au moins un point extrême.
\end{proposition}
\label{prop:2.2}

Le résultat général s'énonce ainsi :

\begin{theorem}
	L'optimum de la fonction objectif $f$ sur $\mathbf{D}_{R}$, s'il existe, est atteint en au moins un sommet de $\mathbf{D}_{R}$.
\end{theorem}
\label{thm:2.2}

%	\textbf{Démonstration} Voir les propriétés géométriques des polyèdres.

Pour résoudre un PL sous forme standard, il suffit de se restreindre aux solutions de base réalisables (les sommets de $\mathbf{D}_{R}$). Tout se passe donc avec les solutions de base.

L'ensemble $\mathbf{D}_{R}$ n'est pas nécessairement borné. En fait pour un PL, 3 situations (et seulement 3) peuvent se produire :

\begin{enumerate}
	\item $\mathbf{D}_{R} = \varnothing$ : le PL n'a pas de solution.
	\item $\mathbf{D}_{R} \neq \varnothing$ mais la fonction objectif $f$ n'est pas majorée sur $\mathbf{D}_{R}$ : le maximum de $f$ vaut $+\infty$. Si $\mathbf{D}_{R}$ est borné, ce cas est exclu.
	\item $\mathbf{D}_{R} \neq \varnothing$ et la fonction objectif $f$ est majorée sur $\mathbf{D}_{R}$ : le PL admet une solution optimale (non nécessairement unique).
\end{enumerate}

%\begin{remarque}
%	Il y a au plus $C_{n}^{m}$ solutions de base réalisables. Pour déterminer une solution de base, on doit résoudre un système linéaire $(x_{B} = A_{B}^{-1} b)$. La résolution d'un système linéaire par une méthode directe de type Gauss/LU requiert de l'ordre de $m^{3}$ opérations. Si on explore toutes les solutions de base et que l'on compare les coûts correspondants, on effectue de l'ordre de $m^{3} C_{n}^{m}$ opérations. Ce nombre est vite très grand avec $n$ et $m$. Par exemple, avec $n=20$ et $m=10$, on a $3 \times 10^{8}$ opérations. Dans la méthode du simplexe, on va explorer seulement les sommets qui permettent d'augmenter la fonction objectif. On va réduire ainsi le nombre de solutions de base à explorer et donc le nombre de systèmes linéaires à résoudre.
%\end{remarque}
%\label{rem:2.6}

\section{Marche à suivre de la méthode algébrique}

La méthode algébrique est une recherche systématique de programmes de base (points sommets) jusqu'à l'obtention d'un programme optimal.

Pour ce faire, il s'agit :

\begin{colorframed}
	\begin{enumerate}
		\item De structurer le problème sous forme standard.
		\item De déterminer les variables de base qui serviront de départ au cheminement vers la solution optimale (programme optimal).
		\item D'expliciter la fonction économique et de déterminer si elle peut être améliorée : recherche de l'éventuelle variable (hors-base) admettant le plus grand coefficient positif.
		\item En introduisant cette variable dans le programme, on choisira la plus petite valeur positive obtenue à l'aide du système d'équations calculé lors de l'étape précédente. Cela induira la variable sortante.
		\item Pour déterminer une nouvelle variable de base, on doit exprimer la fonction économique ainsi que les nouvelles variables de base en fonction des variables hors-base.
		\item Retourner à 3) jusqu'à l'obtention du programme de base optimal.
		\item Donner le programme optimal en précisant la valeur de toutes les variables ainsi que la valeur optimisée de la fonction économique.
	\end{enumerate}
\end{colorframed}
\chapter{Méthode du simplexe}

La présentation algébrique de la méthode du simplexe nous a permis de donner une interprétation pratique et économique aux différentes opérations à effectuer dans la résolution d'un problème de programmation linéaire pouvant comporter deux variables ou plus. Toutefois, la méthode algébrique devient rapidement laborieuse à mesure que le nombre de contraintes et le nombre de variables augmentent.

Nous allons donc systématiser toutes ces opérations algébriques sous forme d'un simple tableau dit « tableau du simplexe ». Ceci permettra de résoudre plus facilement des problèmes ayant beaucoup plus de variables et de contraintes. Malgré l'efficacité de cette méthode, il n'en demeure pas moins que la résolution à l'aide d'un ordinateur devient indispensable à mesure que les problèmes prennent de fortes dimensions.

\section{Principe de l'algorithme du simplexe}

La méthode du simplexe est due à G. Dantzig (1947). Elle comporte 2 phases.
\begin{itemize}
	\item \textbf{Phase I - Initialisation} : Trouver une solution de base réalisable (ou bien détecter l'impossibilité).
	\item \textbf{Phase II - Progression} : On passe d'un sommet à un sommet voisin pour augmenter la fonction objectif (ou bien on détecte une fonction objectif $f$ non majorée).
\end{itemize}

\subsection{L'algorithme du simplexe : la phase 2}

Dans cette section, nous allons présenter la Phase II de la méthode du simplexe. La Phase I qui sert plus à initialiser la Phase II, sera abordée plus tard. Cette phase s'applique à des problèmes du type
\textcolor{red}{
\[
\left\{
\begin{array}{rl}
	\displaystyle\max & c^{t} x \\
	\text{s.t.} & A x \leq b \\
	& x \geq 0
\end{array}
\quad \text{où} \quad
\begin{array}{rl}
	\displaystyle\min_{x} & c^{t} x \\
	\text{s.t.} & A x \leq b \\
	& x \geq 0
\end{array}
\right.
\]
}
où $A$ est une matrice de taille $m \times n$.

On fera l'hypothèse que $b \geq 0$. Cette supposition est cruciale pour la Phase II. Ceci garantit que $0 \in K = (x \geq 0, A x \leq b)$. De plus, nous savons que le point $0$ est un sommet. Ce point servira de point de départ de l'algorithme du simplexe. En gros, l'algorithme va pivoter autour de ce point pour trouver un meilleur sommet. On poursuit l'algorithme jusqu'à l'obtention de la solution optimale.

Supposons qu'on dispose d'une solution de base réalisable $\bar{x}$ d'un PL sous forme standard. À une permutation près des colonnes, la matrice $A$ peut s'écrire
\[
A = \left( A_{B} \mid A_{H} \right)
\]
avec $A_{B}$ une matrice carrée de taille $m \times m$, inversible, correspondant aux variables de base et $A_{H}$ une matrice de taille $m \times (n-m)$, correspondant aux variables hors-base. On décompose également (à la même permutation près des composantes)
\[
\bar{x} = \left( \bar{x}_{B} \mid \bar{x}_{H} \right)^{t}
\]
Le but est de trouver une autre base $B^{*}$ et une solution de base $x^{*}$ associée telles que
\[
f(x^{*}) > f(\bar{x}) \quad \text{( $x^{*}$ est meilleur que $\bar{x}$ )}
\]
La méthode du simplexe consiste à faire entrer une variable hors-base dans la nouvelle base (variable entrante) et faire sortir à la place une variable de base (variable sortante).

\subsubsection{Calcul des coûts réduits et variable entrante}

On écrit la fonction objectif $f$ avec les variables de base/hors-base. On a $b = A x = A_{B} x_{B} + A_{H} x_{H}$ avec $A_{B}$ inversible donc $x_{B} = A_{B}^{-1} (b - A_{H} x_{H})$. On obtient donc
\[
\begin{aligned}
	f(x) &= c^{t} x = c_{B}^{t} x_{B} + c_{H}^{t} x_{H} \quad \text{avec} \quad c = \left( c_{B} \mid c_{H} \right)^{t} \\
	&= c_{B}^{t} A_{B}^{-1} (b - A_{H} x_{H}) + c_{H}^{t} x_{H} \\
	&= c_{B}^{t} A_{B}^{-1} b + \left( c_{H}^{t} - c_{B}^{t} A_{B}^{-1} A_{H} \right) x_{H}
\end{aligned}
\]
Or $x_{B} = A_{B}^{-1} b$ (car $x_{H} = 0$) et $c_{B}^{t} A_{B}^{-1} b = c^{t} x = f(x)$ donc
\[
f(x) = f(x) + \left( c_{H}^{t} - c_{B}^{t} A_{B}^{-1} A_{H} \right) x_{H}
\]
Le vecteur
\[
L_{H}^{t} = c_{H}^{t} - c_{B}^{t} A_{B}^{-1} A_{H}
\]
s'appelle vecteur des coûts réduits.

\begin{itemize}
	\item Si les coûts réduits sont tous négatifs, i.e. $L_{H}^{t} \leq 0$, il n'est alors pas possible d'augmenter la fonction objectif $f$ : l'algorithme se termine normalement, c'est-à-dire qu'on a trouvé une solution de base réalisable optimale.
	\item Dans le cas contraire (i.e. $\exists (L_{H})_{i} > 0$), on a intérêt à faire entrer dans la base la variable hors-base qui a le coût réduit positif le plus grand possible.
\end{itemize}

On note $e \in B$ l'indice de la variable entrante. On choisit $e$ tel que
\[
(L_{H})_{e} = \displaystyle\max_{j} \left\{ (L_{H})_{j}, (L_{H})_{j} > 0 \right\}
\]
ce qu'on note par
\[
e = \operatorname{argmax}_{j} \left\{ (L_{H})_{j}, (L_{H})_{j} > 0 \right\}
\]
\begin{remarque}
	Si on traite d'un problème de minimisation, c'est-à-dire avec $\displaystyle\min f(x)$, alors la variable entrante $x_{e}$ est déterminée par l'indice
	\[
	e = \operatorname{argmin}_{j} \left\{ (L_{H})_{j}, (L_{H})_{j} < 0 \right\}
	\]
\end{remarque}
\label{rem:3.1}

\subsubsection{Variable sortante}

Une fois l'indice $e$ choisi, il faut déterminer quelle variable doit quitter la base. En maintenant la relation $A x = b$, on augmente $x_{e}$ jusqu'à annuler une des variables de base. Cette variable sera alors la variable sortante.
\[
\begin{aligned}
	A x = b &\Longleftrightarrow A_{B} x_{B} + A_{e} x_{e} = b \quad \text{où $A_{e}$ désigne la $e$-ième colonne de $A$} \\
	&\Longleftrightarrow x_{B} = A_{B}^{-1} (b - A_{e} x_{e}) \\
	&\Longleftrightarrow x_{B} = \bar{x}_{B} - A_{B}^{-1} A_{e} x_{e} \\
	&\Longleftrightarrow x_{B} = \bar{x}_{B} - z x_{e}
\end{aligned}
\]
où
\[
z = A_{B}^{-1} A_{e} \in \mathbb{R}^{m}
\]
\begin{itemize}
	\item Si $z \leq 0$, on peut augmenter $x_{e}$ autant qu'on veut, on aura toujours la positivité de la variable de base $x_{B}$. La fonction objectif n'est pas majorée sur $\mathbf{D}_{R}$, i.e. le maximum de $f$ vaut $+\infty$. Dans ce cas, l'algorithme s'arrête.
	\item Sinon (i.e. il existe $z_{i} > 0$), pour avoir la positivité $(\bar{x}_{B})_{i} - z_{i} x_{e} \geq 0$ pour tout $i$, on choisit la variable sortante pour laquelle le rapport $\dfrac{(\bar{x}_{B})_{i}}{z_{i}}$ pour $i = 1, \ldots, m$ avec $z_{i} > 0$, est le plus petit possible. On choisit
	\[
	s = \operatorname{argmin}_{i} \left\{ \dfrac{(\bar{x}_{B})_{i}}{z_{i}}, z_{i} > 0 \right\}
	\]
	On a, dans ce cas, $x_{s} = 0$ et $x_{B} \geq 0$. La valeur de la variable entrante est donnée par
	\[
	x_{e} = \displaystyle\min_{i} \left\{ \dfrac{(\bar{x}_{B})_{i}}{z_{i}}, z_{i} > 0 \right\}
	\]
\end{itemize}

%	\section{}
\begin{colorframed}
	{\color{titleblue}\textbf{Algorithme (Résumé de la méthode du simplexe en phase 2 - progression)}}
	
	\medskip
	\begin{enumerate}
		\item \begin{itemize}
			\item Calcul des variables de base réalisables : étant donné $A = (A_{B} \mid A_{H})$, on calcule les variables de base réalisables $x_{B} = A_{B}^{-1} b \geq 0$.
			\item Calcul des coûts réduits : $L_{H}^{t} = c_{H}^{t} - c_{B}^{t} A_{B}^{-1} A_{H}$.
			\item Si $L_{H} \leq 0$ alors $\bar{x}_{B}$ est une solution optimale (arrêt de l'algorithme).
		\end{itemize}
		\item Variable entrante : $e = \operatorname{argmax}_{j} \left\{ (L_{H})_{j}, (L_{H})_{j} > 0 \right\}$.
		\item Variable sortante : Calcul de $z = A_{B}^{-1} A_{e}$ puis $s = \operatorname{argmin}_{i} \left\{ \dfrac{(\bar{x}_{B})_{i}}{z_{i}}, z_{i} > 0 \right\}$.
		\item On obtient une nouvelle base $B^{*}$ et une nouvelle matrice $A_{B^{*}}$ dans laquelle la colonne $A_{e}$ remplace la colonne $A_{s}$. Calcul de $A_{B^{*}}^{-1}$ et retour en 1.
	\end{enumerate}
\end{colorframed}

Nous allons maintenant systématiser cet algorithme en présentant la méthode du simplexe sous la forme d'un tableau.

\section{Représentation tabulaire du simplexe}

Afin de simplifier la présentation des différentes itérations de la méthode du simplexe, nous systématisons cette procédure sous forme de tableaux sur deux types de problèmes : de maximisation et de minimisation.

\subsection{Cas général de la méthode tabulaire du simplexe}

Considérons le programme linéaire suivant, exprimé sous sa forme standard, qui vise à maximiser une fonction objectif sous des contraintes d'égalité :

\[
\left\{
\begin{array}{rcl}
	\displaystyle \displaystyle\max_{\left(x_{1}, \ldots, x_{n}\right)} z & = & c_{1} x_{1} + c_{2} x_{2} + c_{3} x_{3} + \cdots + c_{n} x_{n} \\
	\text{s.t.} & & a_{11} x_{1} + a_{12} x_{2} + \cdots + a_{1n} x_{n} = b_{1} \\
	& & a_{21} x_{1} + a_{22} x_{2} + \cdots + a_{2n} x_{n} = b_{2} \\
	& & \vdots \\
	& & a_{m1} x_{1} + a_{m2} x_{2} + \cdots + a_{mn} x_{n} = b_{m} \\
	& & x_{1}, \ldots, x_{n} \geq 0
\end{array}
\right.
\]

Ce programme peut être reformulé de manière équivalente en introduisant la variable $z$ dans les contraintes, comme suit :

\begin{equation}\label{pro3.2}
	\left\{
	\begin{array}{rcl}
		\displaystyle \displaystyle\max z & & \\
		\text{s.t.} & & a_{11} x_{1} + a_{12} x_{2} + \cdots + a_{1n} x_{n} = b_{1} \\
		& & a_{21} x_{1} + a_{22} x_{2} + \cdots + a_{2n} x_{n} = b_{2} \\
		& & \vdots \\
		& & a_{m1} x_{1} + a_{m2} x_{2} + \cdots + a_{mn} x_{n} = b_{m} \\
		& & c_{1} x_{1} + c_{2} x_{2} + c_{3} x_{3} + \cdots + c_{n} x_{n} - z = 0 \\
		& & x_{1}, \ldots, x_{n} \geq 0
	\end{array}
	\right.
\end{equation}

Sans perte de généralité, nous pouvons supposer que les variables $x_{1}, \ldots, x_{m}$ sont les variables de base. Le problème \eqref{pro3.2} peut donc être représenté par le tableau suivant, dit tableau du simplexe :


\begin{center}
	\begin{tabular}{c|ccccccccc|c|c}
		\toprule
		Variables de base & \multicolumn{4}{c}{$\overbrace{\phantom{xxxxxxxxxxxxxxxxx}}^{\text{variable de base}}$} & \multicolumn{4}{c}{$\overbrace{\phantom{xxxxxxxxxxxxxxxxx}}^{\text{variable hors base}}$}& & & \\
		& $x_1$ & $x_2$ & $\cdots$ & $x_m$ & $x_{m+1}$ & $x_{m+2}$ & $\cdots$ & $x_n$ & $z$ & Termes de droite & $R$ \\
		\midrule
		$x_1$ & $a_{11}$ & $a_{12}$ & $\cdots$ & $a_{1m}$ & $a_{1m+1}$ & $a_{1m+2}$ & $\cdots$ & $a_{1n}$ & 0 & $b_1$ &\\
		$x_2$ & $a_{21}$ & $a_{22}$ & $\cdots$ & $a_{2m}$ & $a_{2m+1}$ & $a_{2m+2}$ & $\cdots$ & $a_{2n}$ & 0 & $b_2$ &\\
		$\vdots$ & $\vdots$ & $\vdots$ & $\ddots$ & $\vdots$ & $\vdots$ & $\vdots$ & $\ddots$ & $\vdots$ & $\vdots$ & $\vdots$& \\
		$x_m$ & $a_{m1}$ & $a_{m2}$ & $\cdots$ & $a_{mm}$ & $a_{mm+1}$ & $a_{mm+2}$ & $\cdots$ & $a_{mn}$ & 0 & $b_m$ &\\
		\hline
		$z$ & $c_1$ & $c_2$ & $\cdots$ & $c_m$ & $c_{m+1}$ & $c_{m+2}$ & $\cdots$ & $c_n$ & $-1$ & 0& \\
		\hline
	\end{tabular}
\end{center}

\subsection{Problème de maximisation}
Pour bien visualiser la similitude qui existe entre l'approche algébrique et celle des tableaux du simplexe, nous allons résoudre à nouveau le problème de production.\\

Le problème, après l'introduction des variables d'écart, était le suivant :

\begin{itemize}
	\item Maximiser la fonction économique :
	\[
	f(x_{1}, x_{2}, e_{1}, e_{2}, e_{3}) = 6x_{1} + 4x_{2} + 0e_{1} + 0e_{2} + 0e_{3}
	\]
	\item Sous les contraintes :
	\[
	\left\{
	\begin{aligned}
		3x_{1} + 9x_{2} + e_{1} &= 81 \\
		4x_{1} + 5x_{2} + e_{2} &= 55 \\
		2x_{1} + x_{2} + e_{3} &= 20 \\
		x_{1}, x_{2}, e_{1}, e_{2}, e_{3} &\geq 0
	\end{aligned}
	\right.
	\]
\end{itemize}

Que l'on peut écrire sous la forme :

\[
\left\{
\begin{array}{rcl}
	3x_1 + 9x_2 + e_1 + 0e_2 + 0e_3 & = & 81 \\
	4x_1 + 5x_2 + 0e_1 + e_2 + 0e_3 & = & 55 \\
	2x_1 + x_2 + 0e_1 + 0e_2 + e_3 & = & 20 \\
	x_1, x_2, e_1, e_2, e_3 & \geq & 0
\end{array}
\right.
\]

%	\subsection*{Remarque 3.2}
\begin{remarque}
	Cette présentation sous forme d'un tableau de nombres est appelée matrice du problème de programmation linéaire. Toutes les propriétés que vous connaissez sur les calculs à propos des matrices peuvent être alors utilisées.
\end{remarque}


\textbf{ÉTAPE no 1 :} Matrice du programme de base no 1 :

\begin{center}
	\begin{tabular}{c|cccccc|c|c}
		\toprule
		Variables de base & $x_{1}$ & $x_{2}$ & $e_{1}$ & $e_{2}$ & $e_{3}$ & $z$ & $b$&R \\
		\midrule
		$e_{1}$ & 3 & 9 & 1 & 0 & 0 & 0 & 81& \\
		$e_{2}$ & 4 & 5 & 0 & 1 & 0 & 0 & 55 &\\
		$e_{3}$ & 2 & 1 & 0 & 0 & 1 & 0 & 20 &\\
		\hline
		$z$ & 6 & 4 & 0 & 0 & 0 & -1 & 0 &\\
		\bottomrule
	\end{tabular}
\end{center}

\textbf{ÉTAPE no 2 :} Matrice du programme de base no 2 :


Pour augmenter le plus rapidement possible la valeur de la fonction économique, il faut donner une valeur positive à $ x_{1} $ puisque c'est qui a le plus grand coefficient positif sur la dernière ligne du tableau. Quelle est la plus grande valeur que l'on peut attribuer à $ x_{1} $ ?	Pour la déterminer, on prend les rapports des éléments de la colonne de b sur les éléments de la colonne des $ x_{1} $. On trouve alors :
\begin{center}
	\begin{tabular}{c|>{\columncolor{gray!25}}cccccc|c|c}
		\toprule
		Variables de base & \textcolor{blue}{$x_{1}$} & $x_{2}$ & $e_{1}$ & $e_{2}$ & $e_{3}$ & $z$ & $b$&R \\
		\midrule
		$e_{1}$ & 3 & 9 & 1 & 0 & 0 & 0 & 81 & $\dfrac{81}{3}=27$ \\
		$e_{2}$ & 4 & 5 & 0 & 1 & 0 & 0 & 55 & $\dfrac{55}{4}$ \\
		\rowcolor{gray!25}\textcolor{red}{$e_{3}$} & 2 & 1 & 0 & 0 & 1 & 0 & 20 & $\dfrac{20}{2}=10$ \\
		\hline
		$z$ & 6 & 4 & 0 & 0 & 0 & -1 & 0 & \\
		\bottomrule
	\end{tabular}
\end{center}

Le plus petit de ces rapports est $ 10 $ et c'est la plus grande valeur qu'on peut attribuer à $ x_{1} $. Nous utiliserons donc l'élément de la troisième ligne première colonne; encadrons-le pour le mettre en évidence. Nous l'appellerons pivot de la 1ère étape.
\begin{center}
	\begin{tabular}{c|>{\columncolor{gray!25}}cccccc|c|c}
		\toprule
		Variables de base & \textcolor{blue}{$x_{1}$} & $x_{2}$ & $e_{1}$ & $e_{2}$ & $e_{3}$ & $z$ & $b$ &R\\
		\midrule
		$e_{1}$ & 3 & 9 & 1 & 0 & 0 & 0 & 81 & $\dfrac{81}{3}=27$ \\
		$e_{2}$ & 4 & 5 & 0 & 1 & 0 & 0 & 55 & $\dfrac{55}{4}$ \\
		\rowcolor{gray!25}\textcolor{red}{$e_{3}$} & \tikz[baseline=(char.base)]{\node[draw=red,circle,very thick,inner sep=2pt] (char) {2};} & 1 & 0 & 0 & 1 & 0 & 20 & $\dfrac{20}{2}=10$ \\
		\hline
		$z$ & 6 & 4 & 0 & 0 & 0 & -1 & 0 & \\
		\bottomrule
	\end{tabular}
\end{center}

Nous devons maintenant annuler les autres éléments de la première colonne. Pour ce faire, nous effectuons des opérations sur les lignes. En premier lieu, nous allons rendre le pivot unitaire en divisant les termes de la troisième ligne par 2 :

\[
L_{3} \leftarrow \dfrac{1}{2} L_{3}
\]
\begin{center}
	\begin{tabular}{c|cccccc|c|c}
		\toprule
		Variables de base & $x_{1}$ & $x_{2}$ & $e_{1}$ & $e_{2}$ & $e_{3}$ & $z$ & $b$&R \\
		\midrule
		$e_{1}$ & 3 & 9 & 1 & 0 & 0 & 0 & 81& \\
		$e_{2}$ & 4 & 5 & 0 & 1 & 0 & 0 & 55 &\\
		$x_{1}$ & 1 & $\dfrac{1}{2}$ & 0 & 0 & $\dfrac{1}{2}$ & 0 & 10 &\\
		\hline
		$z$ & 6 & 4 & 0 & 0 & 0 & -1 & 0& \\
		\bottomrule
	\end{tabular}
\end{center}

En deuxième lieu, on soustrait 4 fois la troisième ligne à la deuxième, 3 fois la troisième à la
première et finalement 6 fois la troisième à la dernière :
\[
\begin{aligned}
	L_{1} &\leftarrow L_{1} - 3L_{3} \\
	L_{2} &\leftarrow L_{2} - 4L_{3} \\
	L_{4} &\leftarrow L_{4} - 6L_{3}
\end{aligned}
\]

\begin{center}
	\begin{tabular}{c|cccccc|c|c}
		\toprule
		Variables de base & $x_{1}$ & $x_{2}$ & $e_{1}$ & $e_{2}$ & $e_{3}$ & $z$ & $b$&R \\
		\midrule
		$e_{1}$ & 0 & $\dfrac{15}{2}$ & 1 & 0 & $-\dfrac{3}{2}$ & 0 & 51& \\
		$e_{2}$ & 0 & 3 & 0 & 1 & -2 & 0 & 15& \\
		$x_{1}$ & 1 & $\dfrac{1}{2}$ & 0 & 0 & $\dfrac{1}{2}$ & 0 & 10& \\
		\hline
		$z$ & 0 & 1 & 0 & 0 & -3 & -1 & -60& \\
		\bottomrule
	\end{tabular}
\end{center}
Ce tableau représente alors le système d'équations que nous avions obtenu après la deuxième étape de la résolution algébrique \eqref{equ2.7}. On constate également que la valeur obtenue en bas à droite du tableau correspond à l'opposé de la valeur numérique obtenue dans la fonction économique :
\[
z = 60 + x_{2} - 3e_{3}
\]
On constate que l'on peut encore augmenter la valeur de la fonction économique en introduisant une valeur pour $ x_{2} $.


\textbf{ÉTAPE no 3 :} Matrice du programme de base no 3 :

Nous avons
\begin{center}
	\begin{tabular}{c|c>{\columncolor{gray!25}}ccccc|c|c}
		\toprule
		Variables de base & $x_{1}$ & \textcolor{blue}{$x_{2}$} & $e_{1}$ & $e_{2}$ & $e_{3}$ & $z$ & $b$&R \\
		\midrule
		$e_{1}$ & 0 & $\dfrac{15}{2}$ & 1 & 0 & $-\dfrac{3}{2}$ & 0 & 51 & $\dfrac{51}{\frac{15}{2}}= \dfrac{102}{15}$ \\
		\rowcolor{gray!25}\textcolor{red}{$e_{2}$} & 0 & \tikz[baseline=(char.base)]{\node[draw=red,circle,very thick,inner sep=2pt] (char) {3};} & 0 & 1 & -2 & 0 & 15 & $\dfrac{15}{3}=5$ \\
		$x_{1}$ & 1 & $\dfrac{1}{2}$ & 0 & 0 & $\dfrac{1}{2}$ & 0 & 10 & $\dfrac{10}{\dfrac{1}{2}}=20$ \\
		\hline
		$z$ & 0 & 1 & 0 & 0 & -3 & -1 & -60 & \\
		\bottomrule
	\end{tabular}
\end{center}

Le plus petit rapport étant sur la deuxième ligne, notre pivot sera l'élément de la deuxième ligne, deuxième colonne. Rendons-le unitaire en le divisant par 3 :
\[
L_{2} \leftarrow \dfrac{1}{3} L_{2}
\]

\begin{center}
	\begin{tabular}{c|ccccc|c|c}
		\toprule
		Variables de base & $x_{1}$ & $x_{2}$ & $e_{1}$ & $e_{2}$ & $e_{3}$ & $z$ & $b$ \\
		\midrule
		$e_{1}$ & 0 & $\dfrac{15}{2}$ & 1 & 0 & $-\dfrac{3}{2}$ & 0 & 51 \\
		$x_{2}$ & 0 & 1 & 0 & $\dfrac{1}{3}$ & $-\dfrac{2}{3}$ & 0 & 5 \\
		$x_{1}$ & 1 & $\dfrac{1}{2}$ & 0 & 0 & $\dfrac{1}{2}$ & 0 & 10 \\
		\hline
		$z$ & 0 & 1 & 0 & 0 & -3 & -1 & -60 \\
		\bottomrule
	\end{tabular}
\end{center}

Effectuons à nouveau les opérations sur les lignes afin d'annuler les autres éléments de la 2ème colonne :
\[
\begin{aligned}
	L_{1} &\leftarrow L_{1} - \dfrac{15}{2} L_{2} \\
	L_{3} &\leftarrow L_{3} - \dfrac{1}{2} L_{2} \\
	L_{4} &\leftarrow L_{4} - 1 L_{2}
\end{aligned}
\]

\begin{center}
	\begin{tabular}{c|ccccc|c|c}
		\toprule
		Variables de base & $x_{1}$ & $x_{2}$ & $e_{1}$ & $e_{2}$ & $e_{3}$ & $z$ & $b$ \\
		\midrule
		$e_{1}$ & 0 & 0 & 1 & $-\dfrac{5}{2}$ & $\dfrac{7}{2}$ & 0 & $\dfrac{27}{2}$ \\
		$x_{2}$ & 0 & 1 & 0 & $\dfrac{1}{3}$ & $-\dfrac{2}{3}$ & 0 & 5 \\
		$x_{1}$ & 1 & 0 & 0 & $-\dfrac{1}{6}$ & $\dfrac{1}{3}$ & 0 & $\dfrac{15}{2}$ \\
		\hline
		$z$ & 0 & 0 & 0 & $-\dfrac{1}{3}$ & $-\dfrac{7}{3}$ & -1 & -65 \\
		\bottomrule
	\end{tabular}
\end{center}

Il n'y a plus de coefficients positifs sur la dernière ligne du tableau et il n'est donc plus possible d'augmenter la valeur de la fonction économique. La solution optimale est donc
\[
(x_{1}, x_{2}, e_{1}, e_{2}, e_{3})^{t} = \left( \dfrac{15}{2}, 5, \dfrac{27}{2}, 0, 0 \right)^{t}
\]
qui correspond à un bénéfice maximum de $ z = 65 $ produits en produisant $ \dfrac{15}{2} $ unités du produit $ P_{1} $ et 5 unités du produit $ P_{2} $.

%\subsection{Problème de minimisation}
%
%Nous avons vu dans que tout programme de minimisation peut se transformer en un problème de maximisation.
%Pour la résolution d'un problème de minimisation par l'algorithme du simplexe, nous devons adapter le critère d'optimalité. Ainsi pour que la solution de base soit optimale il faut que les coûts marginaux de toutes les variables hors base soient positifs ou nuls.
%Pour mieux comprendre ce cas, je vous propose l'exemple suivant :
%
%Considérons le programme de minimisation sous forme canonique suivant :
%\[
%\left\{
%\begin{array}{r}
%	\displaystyle\min -x_{1} - 2x_{2} \\
%	\text{s.t.} \quad -3x_{1} + 2x_{2} \leq 2 \\
%	-x_{1} + 2x_{2} \leq 4 \\
%	x_{1} + x_{2} \leq 5 \\
%	x_{1} \geq 0, \quad x_{2} \geq 0
%\end{array}
%\right.
%\]
%
%Après introduction des variables d'écart :
%\[
%z = -x_{1} - 2x_{2} + 0e_{1} + 0e_{2} + 0e_{3}
%\]
%Sous les contraintes :
%\[
%\left\{
%\begin{aligned}
%	-3x_{1} + 2x_{2} + e_{1} &= 2 \\
%	-x_{1} + 2x_{2} + e_{2} &= 4 \\
%	x_{1} + x_{2} + e_{3} &= 5 \\
%	x_{1} + 2x_{2} - z &= 0 \\
%	x_{1}, x_{2}, e_{1}, e_{2}, e_{3} &\geq 0
%\end{aligned}
%\right.
%\]
%
%\textbf{ÉTAPE no 1 :} Matrice du programme de base no 1 : On a donc le tableau du simplexe suivant
%où les variables $ e_{1} $, $ e_{2} $ et $ e_{3} $ jouent le rôle des variables de base :
%
%\begin{center}
%	\begin{tabular}{c|ccccc|c|c}
%		\toprule
%		Variables de base & $x_{1}$ & $x_{2}$ & $e_{1}$ & $e_{2}$ & $e_{3}$ & $-z$ & $b$ \\
%		\midrule
%		$e_{1}$ & -3 & 2 & 1 & 0 & 0 & 0 & 2 \\
%		$e_{2}$ & -1 & 2 & 0 & 1 & 0 & 0 & 4 \\
%		$e_{3}$ & 1 & 1 & 0 & 0 & 1 & 0 & 5 \\
%		\hline
%		$-z$ & 1 & 2 & 0 & 0 & 0 & 1 & 0 \\
%		\bottomrule
%	\end{tabular}
%\end{center}
%Cette base n'est pas optimale puisque les coûts relatifs des variables hors base dans la dernière
%ligne ne sont pas positifs ou nuls.
%\[\displaystyle\min\left\lbrace\dfrac{b_{i}}{a_{i2}}, \, a_{i2}>0\right\rbrace   \]
%\textbf{ÉTAPE no 2 :}  Matrice du programme de base no 2 :
%
%Choisissons comme variable entrante dans la nouvelle base la variable hors base ayant le plus petit coût négatif négatif, c'est à dire $ x_{2} $. cherchons la variable sortante en appliquant la règle
%du plus petit rapport :
%
%\begin{center}
%	\begin{tabular}{c|c|>{\columncolor{gray!25}}c|c|c|c|c|c|c}
%		\hline
%		Variables de base & $x_{1}$ & \textcolor{blue}{$x_{2}$} & $e_{1}$ & $e_{2}$ & $e_{3}$ & $-z$ & $b$ & R \\
%		\hline
%		\rowcolor{gray!25}
%		\textcolor{red}{$e_{1}$} & $-3$ & \tikz[baseline=(char.base)]{\node[draw=red,circle,very thick,inner sep=2pt] (char) {2};} & 1 & 0 & 0 & 0 & 2 & $\dfrac{2}{2}=1$ \\
%		\hline
%		$e_{2}$ & $-1$ & 2 & 0 & 1 & 0 & 0 & 4 & $\dfrac{4}{2}=2$ \\
%		\hline
%		$e_{3}$ & 1 & 1 & 0 & 0 & 1 & 0 & 5 & $\dfrac{5}{1}=5$ \\
%		\hline
%		$-z$ & 1 & 2 & 0 & 0 & 0 & 1 & 0 & \\
%		\hline
%	\end{tabular}
%\end{center}
%
%
%Le plus petit de ces rapports est 1 et c'est la plus petite valeur qu'on peut attribuer à $ x_{2} $. Nous utiliserons donc l'élément de la première ligne deuxième colonne, autrement, la variable $ e_{1} $ quitte la base et sera remplacée par la variable $ x_{2} $, ainsi on obtient la tableau suivant.
%\[
%L_{1} \leftarrow \dfrac{1}{2} L_{1}
%\]
%
%\begin{center}
%	\begin{tabular}{c|ccccc|c|c}
%		\toprule
%		Variables de base & $x_{1}$ & $x_{2}$ & $e_{1}$ & $e_{2}$ & $e_{3}$ & $-z$ & $b$ \\
%		\midrule
%		$x_{2}$ & $-\dfrac{3}{2}$ & 1 & $\dfrac{1}{2}$ & 0 & 0 & 0 & 1 \\
%		$e_{2}$ & -1 & 2 & 0 & 1 & 0 & 0 & 4 \\
%		$e_{3}$ & 1 & 1 & 0 & 0 & 1 & 0 & 5 \\
%		\hline
%		$-z$ & 1 & 2 & 0 & 0 & 0 & 1 & 0 \\
%		\bottomrule
%	\end{tabular}
%\end{center}
%
%Nous devons maintenant annuler les autres éléments de la deuxième colonne.
%\[
%\begin{aligned}
%	L_{2} &\leftarrow L_{2} - 2 L_{1} \\
%	L_{3} &\leftarrow L_{3} - 1 L_{1} \\
%	L_{4} &\leftarrow L_{4} - 2 L_{1}
%\end{aligned}
%\]
%
%\begin{center}
%	\begin{tabular}{c|ccccc|c|c}
%		\toprule
%		Variables de base & $x_{1}$ & $x_{2}$ & $e_{1}$ & $e_{2}$ & $e_{3}$ & $-z$ & $b$ \\
%		\midrule
%		$x_{2}$ & $-\dfrac{3}{2}$ & 1 & $\dfrac{1}{2}$ & 0 & 0 & 0 & 1 \\
%		$e_{2}$ & 2 & 0 & -1 & 1 & 0 & 0 & 2 \\
%		$e_{3}$ & $\dfrac{5}{2}$ & 0 & $-\dfrac{1}{2}$ & 0 & 1 & 0 & 4 \\
%		$-z$ & 4 & 0 & -1 & 0 & 0 & 1 & -2 \\
%		\bottomrule
%	\end{tabular}
%\end{center}
%%==================================================================================================================================================================================================
%La solution de base correspondante n'est pas optimale car $c_{1}=-4<0$. la variable $x_{1}$ va donc entrer dans la base.\\
%
%\textbf{ÉTAPE No 3 :} Matrice du programme de base no 3 :\\
%
%\begin{center}
%	\begin{tabular}{c|>{\columncolor{gray!25}}cccccc|c|c}
%		\hline
%		Variables de base & \textcolor{blue}{$x_{1}$} & $x_{2}$ & $e_{1}$ & $e_{2}$ & $e_{3}$ & $-z$ & $b$ & $R$ \\
%		\hline
%		$x_{2}$ & $\dfrac{-3}{2}$ & 1 & $\dfrac{1}{2}$ & 0 & 0 & 0 & 1 &  \\
%		\rowcolor{gray!25}\textcolor{red}{$e_{2}$} & \tikz[baseline=(char.base)]{\node[draw=red,circle,very thick,inner sep=2pt] (char) {2};} & 0 & -1 & 1 & 0 & 0 & 2 & $\dfrac{2}{2}=1$ \\
%		$e_{3}$ & $\dfrac{5}{2}$ & 0 & $-\dfrac{1}{2}$ & 0 & 1 & 0 & 4 & $\dfrac{4}{\dfrac{5}{2}}=\dfrac{8}{5}$ \\
%		\hline
%		$-z$ & -4 & 0 & 1 & 0 & 0 & 1 & 2 &  \\
%		\hline
%	\end{tabular}
%\end{center}
%
%Le plus petit rapport étant sur la deuxième ligne donc $e_{2}$ est une variable sortante, notre pivot sera l'élément de la deuxième ligne, première colonne. Rendons-le unitaire en le divisant par 2 :
%
%$L_{2} \leftarrow \dfrac{1}{2} L_{2} \quad$\begin{tabular}{c|cccccc|c|c}
%	\hline
%	Variables de base & $x_{1}$ & $x_{2}$ & $e_{1}$ & $e_{2}$ & $e_{3}$ & $-z$ & $b$ & $R$ \\
%	\hline
%	$x_{2}$ & $\dfrac{-3}{2}$ & 1 & $\dfrac{1}{2}$ & 0 & 0 & 0 & 1 &  \\
%	$x_{1}$ & 1 & 0 & $-\dfrac{1}{2}$ & $\dfrac{1}{2}$ & 0 & 0 & 1 &  \\
%	$e_{3}$ & $\dfrac{5}{2}$ & 0 & $-\dfrac{1}{2}$ & 0 & 1 & 0 & 4 &  \\
%	\hline
%	$-z$ & -4 & 0 & 1 & 0 & 0 & 1 & 2 &  \\
%	\hline
%\end{tabular}
%
%Ainsi,
%
%\begin{center}
%	\begin{tabular}{ll|llllll|l|l}
%		\cline{2-10}
%		\multirow{3}{*}{$L_{1} \leftarrow L_{1}+\dfrac{3}{2} L_{2}$} & Variables de base & $x_{1}$ & $x_{2}$ & $e_{1}$ & $e_{2}$ & $e_{3}$ & -z & $b$ & $R$ \\
%		\cline{2-10}
%		& \multirow{3}{*}{\begin{tabular}{l}
%				$x_{2}$ \\
%				$x_{1}$ \\
%				$e_{3}$ \\
%		\end{tabular}} & 0 & 1 & $-\dfrac{1}{4}$ & $\dfrac{3}{4}$ & 0 & 0 & $\dfrac{5}{2}$ & \multirow{3}{*}{} \\
%		&  & 1 & 0 & $-\dfrac{1}{2}$ & $\dfrac{1}{2}$ & 0 & 0 & 1 &  \\
%		$L_{3} \leftarrow L_{3}-\dfrac{5}{2} L_{2}$ &  & 0 & 0 & $\dfrac{3}{4}$ & $-\dfrac{5}{4}$ & 1 & 0 & $\dfrac{3}{2}$ &  \\
%		\cline{2-10}
%		$L_{4} \leftarrow L_{4}+4 L_{2}$ & $-z$ & 0 & 0 & -1 & 2 & 0 & 1 & 6 &  \\
%		\cline{2-10}
%	\end{tabular}
%\end{center}
%
%La solution de base obtenue a pour coût $z=-6$, cependant elle n'est pas optimale car $c_{3}=-1<0$, la variable $e_{1}$ va donc entrer dans la base.\\
%
%\textbf{ÉTAPE No 4 :} Matrice du programme de base no 4 :\\
%
%\begin{center}
%	\begin{tabular}{c|cc>{\columncolor{gray!25}}cccc|c|c}
%		\toprule
%		Variables de base & $x_{1}$ & $x_{2}$ & \textcolor{blue}{$e_{1}$} & $e_{2}$ & $e_{3}$ & $-z$ & $b$ & $R$ \\
%		\midrule
%		$x_{2}$ & 0 & 1 & $-\dfrac{1}{4}$ & $\dfrac{3}{4}$ & 0 & 0 & $\dfrac{5}{2}$ &  \\
%		$x_{1}$ & 1 & 0 & $-\dfrac{1}{2}$ & $\dfrac{1}{2}$ & 0 & 0 & 1 &  \\
%		\rowcolor{gray!25}\textcolor{red}{$e_{3}$} & 0 & 0 & \tikz[baseline=(char.base)]{\node[draw=red,circle,very thick,inner sep=2pt] (char) {$ \dfrac{3}{4} $};} & $-\dfrac{5}{4}$ & 1 & 0 & $\dfrac{3}{2}$ & $\dfrac{\dfrac{3}{2}}{\dfrac{3}{4}}=2$ \\
%		\hline
%		$-z$ & 0 & 0 & -1 & 2 & 0 & 1 & 6 &  \\
%		\hline
%	\end{tabular}
%\end{center}
%
%Par suite, on obtient
%
%\begin{tabular}{cc|cccccc|c|c}
%	\cline { 2 - 10 }
%	&Variables de base & $x_{1}$ & $x_{2}$ & $e_{1}$ & $e_{2}$ & $e_{3}$ & $-z$ & $b$ & $R$ \\
%	\cline { 2 - 10 }
%	& $x_{2}$&0 & 1 & $-\dfrac{1}{4}$ & $\dfrac{3}{4}$ & 0 & 0 & $\dfrac{5}{2}$& \\
%	& $x_{1}$& 1 & 0 & $-\dfrac{1}{2}$ & $\dfrac{1}{2}$ & 0 & 0 & 1& \\
%	$L_{3} \leftarrow \dfrac{4}{3} L_{3}$&$e_{1}$ & 0 & 0 & 1 & $-\dfrac{5}{3}$ & $\dfrac{4}{3}$ & 0 & 2 &  \\
%	\cline { 2 - 10 }
%	&$-z$&0 & 0 & -1 & 2 & 0 & 1 & 6& \\
%	\cline { 2 - 10 }
%\end{tabular}
%
%rendons la colonne associé a $e_{1}$ unitaire
%
%\begin{center}
%	\begin{tabular}{cc|cccccc|c|c}
%		\cline { 2 - 10 }
%		& Variables de base & $x_{1}$ & $x_{2}$ & $e_{1}$ & $e_{2}$ & $e_{3}$ & $-z$ & $b$ & $R$ \\
%		\cline { 2 - 10 }
%		$L_{1} \leftarrow L_{1}+\dfrac{1}{4} L_{3}$& $x_{2}$ & 0 & 1 & 0 & $\dfrac{1}{3}$ & $\dfrac{1}{3}$ & 0 & 3 &  \\
%		$L_{2} \leftarrow L_{2}+\dfrac{1}{2} L_{3}$ & $x_{1}$ & 1 & 0 & 0 & $-\dfrac{1}{3}$ & $\dfrac{2}{3}$ & 0 & 2 &  \\
%		& $e_{1}$ & 0 & 0 & 1 & $-\dfrac{5}{3}$ & $\dfrac{4}{3}$ & 0 & 2 &  \\
%		\cline { 2 - 10 }
%		$L_{4} \leftarrow L_{4}+L_{3}$& $-z$ & 0 & 0 & 0 & $\dfrac{1}{3}$ & $\dfrac{4}{3}$ & 1 & 8 &  \\
%		\cline { 2 - 10 }
%	\end{tabular}
%\end{center}
%
%La nouvelle solution de base est
%
%$$
%\left(x_{1}, x_{2}, e_{1}, e_{2}, e_{3}\right)^{t}=(2,3,2,0,0)^{t}
%$$
%
%qui correspond à un coût minimum de $z=-8$, comme les coût réduit de toutes les variables sont positifs ou nuls, cette solution est optimale et l'algorithme du simplexe s'arrête.

\section{Principes de la Méthode du Rectangle}

Pour faciliter les calculs dans les tableaux du simplexe, on peut utiliser la méthode du rectangle :

\begin{enumerate}
	\item Si une ligne contient un zéro (0) qui est à l'intersection avec la colonne pivot, cette ligne ne sera pas modifiée.
	
	\item Si une colonne contient un zéro (0) qui est à l'intersection avec la ligne pivot, cette colonne ne sera pas modifiée.
	
	\item Pour les cellules qui ne se trouvent pas sur la ligne et la colonne pivot, appliquer la méthode du rectangle :
	$$a' \leftarrow a - \frac{b \times c}{\text{pivot}}$$
	où $a'$ est la nouvelle valeur de $a$, $b$ l'élément situé sur la même ligne que $a$ mais dans la colonne du pivot, $c$ l'élément situé sur la même colonne que $a$ mais dans la ligne du pivot.
	
	\item Pour les éléments situés sur la ligne du pivot, ils seront divisés par le pivot.
	
	\item Pour les éléments situés sur la colonne du pivot, ceux situés en dessous et au dessus du pivot seront réduits à zéro (0).
\end{enumerate}

%\subsection{Conclusion}
%L'algorithme du simplexe consiste à se déplacer d'un point extrême du polytope $K=\{x: A x= b, x \geq 0\}$ à un autre point extrême adjacent jusqu'à ce qu'on se rende compte que la solution réalisable actuelle est optimale ou que le problème n'est pas réalisable ou que le problème n'est pas borné supérieurement (ou inférieurement dans le cas de minimisation).

\section{Méthodes des deux phases}
Nous avons vu comment trouver une base de départ lorsque les contraintes du problème sont toutes de type $\leq$ : les variables d'écart constituent alors une base de départ naturelle, leurs colonnes formant a une permutation près une matrice identité.

En pratique il est fréquent que le PL comporte également des contraintes de type $\geq$ ou de type $=$. Dans ce cas, nous n'avons plus de base de départ évidente, il faut en chercher une. Ce sera la première phase de la méthode du simplexe.

La méthode du simplexe consiste donc en deux phases :\\

\begin{description}
	\item[Phase I:] recherche d'une base de départ par la résolution par l'algorithme du simplexe du programme auxiliaire (définit plus loin) ou mise en évidence que le problème de départ n'admet aucune solution admissible.
	
	\item[Phase II :] application de l'algorithme du simplexe au problème de départ cette fois, en partant de la base mise en évidence a la fin de la phase I
\end{description}

%Exemple 3.1 
\begin{exemple}
	Comment, par exemple trouver une solution de base réalisable initiale au problème suivant
	
	$$
	(P):\left\{\begin{aligned}
		\displaystyle\max_{x} \quad 2 x_{1}+3 x_{2} & \\
		\text { s.t. } \quad x_{1}+x_{2} & \leq 10, \\
		5 x_{1}+4 x_{2} & \geq 20 \\
		x_{1}, x_{2} & \geq 0 .
	\end{aligned}\right.
	$$
	
	On transforme ce problème sous la forme canonique $A x \leq b$ et ensuite on ajoute les variables d'écart.
	
	\begin{equation}\label{eq3.3}
		\left\{\begin{aligned}
			\max_{x} \quad 2 x_{1}+3 x_{2} & \\
			\text { s.t. } \quad x_{1}+x_{2}+e_{1} & =10, \\
			-5 x_{1}-4 x_{2}+e_{2} & =-20 \\
			x_{1}, x_{2}, e_{1}, e_{2} & \geq 0 .
		\end{aligned}\right.
	\end{equation}
	
	Clairement le point ( $0,0,10,-20$ ), n'est pas admissible.
\end{exemple}

\subsection{Recherche d'une base de départ : Phase I}
L'objectif est de trouver une solution de base admissible qui servira de point de départ pour l'algorithme du simplexe. L'idée est de résoudre un problème intermédiaire de minimisation dont la solution fournira le point de départ de la méthode du simplexe. Ce problème intermédiaire porte le nom de Phase I du simplexe.


\begin{colorframed}
	{\color{titleblue}\textbf{Démarche à suivre :}}
	
	\medskip
	
	On commence par écrire le {\color{keywordred}programme auxiliaire (PA)} associé au {\color{keywordred}problème de départ (P)} :
	
	\begin{enumerate}[label=\arabic*.]
		\item On s'assure que les {\color{vargreen}membres de droite $b_i$} sont tous {\color{keywordred}positifs ou nuls}, quitte à multiplier la contrainte correspondante par {\color{mathblue}$-1$} si son membre de droite est {\color{keywordred}négatif} (sinon on aboutirait à des {\color{keywordred}solutions de base non réalisables}).
		
		\item On introduit une {\color{vargreen}variable d'écart $e_i$} pour chaque contrainte $i$ de type {\color{mathblue}$\leq$}, affectée d’un signe {\color{mathblue}$+$} dans la contrainte réécrite, comme on le faisait jusqu'à présent.
		
		\item On introduit aussi une {\color{vargreen}variable d'écart $e_i$} pour chaque contrainte $i$ de type {\color{mathblue}$\geq$}, mais affectée d’un signe {\color{mathblue}$-$} dans la contrainte réécrite : l’écart qu’elle mesure est en effet {\color{keywordred}négatif ou nul}, mais la variable d’écart doit rester {\color{keywordred}positive ou nulle}.
		
		\item Le problème ne comporte plus que des contraintes d’{\color{keywordred}égalité}, mais seules les {\color{vargreen}variables d’écart associées aux contraintes initiales de type {\color{mathblue}$\leq$}} forment des colonnes de la matrice identité. Pour chacune des autres contraintes (c’est-à-dire les contraintes initiales {\color{mathblue}$\geq$} réécrites avec une variable d’écart affectée d’un signe {\color{mathblue}$-$}, et les contraintes initiales {\color{mathblue}$=$}), on introduit une {\color{vargreen}variable artificielle $t_i$} {\color{keywordred}positive ou nulle} de la même façon qu’on introduisait les variables d’écart.
		
		\item La {\color{keywordred}fonction objectif de (PA)} est alors la {\color{vargreen}somme des variables artificielles}. On cherche à la {\color{keywordred}minimiser} (en fait à l’{\color{keywordred}annuler}).
	\end{enumerate}
\end{colorframed}

\subsection{Illustration de la méthode sur un exemple}
Dans l'exemple ci-dessus, il s'agit d'introduire deux variables d'écart $e_{1}$ et $e_{2}$ et une variable artificielle $t_{1}$ et de considérer le problème de minimisation associé
\[(P): \left\{
\begin{array}{ll}
	\underset{x}{\displaystyle\max} & 2x_1 + 3x_2 \\
	\text{s.t.} & x_1 + x_2 \leq 10, \\
	& 5x_1 + 4x_2 \geq 20 \\
	& x_1, x_2 \geq 0.
\end{array}
\right.
\overset{\text{Var. d'écart}}{\Longrightarrow}
\begin{array}{c}
	\left\{
	\begin{array}{ll}
		\underset{x}{\displaystyle\max} & 2x_1 + 3x_2 \\
		\text{s.t.} & x_1 + x_2 + e_1 = 10, \\
		& 5x_1 + 4x_2 - e_2 = 20 \\
		& x_1, x_2, e_1, e_2 \geq 0.
	\end{array}
	\right.
	\\
	\\
	\Big\Downarrow \substack{\textcolor{blue}{\text{Var. artificiel}} \\ \text{+ new obj}}
	\\
	\\
	(PA)\left\{
	\begin{array}{l}
		\displaystyle\min \quad z_t = \textcolor{red}{t_1(20-5x_1-4x_2+e_2)} \\
		\begin{array}{ll}
			\text{s.t.} & x_1 + x_2 + e_1 = 10, \\
			& 5x_1 + 4x_2 - e_2 + t_1 = 20 \\
			& x_1, x_2, e_1, e_2, t_1 \geq 0.
		\end{array}
	\end{array}
	\right.
\end{array}\]


\begin{description}
	\item[Phase I]\quad\\
	Le problème (PA) admet toujours une solution admissible, il suffit de prendre le point ( $x_{1}, x_{2}, e_{1}, e_{2}, t_{1}$ ) = ( $0,0,10,0,20$ ).
	
	On forme le tableau initial du simplexe pour ce problème.
	
	\begin{center}
		\begin{tabular}{|c|cccccc|c|c|}
			\hline
			Variables de base & $x_{1}$ & $x_{2}$ & $e_{1}$ & $e_{2}$ & $t_{1}$ & $-z_{t}$ & $b$ & $R$ \\
			\hline
			$e_{1}$ & 1 & 1 & 1 & 0 & 0 & 0 & 10 &  \\
			$t_{1}$ & 5 & 4 & 0 & -1 & 1 & 0 & 20 &  \\
			\hline
			$-z_{t}$ & 0 & 0 & 0 & 0 & 1 & 1 & 0 &  \\
			\hline
		\end{tabular}
	\end{center}
	
	Pour que $t_{1}$ joue le rôle d'une variable de base, il faut faire apparaître un 0 sur la dernière ligne de la colonne de $t_{1}$, pas suite, on transforme le tableau en soustrayant la deuxième ligne de la dernière. Nous obtenons le tableau suivant :
	
	\begin{center}
		\begin{tabular}{|c|>{\columncolor{gray!25}}cccccc|c|c|}
			\hline
			Variables de base & \textcolor{blue}{$x_{1}$} & $x_{2}$ & $e_{1}$ & $e_{2}$ & $t_{1}$ & $-z_{t}$ & $b$ & $R$ \\
			\hline
			$e_{1}$ & 1 & 1 & 1 & 0 & 0 & 0 & 10 & $\dfrac{10}{1}$ \\
			\rowcolor{gray!25}\textcolor{red}{$t_{1}$} & \tikz[baseline=(char.base)]{\node[draw=red,circle,very thick,inner sep=2pt] (char) {5};} & 4 & 0 & -1 & 1 & 0 & 20 & $\dfrac{20}{5}=4$ \\
			\hline
			$-z_{t}$ & -5 & -4 & 0 & 1 & 0 & 1 & -20 &  \\
			\hline
		\end{tabular}
	\end{center}
	
	Cette solution de base n'étant pas optimale, on poursuit avec la méthode standard du simplexe. Dans notre cas, on choisit la variable entrante $x_{1}$ et la variable sortante $t_{1}$, on obtient le tableau suivant
	
	\begin{center}
		\begin{tabular}{|c|cccccc|c|c|}
			\hline
			Variables de base & $x_{1}$ & $x_{2}$ & $e_{1}$ & $e_{2}$ & $t_{1}$ & $-z_{t}$ & $b$ & $R$ \\
			\hline
			$e_{1}$ & 0 & $\dfrac{1}{5}$ & 1 & $\dfrac{1}{5}$ & $-\dfrac{1}{5}$ & 0 & 6 &  \\
			$x_{1}$ & 1 & $\dfrac{4}{5}$ & 0 & $-\dfrac{1}{5}$ & $\dfrac{1}{5}$ & 0 & 4 &  \\
			\hline
			$-z_{t}$ & 0 & 0 & 0 & 0 & 1 & 1 & 0 &  \\
			\hline
		\end{tabular}
	\end{center}
	
	On ne peut améliorer ce résultat car la fonction objective $z=t_{1}=0$.\\
	Nous avons obtenu la solution optimale de la Phase $I(4,0,6,0,0)$ avec la base $\left\{e_{1}, x_{1}\right\}$. De plus, la variable $t_{1}=0$ devient inutile car hors-base. Ainsi, on peut enlever la colonne correspondante à $t_{1}$.
	
	Donc $(4,0,6,0,0)$ constitue une solution de base initiale du problème $(P)$, ce qui va nous permettre de construire le tableau du simplexe de (P) pour démarrer la phase II.
	
	\item[Phase II]\quad\\
	Après avoir terminer la phase I avec une solution optimale de valeur optimale nulle et où les variables artificielles sont toutes des variables hors base. Pour commencer la première itération de la phase II, il faudra remplacer les coefficients de la fonction objectif $z_{t}$ par ceux de la fonction objectif original $\mathcal{z}$, puis recalculer les coûts relatifs.
	
	\begin{center}
		\begin{tabular}{|c|ccccc|c|c|}
			\hline
			Variables de base & $x_{1}$ & $x_{2}$ & $e_{1}$ & $e_{2}$ & $z$ & $b$ & $R$ \\
			\hline
			$e_{1}$ & 0 & $\dfrac{1}{5}$ & 1 & $\dfrac{1}{5}$ & 0 & 6 &  \\
			$x_{1}$ & 1 & $\dfrac{4}{5}$ & 0 & $-\dfrac{1}{5}$ & 0 & 4 &  \\
			\hline
			$z$ & 2 & 3 & 0 & 0 & -1 & 0 &  \\
			\hline
		\end{tabular}
	\end{center}
	
	$x_{1}$ étant une variable de base nous devons remplacer 2 par 0 dans la dernière ligne du tableau par l'opération suivante $L_{3} \leftarrow L_{3}-2 L_{2}$
	
	\begin{center}
		\begin{tabular}{|c|c>{\columncolor{gray!25}}cccc|c|c|}
			\hline
			Variables de base & $x_{1}$ & \textcolor{blue}{$x_{2}$} & $e_{1}$ & $e_{2}$ & $z$ & $b$ & $R$ \\
			\hline
			$e_{1}$ & 0 & $\dfrac{1}{5}$ & 1 & $\dfrac{1}{5}$ & 0 & 6 & $6 \times 5=30$ \\
			\rowcolor{gray!25}\textcolor{red}{$x_{1}$} & 1 & \tikz[baseline=(char.base)]{\node[draw=red,circle,very thick,inner sep=2pt] (char) {$ \dfrac{4}{5} $};} & 0 & $-\dfrac{1}{5}$ & 0 & 4 & $4 \times \dfrac{5}{4}=5$ \\
			\hline
			$z$ & 0 & $\dfrac{7}{5}$ & 0 & $\dfrac{2}{5}$ & -1 & -8 &  \\
			\hline
		\end{tabular}
	\end{center}
	
	La variable $x_{2}$ entre dans la base et la variable $x_{1}$ sort de la base.
	
	\begin{center}
		\begin{tabular}{|c|c>{\columncolor{gray!25}}cccc|c|c|}
			\hline
			Variables de base & $x_{1}$ & \textcolor{blue}{$x_{2}$} & $e_{1}$ & $e_{2}$ & $z$ & $b$ & $R$ \\
			\hline
			$e_{1}$ & 0 & $\dfrac{1}{5}$ & 1 & $\dfrac{1}{5}$ & 0 & 6 & $6 \times 5=30$ \\
			\rowcolor{gray!25}\textcolor{red}{$x_{1}$} & $\dfrac{5}{4}$ & 1 & 0 & $-\dfrac{1}{5}$ & 0 & 4 & $4 \times \dfrac{5}{4}=5$ \\
			\hline
			$z$ & 0 & $\dfrac{7}{5}$ & 0 & $\dfrac{2}{5}$ & -1 & -8 &  \\
			\hline
		\end{tabular}
	\end{center}
	
	On poursuit avec la méthode standard du simplexe
	
	\begin{center}
		\begin{tabular}{|c|rrr>{\columncolor{gray!25}}rr|c|c|}
			\hline
			Variables de base & $x_{1}$ & $x_{2}$ & $e_{1}$ & \textcolor{blue}{$e_{2}$} & $z$ & $b$ & $R$ \\
			\hline
			\rowcolor{gray!25}\textcolor{red}{$e_{1}$} & $-\dfrac{1}{4}$ & 0 & 1 & $\dfrac{1}{4}$ & 0 & 5 & $5 \times 4=20$ \\
			$x_{2}$ & $\dfrac{5}{4}$ & 1 & 0 & $-\dfrac{1}{4}$ & 0 & 5 &  \\
			\hline
			$z$ & $-\dfrac{7}{5}$ & 0 & 0 & $\dfrac{3}{4}$ & -1 & -15 &  \\
			\hline
		\end{tabular}
	\end{center}
	
	La variable $e_{2}$ entre dans la base et la variable $e_{1}$ sort de la base.
	
	\begin{center}
		\begin{tabular}{|c|ccccc|c|c|}
			\hline
			Variables de base & $x_{1}$ & $x_{2}$ & $e_{1}$ & $e_{2}$ & $z$ & $b$ & $R$ \\
			\hline
			$e_{2}$ & -1 & 0 & 4 & 1 & 0 & 20 &  \\
			$x_{2}$ & 1 & 1 & 1 & 0 & 0 & 10 &  \\
			\hline
			$z$ & -1 & 0 & -3 & 0 & -1 & -30 &  \\
			\hline
		\end{tabular}
	\end{center}
	
	L'algorithme se termine à cette étape. La solution optimale est ( $0,10,0,20$ ) avec $z=30$. Ceci est bien conforme au graphique ci-dessous.
\end{description}

%\section{La méthode de "The big $M$ "}
%L'idée ici est de remplacer le programme $ ( P ) $ par un autre qui a la même solution. Pratiquement il s'agit de faire deux choses :
%
%\begin{figure}[h]
%	\begin{center}
%		\includegraphics[scale=0.5]{2025_10_05_84c89f9c4234670b6556g-06}
%		%		\captionsetup{labelformat=empty}
%		\caption{Le graphe associé à l'exemple \ref{eq3.3}}
%	\end{center}
%\end{figure}
%
%\begin{itemize}
%	\item créer la base artificiellement en ajoutant les colonnes manquantes de la matrice identité, cela revient à introduire des variables supplémentaires, appelées "variables artificielles",
%	\item modifier la fonction objective de sorte que les variables artificielles quittent la base au fur et à mesure des itérations.\\
%	
%\end{itemize}
%
%Pour le méthode du grand $M$ on pénalise les variables artificielles dans la fonction objective avec un coefficient négatif très grand noté $-M$ pour un problème de maximisation ( $+M$ pour un problème de minimisation). Ainsi, la maximisation de la fonction objective passe nécessairement par forcer les variables artificielles à quitter la base. Si à l'optimum, certaines variables artificielles sont dans la base, avec une valeur non nulle, alors le programme n'à pas de solution réalisable.\\
%
%Pour l'exemple donné plus haut, le programme associé, noté $M$, est le suivant.
%
%$$
%(P A)\left\{\begin{aligned}
%	\displaystyle\max \quad z=2 x_{1}+3 x_{2}-M t_{1} & \\
%	\text { s.t. } \quad x_{1}+x_{2}+e_{1} & =10, \\
%	5 x_{1}+4 x_{2}-e_{2}+t_{1} & =20 \\
%	x_{1}, x_{2}, e_{1}, e_{2}, t_{1} & \geq 0 .
%\end{aligned}\right.
%$$
%
%Le tableau initial du simplexe est :
%
%\begin{center}
%	\begin{tabular}{|c|cccccc|c|c|}
%		\hline
%		Variables de base & $x_{1}$ & $x_{2}$ & $e_{1}$ & $e_{2}$ & $t_{1}$ & $z$ & $b$ & $R$ \\
%		\hline
%		$e_{1}$ & 1 & 1 & 1 & 0 & 0 & 0 & 10 &  \\
%		$t_{1}$ & 5 & 4 & 0 & -1 & 1 & 0 & 20 &  \\
%		\hline
%		$z$ & 2 & 3 & 0 & 0 & $-M$ & -1 & 0 &  \\
%		\hline
%	\end{tabular}
%\end{center}
%
%Pour que $t_{1}$ joue le rôle d'une variable de base, il faut faire apparaître un 0 sur la dernière ligne de la colonne de $t_{1}$, pas suite, on transforme le tableau en soustrayant $M$ fois la deuxième ligne de la dernière. Nous obtenons le tableau suivant:
%
%\begin{center}
%	\begin{tabular}{|c|>{\columncolor{gray!25}}cccccc|c|c|}
%		\hline
%		Variables de base & \textcolor{blue}{$x_{1}$} & $x_{2}$ & $e_{1}$ & $e_{2}$ & $t_{1}$ & $z$ & $b$ & $R$ \\
%		\hline
%		$e_{1}$ & 1 & 1 & 1 & 0 & 0 & 0 & 10 & $\dfrac{10}{1}$ \\
%		\rowcolor{gray!25}\textcolor{red}{$t_{1}$} & \tikz[baseline=(char.base)]{\node[draw=red,circle,very thick,inner sep=2pt] (char) {5};} & 4 & 0 & -1 & 1 & 0 & 20 & $\dfrac{20}{5}=4$ \\
%		\hline
%		$z$ & $2+5 M$ & $3+4 M$ & 0 & $-M$ & 0 & -1 & $20 M$ &  \\
%		\hline
%	\end{tabular}
%\end{center}
%
%La variable artificiel $t_{1}$ est toujours dans la base, nous devons effectuer une autre itération
%
%\begin{center}
%	\begin{tabular}{|c|cccccc|c|c|}
%		\hline
%		Variables de base & $x_{1}$ & $x_{2}$ & $e_{1}$ & $e_{2}$ & $t_{1}$ & $z$ & $b$ & $R$ \\
%		\hline
%		$e_{1}$ & 0 & $\dfrac{1}{5}$ & 1 & $\dfrac{1}{5}$ & $-\dfrac{1}{5}$ & 0 & 6 &  \\
%		$x_{1}$ & 1 & $\dfrac{4}{5}$ & 0 & $-\dfrac{1}{5}$ & $\dfrac{1}{5}$ & 0 & 4 &  \\
%		\hline
%		$z$ & 0 & $\dfrac{7}{5}$ & 0 & $\dfrac{2}{5}$ & $-2-5 M$ & -1 & -8 &  \\
%		\hline
%	\end{tabular}
%\end{center}
%
%la variable artificiel $t_{1}$ est hors base avec un fort coefficient négatif et ne peut plus entrer dans la base. La phase I s'achève en éliminant la colonne de la variable $t_{1}$ du tableau. Nous retrouvons ainsi le tableau suivant avec lequel on va démarrer la phase II du simplexe comme on l'a déjà fait dans la section précédente.
%
%%Remarque 3.3 - 
%\begin{remarque}
%	Il se peut que la phase I se termine avec une solution de base optimale comportant une ou plusieurs variables artificielles.
%	
%	\begin{itemize}
%		\item Si l'une de ces variables artificielles de base prend une valeur négative, alors le problème initial est non réalisable.
%		\item Si toutes les variables artificielles de base ont une valeur nulle, il faut les sortir de la base pour trouver le tableau initial du simplexe pour le problème $ ( P ) $. Pour cela, on doit remplacer chacune de ces variables artificielles par l'une des variables hors base, qui ne doit pas être une variable artificielle. Puisque la variable de sortie a une valeur nulle, alors la nouvelle variable de base aura aussi une valeur nulle.
%	\end{itemize}
%\end{remarque}
%
%\section{Problèmes irréguliers}
%\subsection{Cas de solution non bornée}
%Graphiquement, ce problème est caractérisé par le fait qu'on peut déplacer la droite de la fonction objectif indéfiniment de manière à accroître la valeur, en gardant toujours une intersection non vide avec l'ensemble des solutions réalisables.
%
%Avec la méthode de simplexe, on reconnaît ce problème lorsque la variable entrante n'admet aucune limite sur sa valeur d'entrée, c'est à dire que tous les ratios $\dfrac{b_{i}}{a_{i j}}$ sont négatifs ou nuls.
%
%On considère le problème suivant :
%
%\[
%\left\{
%\begin{aligned}
%	& \displaystyle\min \; -2x_1 - 3x_2 + x_3 \\
%	& -x_1 - x_2 - x_3 \leq 3 \\
%	& x_1 - x_2 + x_3 \leq 4 \\
%	& -x_1 + x_2 + 2x_3 \leq 1 \\
%	& x_1, x_2, x_3 \geq 0
%\end{aligned}
%\right.
%\overset{\text{rajout des variables d'écart}}{\Longrightarrow}
%\left\{
%\begin{aligned}
%	& \displaystyle\min \; z \\
%	& -x_1 - x_2 - x_3 + e_1 = 3 \\
%	& x_1 - x_2 + x_3 + e_2 = 4 \\
%	& -x_1 + x_2 + 2x_3 + e_3 = 1 \\
%	& -2x_1 - 3x_2 + x_3 - z = 0 \\
%	& x_1, x_2, x_3, e_1, e_2, e_3 \geq 0
%\end{aligned}
%\right.
%\]
%
%Le tableau initial s'écrit
%
%\begin{center}
%	\begin{tabular}{|c|c>{\columncolor{gray!25}}cccccc|c|c|}
%		\hline
%		Variables de base & $x_{1}$ & \textcolor{blue}{$x_{2}$} & $x_{3}$ & $e_{1}$ & $e_{2}$ & $e_{3}$ & $-z$ & $b$ & $R$ \\
%		\hline
%		$e_{1}$ & -1 & -1 & -1 & 1 & 0 & 0 & 0 & 3 &  \\
%		$e_{2}$ & 1 & -1 & 1 & 0 & 1 & 0 & 0 & 4 &  \\
%		\rowcolor{gray!25}\textcolor{red}{$e_{3}$} & -1 & \tikz[baseline=(char.base)]{\node[draw=red,circle,very thick,inner sep=2pt] (char) {1};} & 2 & 0 & 0 & 1 & 0 & 1 &  \\
%		\hline
%		$-z$ & -2 & -3 & 1 & 0 & 0 & 3 & 1 & 0 &  \\
%		\hline
%	\end{tabular}
%\end{center}
%
%Variable sortante : $e_{3}$.\\
%Variable entrante : $x_{2}$.
%
%On pivote autour de de la troisième ligne et on trouve le tableau suivant
%
%\begin{center}
%	\begin{tabular}{|c|>{\columncolor{gray!25}}ccccccc|c|c|}
%		\toprule
%		Variables de base & \textcolor{blue}{$x_{1}$} & $x_{2}$ & $x_{3}$ & $e_{1}$ & $e_{2}$ & $e_{3}$ & $-z$ & $b$ & $R$ \\
%		\midrule
%		$e_{1}$ & -2 & 0 & 1 & 1 & 0 & 0 & 0 & 4 &  \\
%		$e_{2}$ & 0 & 0 & 3 & 0 & 1 & 0 & 0 & 5 &  \\
%		$e_{3}$ & -1 & 1 & 2 & 0 & 0 & 1 & 0 & 1 &  \\
%		\hline
%		$-z$ & -5 & 0 & 7 & 0 & 0 & 0 & 1 & 3 &  \\
%		\hline
%	\end{tabular}
%\end{center}
%
%La variable entrante est $x_{1}$. Toutefois, le critère du quotient ne s'applique plus car toutes les entrées (lignes 1 à 3 ) de la colonne 1 sont négatives ou nulles.
%
%%Remarque 3.4 
%\begin{remarque}
%	Selon cet exemple, on peut conclure que s'il n'est pas possible de calculer la ligne de pivot pour un choix de colonne de pivot $j$, le problème n'admet pas de solution optimale car il sera toujours non borné.
%\end{remarque}
%
%\section{Les problèmes impossibles}
%Graphiquement, on a caractérisé ces problèmes par un ensemble de solutions réalisables vide. Avec la méthode de simplexe, on reconnaît que le problème est impossible si une ou plusieurs variables artificielles sont présentes dans la base dans le tableau de simplexe optimal, ce qui signifie que la solution donnée par ce tableau n'est pas réellement réalisable.
%
%Vérifions à l'aide de la méthode de simplexe, que le problème suivant est réellement impossible:
%
%$$
%\left\{\begin{array}{c}
%	\displaystyle\max \quad 4 x_{1}+3 x_{2} \\
%	x_{1}+x_{2} \leq 2 \\
%	3 x_{1}+x_{2} \geq 10 \\
%	x_{1}, x_{2} \geq 0
%\end{array}\right.
%$$
%
%En introduisant les variables d'écarts et les variables artificielles le programme s'écrit :
%
%$$
%\left\{\begin{array}{c}
%	\displaystyle\max \quad 4 x_{1}+3 x_{2}-M t \\
%	x_{1}+x_{2}+e_{1}=2 \\
%	3 x_{1}+x_{2}-e_{2}+t=10 \\
%	x_{1}, x_{2}, e_{1}, e_{2}, t \geq 0
%\end{array}\right.
%$$
%
%La résolution par le simplexe donne les tableaux suivant
%
%\begin{center}
%	\begin{tabular}{|c|cccccc|c|c|}
%		\hline
%		Variables de base & $x_{1}$ & $x_{2}$ & $e_{1}$ & $e_{2}$ & $t_{1}$ & $z$ & $b$ & $R$ \\
%		\hline
%		$e_{1}$ & 1 & 1 & 1 & 0 & 0 & 0 & 2 &  \\
%		$t_{1}$ & 3 & 1 & 0 & -1 & 1 & 0 & 10 &  \\
%		\hline
%		$z$ & 4 & 3 & 0 & 0 & $-M$ & -1 & 0 &  \\
%		\hline
%	\end{tabular}
%\end{center}
%
%Ainsi,
%
%\begin{center}
%	\begin{tabular}{|c|cccccc|c|c|}
%		\hline
%		Variables de base & $x_{1}$ & $x_{2}$ & $e_{1}$ & $e_{2}$ & $t_{1}$ & $z$ & $b$ & $R$ \\
%		\hline
%		$e_{1}$ & 1 & 1 & 1 & 0 & 0 & 0 & 2 &  \\
%		$t_{1}$ & 3 & 1 & 0 & -1 & 1 & 0 & 10 &  \\
%		\hline
%		$z$ & $4+3 M$ & $3+M$ & 0 & $-M$ & 0 & -1 & $10 M$ &  \\
%		\hline
%	\end{tabular}
%\end{center}
%
%D'où,
%
%\begin{center}
%	\begin{tabular}{|c|cccccc|c|c|}
%		\hline
%		Variables de base & $x_{1}$ & $x_{2}$ & $e_{1}$ & $e_{2}$ & $t_{1}$ & $z$ & $b$ & $R$ \\
%		\hline
%		$x_{1}$ & 1 & 1 & 1 & 0 & 0 & 0 & 2 &  \\
%		$t_{1}$ & 0 & -2 & -3 & -1 & 1 & 0 & 10 &  \\
%		\hline
%		$z$ & 0 & $-1-2 M$ & $-4-3 M$ & 0 & $-M$ & -1 & $4 M-8$ &  \\
%		\hline
%	\end{tabular}
%\end{center}
%
%Le tableau de simplexe ci-dessus est optimal avec une variable artificielle dans la base.
%
%\section{Les problèmes à solutions multiples}
%Nous avons vu dans l'algorithme du simplexe que lorsque les coûts marginaux de toutes les variables hors base sont $\leq 0$, alors la solution actuelle est optimale. Mais s'ils sont tous non nuls $ (<0) $, alors la solution optimale est unique. Si par contre, l'un de ces coûts des variables hors base est nul, alors on peut avoir deux ou plusieurs solutions de base optimales. Pour illustrer, reprenons l'exercice de la section résolution graphique qui admet une infinité de solutions optimales.
%
%$$
%\left\{\begin{aligned}
%	\displaystyle\min -50 x_{1}-30 x_{2} & \\
%	\text { s.t } 10 x_{1}+6 x_{2} & \leq 45 \\
%	4 x_{1}+6 x_{2} & \leq 36 \\
%	2 x_{1}+6 x_{2} & \leq 27 \\
%	x & \geq 0
%\end{aligned}\right.
%$$
%
%\begin{center}
%	\begin{tabular}{|c|cccccc|c|c|}
%		\hline
%		Variables de base & $x_{1}$ & $x_{2}$ & $e_{1}$ & $e_{2}$ & $e_{3}$ & $-z$ & $b$ & $R$ \\
%		\hline
%		$e_{1}$ & 10 & 6 & 1 & 0 & 0 & 0 & 45 &  \\
%		$e_{2}$ & 4 & 6 & 0 & 1 & 0 & 0 & 36 &  \\
%		$e_{3}$ & 2 & 6 & 0 & 0 & 1 & 0 & 27 &  \\
%		\hline
%		$-z$ & -50 & -30 & 0 & 0 & 0 & 1 & 0 &  \\
%		\hline
%	\end{tabular}
%\end{center}
%
%La variable $x_{1}$ entre dans la base pour remplacer la variable $e_{1}$
%
%\begin{center}
%	\begin{tabular}{|c|cccccc|c|c|}
%		\hline
%		Variables de base & $x_{1}$ & $x_{2}$ & $e_{1}$ & $e_{2}$ & $e_{3}$ & $-z$ & $b$ & $R$ \\
%		\hline
%		$x_{1}$ & 1 & $\dfrac{3}{5}$ & $\dfrac{1}{10}$ & 0 & 0 & 0 & $\dfrac{9}{2}$ &  \\
%		$e_{2}$ & 0 & $\dfrac{18}{5}$ & $-\dfrac{2}{5}$ & 1 & 0 & 0 & 18 &  \\
%		$e_{3}$ & 0 & $\dfrac{24}{5}$ & $-\dfrac{1}{5}$ & 0 & 1 & 0 & 18 &  \\
%		\hline
%		$-z$ & 0 & 0 & 5 & 0 & 0 & 1 & 225 &  \\
%		\hline
%	\end{tabular}
%\end{center}
%
%Nous obtenons une solution de base optimale dont les variables de base sont $x_{1}, e_{2}$ et $e_{3}$. Mais le coût marginal de la variable hors base $x_{2}$ est nul. Faisons entrer cette variable dans la base et appliquons la règle du plus petit rapport pour choisir la variable de sortie qui est dans ce cas $e_{3}$. Nous obtenons le tableau suivant
%
%\begin{center}
%	\begin{tabular}{|c|cccccc|c|c|}
%		\hline
%		Variables de base & $x_{1}$ & $x_{2}$ & $e_{1}$ & $e_{2}$ & $e_{3}$ & $-z$ & $b$ & $R$ \\
%		\hline
%		$x_{1}$ & 1 & 0 & $\dfrac{1}{8}$ & 0 & $-\dfrac{1}{8}$ & 0 & $\dfrac{9}{4}$ &  \\
%		$e_{2}$ & 0 & 0 & $-\dfrac{1}{4}$ & 1 & $-\dfrac{3}{4}$ & 0 & $\dfrac{9}{2}$ &  \\
%		$x_{2}$ & 0 & 1 & $-\dfrac{1}{24}$ & 0 & $\dfrac{5}{24}$ & 0 & $\dfrac{15}{4}$ &  \\
%		\hline
%		$-z$ & 0 & 0 & 5 & 0 & 0 & 1 & 225 &  \\
%		\hline
%	\end{tabular}
%\end{center}
%
%Nous obtenons donc une autre solution de base optimale dont les variables de base sont $x_{1}, e_{2}$ et $x_{2}$.
%
%\subsection{Les problèmes à solution dégénérée}
%En appliquons la règle du plus petit rapport pour désigner la variable de sortie, il est possible que le plus petit rapport soit atteint pour deux contraintes. En effectuant le changement de base, nous obtenons alors une solution de base avec l'une des composante nulle. Nous disons qu'on a une dégénérescence. Ceci peut entraîner des difficultés d'exécution et d'efficacité de l'algorithme du simplexe. En effet, nous pouvons faire plusieurs changement de base sans améliorer la valeur de la fonction objectif. Nous pouvons même retrouver une base déjà rencontrer à une itération précédente. On dit alors qu'on a un phénomène de cyclage. Afin d'illustrer ce type de dégénérescence. considérons l'exemple suivant :
%
%$$
%\left\{\begin{array}{rl}
%	\displaystyle\min & -4 x_{1}-3 x_{3} \\
%	s . t & x_{1}-x_{2} \\
%	& \leq 2 \\
%	& 2 x_{1}+x_{3} \\
%	& \leq 4 \\
%	x_{1}+x_{2}+x_{3} & \leq 3 \\
%	& x_{1}, x_{2}, x_{3}
%\end{array} \geq 0\right.
%$$
%
%Introduisons les variables d'écart et écrivons le premier tableau du simplexe :
%
%\begin{center}
%	\begin{tabular}{|c|ccccccc|c|c|}
%		\hline
%		Variables de base & $x_{1}$ & $x_{2}$ & $x_{3}$ & $e_{1}$ & $e_{2}$ & $e_{3}$ & $-z$ & $b$ & $R$ \\
%		\hline
%		$e_{1}$ & 1 & -1 & 0 & 1 & 0 & 0 & 0 & 2 &  \\
%		$e_{2}$ & 2 & 0 & 1 & 0 & 1 & 0 & 0 & 4 &  \\
%		$e_{3}$ & 1 & 1 & 1 & 0 & 0 & 1 & 0 & 3 &  \\
%		\hline
%		$-z$ & -4 & 0 & -3 & 0 & 0 & 0 & 1 & 0 &  \\
%		\hline
%	\end{tabular}
%\end{center}
%
%La variable $x_{1}$ entre dans la base, mais en appliquons la règle du plus petit rapport, nous constatons que le minimum des rapports est atteint pour les deux premières contraintes, alors que seulement une des deux doit quitter la base. Choisissons arbitrairement $e_{1}$ comme variable sortante. nous obtenons alors le tableau suivant :
%
%\begin{center}
%	\begin{tabular}{|c|ccccccc|c|c|}
%		\hline
%		Variables de base & $x_{1}$ & $x_{2}$ & $x_{3}$ & $e_{1}$ & $e_{2}$ & $e_{3}$ & $-z$ & $b$ & $R$ \\
%		\hline
%		$x_{1}$ & 1 & -1 & 0 & 1 & 0 & 0 & 0 & 2 &  \\
%		$e_{2}$ & 0 & 2 & 1 & -2 & 1 & 0 & 0 & 0 &  \\
%		$e_{3}$ & 0 & 2 & 1 & -1 & 0 & 1 & 0 & 1 &  \\
%		\hline
%		$-z$ & 0 & -4 & -3 & 4 & 0 & 0 & 1 & 8 &  \\
%		\hline
%	\end{tabular}
%\end{center}
%
%cette solution de base n'est pas optimale puisque les coûts marginaux des variables hors base ne sont pas tous positifs. Cette solution de base est dite dégénérée puisque l'une des variables de base est nulle ( $e_{2}=0$ ). La variable $x_{2}$ entre dans la base. En appliquons la règle du plus petit rapport, nous constatons que la variable $e_{2}$ quitte la base et nous obtenons le tableau suivant
%
%\begin{center}
%	\begin{tabular}{|c|ccccccc|c|c|}
%		\hline
%		Variables de base & $x_{1}$ & $x_{2}$ & $x_{3}$ & $e_{1}$ & $e_{2}$ & $e_{3}$ & $-z$ & $b$ & $R$ \\
%		\hline
%		$x_{1}$ & 1 & 0 & $\dfrac{1}{2}$ & 0 & $\dfrac{1}{2}$ & 0 & 0 & 2 &  \\
%		$x_{2}$ & 0 & 1 & $\dfrac{1}{2}$ & -1 & $\dfrac{1}{2}$ & 0 & 0 & 0 &  \\
%		$e_{3}$ & 0 & 0 & 0 & 1 & -1 & 1 & 0 & 1 &  \\
%		\hline
%		$-z$ & 0 & 0 & -1 & 0 & 2 & 0 & 1 & 8 &  \\
%		\hline
%	\end{tabular}
%\end{center}
%
%Nous constatons que la nouvelle solution de base est non optimale et qu'elle est aussi dégénérée $\left(x_{2}=0\right)$. Ce changement de base n'a pas modifié la valeur de la fonction objectif $(z=-8)$. A l'itération suivante de l'algorithme du simplexe, cette variable $x_{2}$ quitte la base et sera remplacée par la variable $x_{3}$
%
%\begin{center}
%	\begin{tabular}{|c|c|c|c|c|c|c|c|c|c|}
%		\hline
%		Variables de base & $x_{1}$ & $x_{2}$ & $x_{3}$ & $e_{1}$ & $e_{2}$ & $e_{3}$ & $-z$ & $b$ & $R$ \\
%		\hline
%		$x_{1}$ & 1 & -1 & 0 & 1 & 0 & 0 & 0 & 2 &  \\
%		$x_{3}$ & 0 & 2 & 1 & -2 & 1 & 0 & 0 & 0 &  \\
%		$e_{3}$ & 0 & 0 & 0 & 1 & -1 & 1 & 0 & 1 &  \\
%		\hline
%		$-z$ & 0 & 2 & 0 & -2 & 3 & 0 & 1 & 8 &  \\
%		\hline
%	\end{tabular}
%\end{center}
%
%Encore une fois, ce changement de base n'a pas modifié la valeur de la fonction objectif. En plus la nouvelle solution de base est aussi non optimale. Nous pouvons donc se poser la question de savoir si nous allons continuer ainsi sans pouvoir trouver la solution de base optimale. Heureusement ce phénomène de dégénérescence s'arrête à l'itération suivante
%
%\begin{center}
%	\begin{tabular}{|c|c|c|c|c|c|c|c|c|c|}
%		\hline
%		Variables de base & $x_{1}$ & $x_{2}$ & $x_{3}$ & $e_{1}$ & $e_{2}$ & $e_{3}$ & $-z$ & $b$ & $R$ \\
%		\hline
%		$x_{1}$ & 1 & -1 & 0 & 0 & 1 & -1 & 0 & 1 &  \\
%		$x_{3}$ & 0 & 2 & 1 & 0 & -1 & 2 & 0 & 2 &  \\
%		$e_{1}$ & 0 & 0 & 0 & 1 & -1 & 1 & 0 & 1 &  \\
%		$-z$ & 0 & 2 & 0 & 0 & 1 & 2 & 1 & 10 &  \\
%		\hline
%	\end{tabular}
%\end{center}
%
%%\textbf{Remarque 3.5} 
%\begin{remarque}
%	Dans le cas où on a plusieurs variables candidates d'entrer en base et plusieurs variables candidates de sortir de la base, la règle du plus petit indice, découverte par R. Bland, peut nous éviter le cyclage infini. Cette règle consiste à choisir systématiquement, comme variable entrante, parmi les multiples candidats, la variable $x_{e}$ qui correspond au plus petit indice $e$ et comme variable sortante, la variable $x_{s}$ qui correspond au plus petit indice $s$.
%\end{remarque}	
	
	
	\chapter{Dualité}
	
	\section{Motivation et définition}
	Dans l'exemple~\ref{ex:prod}, un acheteur propose un prix unitaire $y_i\ge0$ pour chaque ressource ; l'entreprise n'accepte que si, pour chaque produit $P_j$, la valeur des ressources consommées est au moins égale au bénéfice $c_j$ ; l'acheteur minimise le coût total :
	\[\min\ 81y_1+55y_2+20y_3\quad\text{s.c.}\quad3y_1+4y_2+2y_3\ge6,\ \ 9y_1+5y_2+y_3\ge4,\ \ y\ge0.\]
	\begin{definition}
		Au \textbf{primal} $(PL)$ : $\max f(x)=c\T x$ s.c. $Ax\le b$, $x\ge0$, on associe le \textbf{dual} $(PLD)$ : $\min g(y)=b\T y$ s.c. $A\T y\ge c$, $y\ge0$.
	\end{definition}
	Le primal a $m$ contraintes $\le$ et $n$ variables positives ; le dual a $m$ variables positives et $n$ contraintes $\ge$. Pour une forme canonique mixte :
	\begin{center}\small
		\begin{tabular}{|c|c|}\hline
			\textbf{Primal} (max) & \textbf{Dual} (min)\\\hline
			contrainte $i\in I_1$ : $\sum_ja_{ij}x_j\le b_i$ & $y_i\ge0$\\
			contrainte $i\in I_2$ : $\sum_ja_{ij}x_j=b_i$ & $y_i$ de signe quelconque\\
			$x_j\ge0$ ($j\in J_1$) & contrainte $\sum_ia_{ij}y_i\ge c_j$\\
			$x_j$ libre ($j\in J_2$) & contrainte $\sum_ia_{ij}y_i=c_j$\\\hline
		\end{tabular}
	\end{center}
	\begin{proposition}
		Le dual du dual est le primal. (On écrit $\min b\T y,\ A\T y\ge c$ comme $-\max(-b)\T y,\ A\T y\le-c$ ; son dual est $\min(-c)\T x$ s.c. $Ax\ge-b$, $x\ge0$, soit $\max c\T x$, $Ax\le b$, $x\ge0$.)
	\end{proposition}
	
	\begin{exemple}[Écrire un dual, mode d'emploi]
		Pour un primal \emph{min} avec contraintes $\ge$ (et $x\ge0$), le dual est un \emph{max} avec contraintes $\le$ et $y\ge0$ : les rôles de $b$ et $c$ s'échangent et on transpose $A$. Problème des pilules : $\min x_1+x_2$ s.c. $2x_1+x_2\ge12$, $5x_1+8x_2\ge74$, $x_1+6x_2\ge24$, $x\ge0$. Une variable duale $y_i\ge0$ par contrainte ; une contrainte duale par variable primale :
		\[\max\ 12y_1+74y_2+24y_3\quad\text{s.c.}\quad2y_1+5y_2+y_3\le1,\ \ y_1+8y_2+6y_3\le1,\ \ y\ge0.\]
		(Pour retrouver ces règles, on réécrit le primal en maximisation avec des contraintes $\le$, on applique la définition, puis on remet le dual en forme $\max$.)
	\end{exemple}
	
	\section{Théorèmes de dualité}
	\begin{theorem}[Dualité faible]
		Soit $x$ réalisable pour $(PL)$ et $y$ réalisable pour $(PLD)$. Alors (1) $f(x)\le g(y)$ ; (2) si $f(x)=g(y)$, $x$ et $y$ sont optimaux.
	\end{theorem}
	\begin{proof}
		\textbf{(1)} Les hypothèses utilisées : $A\T y\ge c$ et $y\ge0$ (dual réalisable) ; $Ax\le b$ et $x\ge0$ (primal réalisable).
		\emph{Étape 1.} Comme $c\le A\T y$ composante par composante et $x\ge0$, on peut multiplier par $x_j\ge0$ sans changer le sens : $c\T x\le(A\T y)\T x$.
		\emph{Étape 2.} $(A\T y)\T x=y\T Ax$ (propriété de la transposée).
		\emph{Étape 3.} Comme $Ax\le b$ et $y\ge0$, on multiplie par $y_i\ge0$ : $y\T Ax\le y\T b$.
		\emph{Bilan} : $f(x)=c\T x\le y\T Ax\le b\T y=g(y)$.
		\textbf{(2)} Soit $x^*,y^*$ réalisables avec $f(x^*)=g(y^*)$. Pour tout $x$ réalisable, (1) donne $f(x)\le g(y^*)=f(x^*)$ : $x^*$ est optimal. Pour tout $y$ réalisable, $g(y)\ge f(x^*)=g(y^*)$ : $y^*$ est optimal.
	\end{proof}
	\begin{theorem}[Dualité forte]
		Si $(PL)$ admet une solution optimale $x^*$, alors $(PLD)$ admet une solution optimale $y^*$ et $f(x^*)=g(y^*)$.
	\end{theorem}
	\begin{proof}
		\emph{Étape 1 (forme standard).} On ajoute des variables d'écart : $(A\mid I_m)\binom xe=b$. On note $\tilde A=(A\mid I_m)$ cette matrice et $\tilde c=(c,0)$ le vecteur de coûts (les écarts ont un coût nul).
		\emph{Étape 2 (base optimale).} Le primal a un optimum ; il est atteint en une solution de base réalisable (chapitre~6). Soit $B^*$ la base correspondante, $x^*_B=\tilde A_{B^*}^{-1}b$ et $f(x^*)=\tilde c_{B^*}\T\tilde A_{B^*}^{-1}b$. Notons $H^*$ les indices hors base. Le critère d'optimalité du simplexe (chapitre~7) dit que les coûts réduits sont $\le0$ : $\tilde c_{H^*}\T-\tilde c_{B^*}\T\tilde A_{B^*}^{-1}\tilde A_{H^*}\le0$.
		\emph{Étape 3 (définition de $y^*$).} Posons $y^*=(\tilde A_{B^*}^{-1})\T\tilde c_{B^*}\in\R^m$, de sorte que $\tilde c_{B^*}\T\tilde A_{B^*}^{-1}=y^{*\top}$. Le critère d'optimalité devient $\tilde c_j-\tilde A_j\T y^*\le0$ pour toute colonne $j$ hors base, soit $\tilde A_j\T y^*\ge\tilde c_j$. Pour une colonne de base, $\tilde A_j\T y^*=\tilde c_j$ (par construction de $y^*$). Donc $\tilde A\T y^*\ge\tilde c$.
		\emph{Étape 4 ($y^*$ est réalisable pour le dual).} $\tilde A\T y^*\ge\tilde c$ se lit en deux blocs. Bloc des colonnes de $A$ : $A\T y^*\ge c$. Bloc des colonnes d'écart ($\tilde A_j=e_i$, $\tilde c_j=0$) : $y_i^*\ge0$. Donc $A\T y^*\ge c$ et $y^*\ge0$ : $y^*$ est réalisable pour $(PLD)$. (C'est ici que la positivité de $y^*$ est justifiée.)
		\emph{Étape 5 (mêmes valeurs).} $g(y^*)=b\T y^*=b\T(\tilde A_{B^*}^{-1})\T\tilde c_{B^*}=\tilde c_{B^*}\T\tilde A_{B^*}^{-1}b=f(x^*)$.
		\emph{Étape 6.} $x^*$ et $y^*$ sont réalisables avec $f(x^*)=g(y^*)$ : par la dualité faible (2), $y^*$ est optimal.
	\end{proof}
	\begin{theorem}[Quatre situations possibles]
		Pour un primal et son dual, seules sont possibles : (a) les deux ont un optimum fini, de même valeur ; (b) primal à optimum infini et dual irréalisable ; (c) primal irréalisable et dual à optimum infini ; (d) les deux irréalisables.
	\end{theorem}
	\begin{center}\small
		\begin{tabular}{l|ccc}
			Primal $\backslash$ Dual & optimum fini & optimum infini & irréalisable\\\hline
			optimum fini & (a) & impossible & impossible\\
			optimum infini & impossible & impossible & (b)\\
			irréalisable & impossible & (c) & (d)
		\end{tabular}
	\end{center}
	
	\begin{proof}[Preuve du théorème des quatre situations]
		\emph{(i) Si le primal a un optimum fini}, la dualité forte donne un dual avec optimum fini : on est en (a), et les cas ``dual infini'' ou ``dual irréalisable'' sont exclus. Symétriquement (le dual du dual est le primal), si le dual a un optimum fini, le primal aussi.
		\emph{(ii) Si le primal est à optimum infini} ($f$ non majorée) : s'il existait $y$ réalisable pour le dual, la dualité faible donnerait $f(x)\le g(y)$ pour tout $x$ réalisable, donc $f$ serait majorée : contradiction. Le dual est irréalisable : cas (b). Symétriquement, un dual à optimum infini ($g$ non minorée) force un primal irréalisable : cas (c).
		\emph{(iii)} Il reste le cas où aucun des deux n'a d'optimum fini et aucun n'est à optimum infini : les deux sont irréalisables, cas (d). Les autres combinaisons du tableau sont donc impossibles.
	\end{proof}
	
	\section{Conditions d'optimalité primal-dual (COPD)}
	\begin{exemple}[Illustration]
		$\min-x_1-x_2$ s.c. $-x_1-3x_2\ge-3$, $-2x_1-x_2\ge-2$, $x\ge0$ (primal) a pour solution $x^*=(\frac35,\frac45)$, $f=-\frac75$. Son dual $\max-3y_1-2y_2$ s.c. $-y_1-2y_2\le-1$, $-3y_1-y_2\le-1$, $y\ge0$ a pour solution $y^*=(\frac15,\frac25)$, $g=-\frac75$ : les valeurs optimales coïncident (dualité forte). Les deux contraintes primales sont saturées en $x^*$ et $y^*>0$ ; $x^*>0$ et les deux contraintes duales sont saturées : les COPD ci-dessous sont vérifiées.
	\end{exemple}
	
	\begin{theorem}[Écarts complémentaires]\label{thm:copd}
		Soient $x$, $y$ réalisables pour $(PL)$, $(PLD)$ (forme canonique pure). Alors $x$ et $y$ sont optimaux $\iff$
		\[\forall i:\ y_i\Bigl(b_i-\sum_ja_{ij}x_j\Bigr)=0,\qquad\forall j:\ x_j\Bigl(\sum_ia_{ij}y_i-c_j\Bigr)=0.\]
		Autrement dit : une contrainte non saturée dans un programme force la variable correspondante de l'autre à $0$ ; une variable strictement positive force la contrainte correspondante de l'autre à être une égalité.
	\end{theorem}
	\begin{proof}
		\emph{Étape 1 (écarts).} Pour $x,y$ réalisables, posons $e=b-Ax\ge0$ (écarts primaux) et $\epsilon=A\T y-c\ge0$ (écarts duaux).
		\emph{Étape 2 (une identité clé).} On remplace $b=Ax+e$ et $c=A\T y-\epsilon$ :
		\[g(y)-f(x)=b\T y-c\T x=(Ax+e)\T y-(A\T y-\epsilon)\T x=y\T Ax+e\T y-y\T Ax+\epsilon\T x=e\T y+\epsilon\T x.\]
		\emph{Étape 3 (somme de termes positifs).} Comme $e,y,\epsilon,x\ge0$, tous les termes $e_iy_i$ et $\epsilon_jx_j$ sont $\ge0$ : leur somme est nulle si et seulement si chacun est nul.
		\emph{Étape 4 (conclusion).} ($\Rightarrow$) Si $x,y$ sont optimaux, la dualité forte donne $f(x)=g(y)$, donc la somme est nulle : $e_iy_i=0$ et $\epsilon_jx_j=0$ pour tous $i,j$ (c'est la complémentarité). ($\Leftarrow$) Si tous ces produits sont nuls, $g(y)-f(x)=0$ : par la dualité faible (2), $x$ et $y$ sont optimaux.
	\end{proof}
	\begin{exemple}[Production]
		Le primal admet $x^*=(\frac{15}2,5)$, $e^*=(\frac{27}2,0,0)$. $e_1^*>0\Rightarrow y_1^*=0$ ; $x_1^*,x_2^*>0\Rightarrow$ les deux contraintes duales sont des égalités :
		$4y_2+2y_3=6$, $5y_2+y_3=4$, d'où $y^*=(0,\frac13,\frac73)$ et $g(y^*)=\frac{55}3+\frac{140}3=65=f(x^*)$.
	\end{exemple}
	\textbf{Lecture du dual dans le tableau final du primal.} Dans le tableau optimal $T_2$ (exemple~\ref{ex:simpl}), les coûts réduits des variables d'écart donnent le dual : $y_i^*=-\bar c_{e_i}$, soit $(0,\frac13,\frac73)$ ; les coûts réduits de $x_j$, nuls ici (variables de base), valent $-\epsilon_j^*$. Réciproquement, le tableau optimal du dual donne les valeurs primales $x_j^*$ et $e_i^*$ (lues dans la ligne des coûts) : $x^*=(\frac{15}2,5)$, $e^*=(\frac{27}2,0,0)$. Les valeurs de $y^*$ sont les \emph{prix fictifs} (ou duaux) des ressources : $y_3^*=\frac73$ est le gain marginal d'une unité de matière première supplémentaire.
	\textbf{Tableau optimal du dual de l'exemple de production} (variables $y_1,y_2,y_3$, d'excès $\epsilon_1,\epsilon_2$ ; problème de minimisation) :
	\begin{center}\small
		\begin{tabular}{c|ccccc|c}
			&$y_1$&$y_2$&$y_3$&$\epsilon_1$&$\epsilon_2$&$b$\\\hline
			$y_3$&$-\frac72$&0&1&$-\frac56$&$\frac23$&$\frac73$\\ $y_2$&$\frac52$&1&0&$\frac16$&$-\frac13$&$\frac13$\\\hline
			$-z$&$\frac{27}2$&0&0&$\frac{15}2$&5&$-65$
		\end{tabular}
	\end{center}
	Les coûts réduits $\frac{27}2,\frac{15}2,5$ de $y_1,\epsilon_1,\epsilon_2$ sont les valeurs primales $e_1^*,x_1^*,x_2^*$. La symétrie est parfaite : le tableau optimal d'un problème contient la solution de l'autre.
	\begin{retenir}
		Dual : on échange $\max\leftrightarrow\min$, $c\leftrightarrow b$, $A\leftrightarrow A\T$. $f(x)\le g(y)$ toujours ; égalité $\iff$ optimalité des deux ; optimum fini d'un côté $\Rightarrow$ optimum fini de l'autre. COPD : on déduit une solution de l'autre.
	\end{retenir}
	
	\chapter*{Fiche récapitulative et conclusion}
	\addcontentsline{toc}{chapter}{Fiche récapitulative et conclusion}
	\textbf{Conditions d'optimalité.}
	\begin{center}\small
		\begin{tabular}{|p{2.6cm}|p{5.3cm}|p{6.3cm}|}\hline
			\textbf{Problème} & \textbf{Nécessaire} (minimum local) & \textbf{Suffisant}\\\hline
			Sans contrainte & $\grad f(x^*)=0$ ; $\hess f(x^*)\succeq0$ & $\grad f=0$ et $\hess f\succ0$ $\Rightarrow$ min local strict ; $f$ convexe : $\grad f=0\Rightarrow$ min global\\\hline
			Égalités $g(x)=0$ & $x^*$ régulier : $\grad f+\sum\lambda_i\grad g_i=0$ & $\grad\M=0$, $x^*$ régulier et $h\T\hess_x\M h>0$ sur $\ker D\T$\\\hline
			Égalités + inégalités & KKT (QCMF/LICQ) : stationnarité, admissibilité, $\mu\ge0$, $\mu_jg_j=0$ & KKT + problème convexe ($\Rightarrow$ min global)\\\hline
		\end{tabular}
	\end{center}
	\textbf{Qualifications.} LICQ $\Rightarrow$ MFCQ ; Slater $\Rightarrow$ MFCQ (convexe) ; LICQ $\Rightarrow$ MFCQ $\Rightarrow$ CPLD $\Rightarrow$ ACQ.\\
	\textbf{Existence/unicité.} continue + coercive $\Rightarrow$ minimum global ; strictement convexe $\Rightarrow$ au plus un minimum. Formes quadratiques : $\grad f=Ax-b$, $\hess f=A$, Sylvester.\\
	\textbf{PL.} Sommet $\iff$ solution de base réalisable ; l'optimum, s'il existe, est atteint en un sommet. Simplexe (max) : entrante $=\max$ coût réduit $>0$ ; sortante $=\min$ ratio $b_i/a_{ie}$ ($a_{ie}>0$) ; optimum si coûts réduits $\le0$ ; colonne $\le0$ $\Rightarrow$ non borné ; artificielle $>0$ à la fin $\Rightarrow$ irréalisable ; coût réduit nul $\Rightarrow$ solutions multiples ; dégénérescence $\Rightarrow$ règle de Bland.\\
	\textbf{Dualité.} $\max c\T x$, $Ax\le b$, $x\ge0$ $\leftrightarrow\min b\T y$, $A\T y\ge c$, $y\ge0$ ; faible, forte, COPD.
	
	\textbf{Formules clés.}
	\begin{center}\small
		\begin{tabular}{|p{4.3cm}|p{9.7cm}|}\hline
			Forme quadratique & $f=\frac12x\T Ax-b\T x+c$ : $\grad f=Ax-b$, $\hess f=A$, minimum $x^*=A^{-1}b$ si $A\succ0$\\\hline
			Taylor ordre 2 & $f(x+h)=f(x)+\grad f(x)\T h+\frac12h\T\hess f(x)h+o(\norm h^2)$\\\hline
			Lagrange (égalités) & $\M=f+\sum\lambda_ig_i$ ; $\grad\M=0$ ; test : $p(t)=\det\begin{pmatrix}A-tI&D\\D\T&0\end{pmatrix}$\\\hline
			KKT & $\grad f+\sum\lambda_i\grad h_i+\sum\mu_j\grad g_j=0$, $h=0$, $g\le0$, $\mu\ge0$, $\mu_jg_j=0$\\\hline
			Solution de base & $x_B=A_B^{-1}b$, $x_H=0$\\\hline
			Coûts réduits & $L_H\T=c_H\T-c_B\T A_B^{-1}A_H$ ; sortante : $\min_i\{(\bar x_B)_i/z_i\mid z_i>0\}$, $z=A_B^{-1}A_e$\\\hline
			Dualité & $\max c\T x,\ Ax\le b,\ x\ge0\ \leftrightarrow\ \min b\T y,\ A\T y\ge c,\ y\ge0$ ; $f(x)\le g(y)$\\\hline
		\end{tabular}
	\end{center}
	\textbf{Conclusion.} Le cours suit une progression unique : \emph{existence} (compacité, coercivité), \emph{conditions d'optimalité} (ordres 1 et 2, convexité), \emph{contraintes} (Lagrange, KKT), puis le cas linéaire où la géométrie des sommets, le simplexe et la dualité donnent un algorithme complet. L'erreur la plus fréquente est de confondre conditions nécessaires et suffisantes ; la seconde est d'oublier la qualification des contraintes.
	
	\chapter*{Annexe : errata du cours source}
	\addcontentsline{toc}{chapter}{Annexe : errata du cours source}
	\footnotesize
	Les points suivants ont été corrigés dans cette synthèse ; ce sont des erreurs du document source, listées pour que rien ne soit modifié silencieusement. Les compléments de ma part (non présents dans la source) sont marqués \emph{(compl.)}.
	\begin{multicols}{2}
		\begin{itemize}[leftmargin=*,itemsep=1pt]
			\item \textbf{Ch.1--2.} ``Strictement convexe $\iff\hess f\succ0$'' est faux ($x^4$) : seule $\Leftarrow$ est vraie. ``Une fonction convexe non stricte peut avoir plusieurs minimums locaux'' : lire plusieurs minimiseurs globaux. Preuve de $\Rightarrow$ du théorème CNS convexe invalide (utiliser la condition du premier ordre) ; ``$-\grad f=0\Rightarrow$'' lire ``$\Leftarrow$''. Condition suffisante d'ordre 2 : hessienne évaluée en $x$ et non en $x^*$ ; preuve non uniforme (utiliser $\lambda_{\min}$).
			\item \textbf{Exo $f_1$} : $\det=192$ (et non 240). \textbf{Exo $f_3$} : $f(\sqrt2,-\sqrt2)=-8$ (et non 0) ; l'argument ``$(1-x^2)^3=-1$ impossible'' est faux ($x^2=2$ convient) : on additionne les deux équations ; en $(0,0)$ le test est inconclusif, il faut étudier $f(t,\pm t)$. Exo 2, question 1 : $\alpha=\beta=-2$ ne prouve pas la coercivité ; $\alpha=2,\beta=-18$ convient \emph{(compl.)}.
			\item \textbf{Ch.3.} Exemple~\ref{ex:lag1} : $\lambda^*=-7$ (et non $7$) ; la conclusion (max, $f=8$) est inchangée. Hessienne bordée écrite avec deux signes différents ; $m$ et $p$ mélangés ; $D$ non défini.
			\item \textbf{Ch.4.} Exemple de Fritz John ($\min x$ sur $y=0$) sans minimum : remplacé ; exemple ``contraintes actives'' sans minimum et contenant un texte de correction : supprimé. ``CRCQ $\Rightarrow$ QCMF'' faux. KKT suffisant en cas convexe : sans qualification. Exemple du triangle : trois gradients ne peuvent être indépendants dans $\R^2$ (la LICQ tient car au plus deux contraintes sont actives). Exemple disque/demi-plan : cas $\mu_1>0$ erronés, $\mu_1<0$ ; l'argument de symétrie est faux ; un seul point KKT. Exemple égalité/inégalité : $h=0$ non convexe, conclure par compacité. Figure du triangle : la région est un demi-plan.
			\item \textbf{Ch.5--6.} $x_j^-=\max(-x_j,0)\ge0$. Exemple de transformation canonique : égalité écrite $\ge$ et contrainte $\ge$ non convertie. $\mathbf D_R$ : polyèdre (non polytope). Équations de la méthode algébrique : $e_2=15-3x_2+2e_3$ ; $x_1=\frac{15}2+\frac16e_2-\frac56e_3$ ; ligne $z$ : $-\frac74$ (et non $-\frac75$) pour $x_1$. Figure de $\mathbf D_R$ : intercepts $(\frac32,1,3)$.
			\item \textbf{Ch.7.} Big $M$ : coût réduit de $t_1$ vaut $-\frac25-M$. Remarque 3.3 : l'artificielle de base est $>0$ (et non ``négative'') si irréalisable. Exemple non borné : colonne $e_3$ mal recopiée ; critère : $a_{ie}\le0$ pour tout $i$. Exemple impossible : second membre de $t_1$ : $4$ (et non $10$). Signes de la ligne $-z$ de l'exemple de minimisation incohérents (unifiés). Exercice~2.1 sans solution rédigée \emph{(compl.)}.
			\item \textbf{Ch.8.} Tableau des situations primal/dual erroné : le cas (d) est possible ; quatre situations. Dualité forte : justifier $y^*\ge0$ par les coûts réduits des écarts. COPD : ``si $\epsilon_i\neq0$ alors $x_i=0$''. Lecture du dual : $y_j=-c_{n+j}$.
			\item \textbf{Forme.} Fichiers \texttt{ULBO} et \texttt{2025\_10\_05\_*} absents : logo remplacé par un titre texte et figures redessinées en TikZ. Références ``Th.~2.2.5/2.2.7'' non résolues.
		\end{itemize}
	\end{multicols}
\end{document}
